% !TEX program = xelatex
\documentclass[12pt,a4paper,fontset=fandol]{ctexbook}

% ==================== 页面 ====================
\usepackage[top=2.5cm,bottom=2.5cm,left=2.5cm,right=2.5cm]{geometry}

% ==================== 字体 ====================
\usepackage{fontspec}
\usepackage{iffont}
% CJK 钉 fandol（两侧都取自 TeX 树里的同一份 FandolSong-Regular.otf）；
% mono 钉 DejaVu Sans Mono：CI 由 fontconfig 命中同名，本机没有该系统字体，
% 就用 MiKTeX dejavu 包里的同一版本 TTF（kpathsea 按文件名命中）。
\IfFontExistsTF{DejaVu Sans Mono}
  {\setmonofont{DejaVu Sans Mono}[Scale=MatchLowercase]}
  {\setmonofont[Extension=.ttf,
      UprightFont=DejaVuSansMono,
      BoldFont=DejaVuSansMono-Bold,
      ItalicFont=DejaVuSansMono-Oblique,
      BoldItalicFont=DejaVuSansMono-BoldOblique]
     {DejaVuSansMono}[Scale=MatchLowercase]}

% ==================== 颜色 ====================
\usepackage[dvipsnames,svgnames,x11names]{xcolor}
\definecolor{codebg}{HTML}{F5F5F5}
\definecolor{codeframe}{HTML}{E0E0E0}
\definecolor{cppkeyword}{HTML}{1A1AFF}
\definecolor{cppcomment}{HTML}{2E8B2E}
\definecolor{cppstring}{HTML}{B22222}
\definecolor{linenumgray}{HTML}{999999}
\definecolor{toolheadbg}{HTML}{2D3E50}
\definecolor{yamlkey}{HTML}{7A3E9D}
\definecolor{diffadd}{HTML}{1A7F37}
\definecolor{diffdel}{HTML}{B62324}

% ==================== listings ====================
\usepackage{listings}

% ---- 所有代码块共用的排版（先定义，各语言 style 组合它）----
\lstdefinestyle{codebase}{
  basicstyle=\small\ttfamily,
  keywordstyle=\color{cppkeyword}\bfseries,
  commentstyle=\color{cppcomment}\itshape,
  stringstyle=\color{cppstring},
  numberstyle=\tiny\color{linenumgray},
  numbers=left,
  frame=single,
  rulecolor=\color{codeframe},
  framesep=6pt,
  breaklines=true,
  tabsize=4,
  columns=flexible,
  keepspaces=true,
  backgroundcolor=\color{codebg},
  showstringspaces=false,
  escapeinside={(*@}{@*)},
}

\lstdefinestyle{cppstyle}{
  style=codebase,
  language=C++,
  morekeywords={constexpr,consteval,constinit,decltype,auto,
    concept,requires,co_await,co_yield,co_return,
    noexcept,override,final,nullptr,static_assert,
    template,typename,namespace,using,mutable,volatile,
    explicit,operator,friend,inline,virtual,
    thread_local,alignas,alignof},
  deletekeywords={int,char,bool,float,double,long,short,unsigned,signed,void,wchar_t},
  emph={size_t,int32_t,uint32_t,int64_t,uint64_t,ssize_t,
    string,vector,map,set,unique_ptr,shared_ptr,optional,variant,any,
    span,string_view,FILE,wchar_t},
  emphstyle=\color{teal!80!black},
}
\lstset{style=cppstyle}

% ---- C：图形 API 的 C ABI 边界（Vulkan 是纯 C 接口）----
\lstdefinestyle{cstyle}{
  style=codebase,
  language=C,
  morekeywords={bool,FILE,size_t,restrict,inline},
  emph={size_t,int32_t,uint32_t,int64_t,uint64_t,ssize_t,
    VkInstance,VkDevice,VkBuffer,VkImage,VkImageSubresourceRange,
    VkCommandBuffer,VkPipeline,VkPipelineLayout,VkShaderModule,VkGraphicsPipelineCreateInfo,
    VkResult,VkFormat,VkExtent2D,VkExtent3D,VkDeviceSize,VkMemoryPropertyFlags,
    VkSubmitInfo,VkFence,VkSemaphore,VkSwapchainKHR,VkAllocationCallbacks},
  emphstyle=\color{teal!80!black},
}

% ---- YAML：CI 描述文件、着色器构建与资产管线配置 ----
\lstdefinestyle{yamlstyle}{
  style=codebase,
  morecomment=[l]{\#},
  morestring=[b]",
  morestring=[d]',
  alsoletter={-,:},
  morekeywords=[1]{Version,Checks,CheckOptions,HeaderFilterRegex,FormatStyle,
    User,WarningsAsErrors,CompilationDatabase,ExtraArgs,Add,Remove,
    Diagnostics,ClangTidy,UnusedIncludes,MissingIncludes,index,background-index,
    clang-tidy,InlayHints,Style,BasedOnStyle,IndentWidth,TabWidth,UseTab,
    ColumnLimit,BreakBeforeBraces,AlignAfterOpenBracket,AllowShortIfStatementsOnASingleLine,
    IncludeBlocks,SortIncludes,SeparateDefinitionBlocks,SpacesInAngles,
    Language,If,Else,Disable,Key,Name,files,exclude,repo,rev,hooks,id,entry,
    language_types,on,push_pull_request,steps,uses,runs-on,defaults},
  keywordstyle=[1]\color{yamlkey}\bfseries,
}

% ---- JSON：编译数据库、着色器变体表与工具配置 ----
\lstdefinestyle{jsonstyle}{
  style=codebase,
  morestring=[b]",
  alsoletter={-},
  morekeywords={directory,file,output,arguments,command,version,
    cmakeMinimumRequired,name,generator,binaryDir,sourceDir,caches,
    toolchainFile,C_CXX_FLAGS,defineConstants,environments,
    editorMode,rulerColumn,tabSize,insertSpaces,files,settings,
    stages,entry_points,spirv_variants,per_variant,defines,entry_name,
    shader_model,target_formats,asset,buffer_view,texture,material},
  keywordstyle=\color{yamlkey}\bfseries,
}

% ---- Shell：命令行与工具调用 ----
\lstdefinelanguage{ShellLD}{
  morekeywords={if,then,else,elif,fi,for,in,do,done,while,case,esac,
    function,return,local,export,source,echo,cd,set,exec,trap,read,
    printf,test,exit,shift,eval,alias},
  sensitive=true,
  morecomment=[l]{\#},
  morestring=[b]",
  morestring=[d]',
}
\lstdefinestyle{shellstyle}{
  style=codebase,
  language=ShellLD,
  keywordstyle=\color{cppkeyword}\bfseries,
  emph={glslangValidator,glslc,dxc,slangc,fxc,spirv-as,spirv-dis,spirv-cross,
    spirv-val,spirv-opt,spirv-remap,vkconfig,vkcube,vulkaninfo,renderdoccmd,
    naga,shaderc,
    grep,sed,awk,xargs,jq,python3,pip,pip3,pwsh,bash,git,
    cmake,ninja,make,ccache,lldb,gdb,code,nvim,vim,emacs,tar,mkdir,touch,ls,cat,
    pre-commit,actions,gh,yamllint,wc,diff,find,sort,uniq,tee,rm},
  emphstyle=\color{teal!80!black}\bfseries,
}

% ---- CMake ----
\lstdefinelanguage{CMakeLD}{
  morekeywords={cmake_minimum_required,project,add_executable,
    add_library,target_link_libraries,target_include_directories,
    target_compile_definitions,target_compile_options,set,if,elseif,
    else,endif,find_package,option,message,install,file,
    target_sources,add_subdirectory,enable_language,set_target_properties,
    FetchContent_Declare,FetchContent_MakeAvailable,list,string},
  sensitive=false,
  morecomment=[l]{\#},
  morestring=[b]",
}
\lstdefinestyle{cmakestyle}{
  style=codebase,
  language=CMakeLD,
  keywordstyle=\color{cppkeyword}\bfseries,
}

% ---- GLSL / HLSL：着色器语言本体 ----
% 数字算作标识符字符，否则 vec2、mat4x4、SV_Target 这类词会被拆断而匹配不上关键字
\lstdefinelanguage{GLSL}{
  alsoletter={0123456789},
  morekeywords=[1]{attribute,const,uniform,varying,buffer,shared,layout,
    centroid,flat,smooth,noperspective,in,out,inout,invariant,precise,discard,
    precision,highp,mediump,lowp,patch,sample,subroutine,coherent,volatile,restrict,
    readonly,writeonly,atomic_uint,asm,void,int,uint,float,bool,double,
    vec2,vec3,vec4,ivec2,ivec3,ivec4,bvec2,bvec3,bvec4,uvec2,uvec3,uvec4,dvec2,dvec3,dvec4,
    mat2,mat3,mat4,mat2x2,mat2x3,mat2x4,mat3x2,mat3x3,mat3x4,mat4x2,mat4x3,mat4x4,
    sampler1D,sampler2D,sampler3D,samplerCube,sampler2DRect,sampler2DArray,
    sampler1DShadow,sampler2DShadow,samplerCubeShadow,samplerBuffer,
    sampler2DMS,sampler2DMSArray,isampler2D,isampler3D,usampler2D,usampler3D,
    image2D,image3D,imageCube,
    if,else,for,while,do,break,continue,return,switch,case,default,true,false},
  morekeywords=[2]{radians,degrees,sin,cos,tan,asin,acos,atan,sinh,cosh,tanh,
    pow,exp,log,exp2,log2,sqrt,inversesqrt,abs,sign,floor,trunc,round,roundEven,
    ceil,fract,mod,modf,min,max,clamp,mix,step,smoothstep,isnan,isinf,
    texture,textureSize,textureLod,textureProj,textureGrad,
    dFdx,dFdy,fwidth,fwidthFine,normalize,length,distance,dot,cross,
    outerProduct,transpose,inverse,determinant,faceforward,reflect,refract,
    packHalf2x16,unpackHalf2x16,atomicAdd,atomicExchange,atomicCompSwap},
  keywordstyle=[1]\color{cppkeyword}\bfseries,
  keywordstyle=[2]\color{purple!70!black},
  sensitive=true,
  morecomment=[l]{//},
  morecomment=[s]{/*}{*/},
  morestring=[b]"
}
\lstdefinestyle{glslstyle}{
  style=codebase,
  language=GLSL,
}

\lstdefinelanguage{HLSL}{
  alsoletter={0123456789_},
  morekeywords=[1]{cbuffer,tbuffer,Buffer,Texture1D,Texture2D,Texture3D,TextureCube,
    Texture2DArray,Texture2DMS,StructuredBuffer,ByteBuffer,RWBuffer,RWStructuredBuffer,
    RWByteBuffer,AppendStructuredBuffer,ConsumeStructuredBuffer,RasterizerState,
    SamplerState,SamplerComparisonState,constantbuffer,groupshared,
    void,bool,int,uint,min16int,min16uint,float,half,double,dword,
    vector,matrix,bool2,bool3,bool4,int2,int3,int4,uint2,uint3,uint4,
    float2,float3,float4,half2,half3,half4,double2,double3,double4,
    float2x2,float3x3,float4x4,float4x3,float3x4,float2x4,float4x2,
    snorm,unorm,precise,linear,sampler,sampler1D,sampler2D,sampler3D,sampler_state,
    if,else,for,while,do,break,continue,return,switch,case,default,true,false,discard,
    in,out,inout,uniform,static,extern,noinline,payload},
  morekeywords=[2]{clip,killed,lerp,step,smoothstep,saturate,round,frac,rsqrt,isnan,
    sin,cos,tan,asin,acos,atan,atan2,sinh,cosh,tanh,pow,exp,log,log2,exp2,sqrt,
    abs,floor,ceil,min,max,clamp,sign,modf,distance,length,dot,cross,normalize,
    reflect,refract,faceforward,mul,transpose,determinant,
    tex2D,tex3D,texCUBE,tex2Dproj,tex2Dlod,Sample,SampleLevel,SampleBias,
    Load,GetDimensions,CalculateLevelOfDetail,
    ddx,ddy,ddx_fine,ddy_fine,width,AnalyticDerivativeGroupQuads,
    wavegetlaneid,wavelaneid,wavematch,waveshflget,waveactivecount,
    atomicadd,atomicand,atomicmax,atomicmin,atomicor,atomicexchange,interlockedadd},
  keywordstyle=[1]\color{cppkeyword}\bfseries,
  keywordstyle=[2]\color{purple!70!black},
  sensitive=false,
  morecomment=[l]{//},
  morecomment=[s]{/*}{*/},
  morestring=[b]"
}
\lstdefinestyle{hlslstyle}{
  style=codebase,
  language=HLSL,
}

% ---- Python：着色器变体批处理、SPIR-V 分析与性能数据脚本 ----
\lstdefinestyle{pythonstyle}{
  style=codebase,
  language=Python,
  morekeywords={self,lambda,async,await,yield,with,as,pass,del,
    global,nonlocal,is,in,not,and,or,True,False,None},
  keywordstyle=\color{cppkeyword}\bfseries,
  emph={np,plt,numpy,matplotlib,json,os,sys,pathlib,struct,subprocess,
    trimesh,pyrender,spirv_dis,re},
  emphstyle=\color{teal!80!black},
}

% ---- 统一 diff：补丁式改动对比 ----
\lstdefinestyle{diffstyle}{
  style=codebase,
  numbers=none,
  frame=none,
  backgroundcolor=\color{white},
  morecomment=[l][\color{diffdel}]{-},
  morecomment=[l][\color{diffadd}]{+},
  morecomment=[l][\color{blue!70!black}\bfseries]{@@},
}

% ==================== tcolorbox ====================
\usepackage[most]{tcolorbox}

\newtcolorbox{guideline}[1][]{
  enhanced,breakable,
  colback=blue!5!white,colframe=blue!75!black,
  fonttitle=\bfseries,title=要点,#1,
  boxed title style={colback=blue!75!black},
  attach boxed title to top left={yshift=-2mm,xshift=2mm},
  arc=2mm,boxrule=0.8mm,
}
\newtcolorbox{compare}[1][]{
  enhanced,breakable,
  colback=green!3!white,colframe=green!50!black,
  fonttitle=\bfseries,title=对比,#1,
  boxed title style={colback=green!50!black},
  attach boxed title to top left={yshift=-2mm,xshift=2mm},
  arc=2mm,boxrule=0.8mm,
}
\newtcolorbox{warning}[1][]{
  enhanced,breakable,
  colback=red!5!white,colframe=red!75!black,
  fonttitle=\bfseries,title=常见陷阱,#1,
  boxed title style={colback=red!75!black},
  attach boxed title to top left={yshift=-2mm,xshift=2mm},
  arc=2mm,boxrule=0.8mm,
}
\newtcolorbox{note}[1][]{
  enhanced,breakable,
  colback=orange!5!white,colframe=orange!80!black,
  fonttitle=\bfseries,title=注意,#1,
  boxed title style={colback=orange!80!black},
  attach boxed title to top left={yshift=-2mm,xshift=2mm},
  arc=2mm,boxrule=0.8mm,
}
\newtcolorbox{bestpractice}[1][]{
  enhanced,breakable,
  colback=purple!5!white,colframe=purple!60!black,
  fonttitle=\bfseries,title=最佳实践,#1,
  boxed title style={colback=purple!60!black},
  attach boxed title to top left={yshift=-2mm,xshift=2mm},
  arc=2mm,boxrule=0.8mm,
}

% ==================== 章节标题 ====================
\usepackage{titlesec}
\usepackage{fancyhdr}
\usepackage{graphicx}
\usepackage{amssymb}
\usepackage{amsmath}
\usepackage{tabularx}
\usepackage{booktabs}
\usepackage{longtable}
\usepackage{array}
\usepackage{multirow}
\usepackage{enumitem}
\usepackage{seqsplit}
\usepackage{float}

\titleformat{\chapter}[display]
  {\normalfont\huge\bfseries\color{toolheadbg}}
  {\CTEXthechapter}{20pt}{\Huge}
\titleformat{\section}
  {\normalfont\Large\bfseries\color{blue!70!black}}
  {\thesection}{1em}{}
\titleformat{\subsection}
  {\normalfont\large\bfseries\color{blue!60!black}}
  {\thesubsection}{1em}{}

\pagestyle{fancy}
\fancyhf{}
\fancyhead[LE,RO]{\thepage}
\fancyhead[RE]{\leftmark}
\fancyhead[LO]{\rightmark}
\setlength{\headheight}{14.5pt}

\setlist[itemize]{leftmargin=2em,itemsep=2pt}
\setlist[enumerate]{leftmargin=2em,itemsep=2pt}
\emergencystretch=6em

% ==================== 超链接 ====================
\usepackage{hyperref}
\hypersetup{
  colorlinks=true,
  linkcolor=blue!70!black,
  urlcolor=blue!60!black,
  bookmarks=true,
  pdfauthor={C++ Reference},
  pdftitle={C/C++ 渲染与着色器参考手册：GPU 管线、着色语言与图形 API},
}

% ==================== 自定义命令 ====================
\newcommand{\cpp}[1]{C++#1}
\newcommand{\Fn}[1]{\texttt{#1()}}
\newcommand{\Tp}[1]{\texttt{#1}}
\newcommand{\Opt}[1]{\texttt{-\/-#1}}
\newcommand{\Bk}{\discretionary{}{}{}}
\newcommand{\Lib}[1]{\textbf{\texttt{#1}}}
\newcommand{\Lang}[1]{\textsf{#1}}
\newcommand{\CL}{\Lang{C}}
\newcommand{\CPP}{\Lang{C++}}
\newcommand{\PY}{\Lang{Python}}
% 域标签：本书讲"两侧对比"时统一用这两个写法，避免同义多形
\newcommand{\DomCPU}{\Lang{CPU}}
\newcommand{\DomGPU}{\Lang{GPU}}
% 着色语言与中间表示缩写：正文、表头与图注统一用这四个，不写全拼或变形
\newcommand{\GSL}{\Lang{GLSL}}
\newcommand{\HSL}{\Lang{HLSL}}
\newcommand{\MSL}{\Lang{MSL}}
\newcommand{\SPIV}{\Lang{SPIR-V}}

% ====================================================================
% book 双页模式默认 \flushbottom，会靠拉伸行距把每页撑满；本仓库框体与长表多，
% 撑页产生的 Underfull \vbox 是排版噪声而非缺陷，这里统一改用小底部对齐。
\raggedbottom

\begin{document}

\frontmatter

\begin{titlepage}
\centering
\vspace*{4cm}
{\Huge\bfseries C/C++ 渲染与着色器参考手册\par}
\vspace{1cm}
{\LARGE GPU 绘图管线、着色语言与图形 API\par}
\vspace{0.4cm}
{\Large OpenGL、Vulkan、Direct3D 与 Metal 之上的 GLSL、HLSL、SPIR-V 接口\par}
\vfill
{\large \today\par}
\end{titlepage}

\tableofcontents

\mainmatter

\cleardoublepage

% ====================================================================
% 章节正文由 frag_chNN.tex 合并插入到下面这行之前（merge_frags.py）

% ============================================================
% ---- frag_ch01.tex ----
\part{为什么渲染的热路径是 C/C++ 与 GPU}
\chapter{分界与成因：CPU 提交与 GPU 执行}
\label{ch:ch01-boundary}

渲染的写法和其它系统编程不一样，根源在于它同时受两套时间约束：
\DomCPU{} 侧的一段循环要在几十微秒内完成录制与提交，\DomGPU{} 侧的一串固定
阶段要在同一帧内把几百万个片元算完。两侧由一条硬件队列隔开，队列让双方并行，
也就让``某一侧省下的时间''未必等于``帧时间变短''。

本册从这条缝出发。本章先把边界画清楚：帧预算怎么来、队列模型为什么把延迟
换成吞吐、\DomCPU{} 热路径为什么塞不下垃圾回收与无界锁、着色器为什么是
``第三种语言''、一次 draw call 的钱究竟花在哪里。第
\ref{ch:ch02-raster-pipeline} 章沿 \DomGPU{} 侧把管线阶段逐个走一遍，
后面各章进入输入装配、变换空间、三种着色语言与 \SPIV{} 工具链
（第 \ref{ch:ch03-vertex-input} 至 \ref{ch:ch07-spirv-toolchain} 章）。

\section{帧预算：由刷新率倒推的硬上限}

帧预算不是``性能目标''，而是由显示刷新率倒推出来的截止期限。常见的几个数量级：

\begin{itemize}
  \item 60 Hz：每帧约 16.7 ms；
  \item 90 Hz（一体机与头显常见）：约 11.1 ms，头显对单帧超时的容忍度更低；
  \item 120 Hz：约 8.3 ms；144 Hz：约 6.9 ms；240 Hz：约 4.2 ms。
\end{itemize}

刷新率翻倍并不只是``预算减半''。减半意味着原本能被平均掉的尖峰不再被平均掉：
一帧里一次几毫秒的堆分配、一次纹理上传、一次着色器模块的即时编译，在 60 Hz 下
只是抖动，在 120 Hz 下就是丢帧。所以渲染层考核的从来不是平均帧时，而是分布的
尾部与帧间隔的一致性（通常叫 frame pacing）。

预算还要在三方之间分配。\DomCPU{} 侧拿到``录制加提交''的时间，\DomGPU{} 侧拿到
``执行完本帧全部命令''的时间，剩下的由合成与扫描输出占据。双缓冲与三缓冲让
这三段\emph{重叠}，因此判据不是两侧时间相加小于 16.7 ms，而是同一帧的两段
各自不越界、并且队列不排空。把两条曲线分开看，是本章所有测量的前提。

\begin{warning}[title={只看一个总帧时会误诊}]
总帧时 16 ms 里 \DomCPU{} 用 12 ms 与 \DomGPU{} 用 12 ms，优化手段完全不同：
前者要减录制与提交开销，后者要减片元数量与带宽。在瓶颈之外的一侧优化，
只会让另一侧的余量继续被浪费。
\end{warning}

\section{提交与执行之间的队列模型}

现代图形 API 的公共骨架是三段式：先在\emph{命令缓冲}（Vulkan 的
\Tp{VkCommandBuffer}、D3D12 的 command list）里把状态与绘制指令录下来，
再整段\emph{提交}到某个队列（\Fn{vkQueueSubmit} 配合 \Tp{VkSubmitInfo}，
或 D3D12 命令队列的 \Fn{ExecuteCommandLists}），最后由硬件按提交顺序执行。
三段之间没有``立刻返回结果''的语义：提交是异步的，写下去的命令缓冲
只是驱动环形缓冲里的一段字节。

双缓冲与三缓冲是在这个模型上做的取舍。两帧缓冲时，\DomCPU{} 想录制下一帧
就要等待即将显示的那张交换链图像被释放；三帧缓冲把等待点往后挪一帧，
两侧更不容易互相堵死，代价是最坏情况下输入到像素的延迟增加约一帧。
Vulkan 把选择权交给应用：FIFO 等价于等待垂直同步的队列行为，
mailbox 允许 \DomGPU{} 提前画完并只保留最新一帧，immediate 不等待但会撕裂。
三者都没有免费的午餐，工程上常见的组合是桌面用 mailbox、头显用 FIFO
加异步时间扭曲。

跨层同步只有两种基本工具：信号量与栅栏。信号量表达队列之间的顺序
（上传纹理的计算命令没跑完，图形队列就不许读它）；栅栏表达 \DomGPU{}
已经跑到哪里这一事实。栅栏的成本是一次显式等待，而真正贵的是\emph{为了拿到
数据而等待}：时间戳查询、像素读回、统计查询都要等 \DomGPU{} 追上，
所以这类采样一律做成``晚一到两帧收割''。

\begin{note}[title={帧内不许阻塞的读回三件套}]
时间戳查询、\Fn{vkMapMemory} 之后立刻读、\Fn{vkWaitForFences} 都会打断流水线。
把它们挪到帧尾并允许结果过期一帧，是绝大多数引擎 \DomCPU{} 侧唯一能接受的形态；
需要即时决定本帧画不画的读回，应改用 \DomGPU{} 侧的条件渲染或间接绘制。
\end{note}

下面这段骨架给出可复用的形状：录制时按帧索引写入两个时间戳，提交之后
在\emph{下一帧}非阻塞地收割，任何一次收割失败只代表这一帧没有样本，
不引入等待。

\begin{lstlisting}[style=cppstyle,caption={主循环帧预算与 GPU 时间戳的非阻塞读回骨架},label={lst:ch01-frame-loop}]
// 每个槽位记录一帧的三段计时；kRing 与交换链图像数同阶
struct FrameTiming {
    uint64_t cpu_begin_ns = 0;
    uint64_t cpu_submit_ns = 0;
    uint64_t gpu_begin = 0;      // GPU tick，不是纳秒
    uint64_t gpu_end   = 0;
    bool     valid     = false;
};

static constexpr int kRing = 3;
static FrameTiming   g_ring[kRing];
static int           g_head = 0;

void begin_frame(FrameContext& fc) {
    fc.frame_begin_ns = fc.clock.now_ns();
    const double period_ns = 1e9 / static_cast<double>(fc.refresh_hz);
    // 预算 = 帧周期 - 合成与扫描输出占用的估计值
    fc.frame_deadline_ns = fc.frame_begin_ns
        + static_cast<uint64_t>(period_ns) - fc.present_cost_ns;
}

void end_frame(FrameContext& fc) {
    // 先打尾部时间戳，再结束录制，最后提交
    vkCmdWriteTimestamp(fc.cmd, VK_PIPELINE_STAGE_BOTTOM_OF_PIPE_BIT,
                        fc.query_pool, 2 * g_head + 1);
    vkEndCommandBuffer(fc.cmd);

    const VkSubmitInfo si = { VK_STRUCTURE_TYPE_SUBMIT_INFO };
    fc.last_submit_ns = fc.clock.now_ns();
    vkQueueSubmit(fc.queue, 1, &si, fc.in_flight[g_head]);
    g_head = (g_head + 1) % kRing;

    const uint64_t used = fc.clock.now_ns() - fc.frame_begin_ns;
    fc.hud.late_cpu = (used > fc.frame_deadline_ns) ? 1 : 0;
    fc.hud.cpu_ms   = used * 1e-6;
}

// 下一帧开头调用：非阻塞收割，收不到就跳过
void harvest(FrameContext& fc) {
    const int prev = (g_head + kRing - 1) % kRing;
    uint64_t raw[4] = {0, 0, 0, 0};
    const VkResult r = vkGetQueryPoolResults(
        fc.device, fc.query_pool, 2 * prev, 2, sizeof(raw), raw,
        sizeof(uint64_t) * 2,
        VK_QUERY_RESULT_64_BIT | VK_QUERY_RESULT_WITH_AVAILABILITY_BIT);
    if (r != VK_SUCCESS || raw[1] == 0 || raw[3] == 0) {
        return;                     // 尚未落地，本帧不采样也不等
    }
    const double ns_per_tick = fc.limits.timestampPeriod;
    g_ring[prev].gpu_begin = raw[0];
    g_ring[prev].gpu_end   = raw[2];
    g_ring[prev].valid     = true;
    fc.hud.gpu_ms = (raw[2] - raw[0]) * ns_per_tick * 1e-6;
}
\end{lstlisting}

\Tp{timestampPeriod} 的含义是``每 tick 多少纳秒''，来自物理设备限制，
不同 \DomGPU{} 的刻度与精度差别很大，因此时间戳只适合在同一设备上
比较\emph{相对}长短，不适合跨设备对齐绝对时间。

\section{\DomCPU{} 侧热路径为什么不容下垃圾回收、反射与无界锁}

录制与提交这段代码有三个特征：执行频率高（每帧每对象一遍）、
单次工作量小、超时即丢失可见结果。三条合起来排除了三类在别处很合理的机制。

\textbf{垃圾回收。} 回收器的工作量由堆的历史决定，不由本帧决定，
于是它的暂停会精确地落在最不该出现尖峰的地方。渲染层的对象生命周期又极不均匀：
网格与纹理活得很稳，临时矩阵与调试字符串活得极短，这正好让分代假设失效。
替代方案是按域切开分配：帧内失效的走帧作用域线性资源
（\Tp{std::pmr} 的单调缓冲区即可），跨帧的走显式设备资源与句柄表。

\textbf{反射。} 真实成本不是``慢一点''，而是它把类型信息推到运行期，
从而切断内联与常量传播。绘制列表若以虚接口逐个遍历对象，编译器无法合并材质
查表与状态比较这些小额操作，也无法在 \cpp{20} 里用编译期表把已知组合预先展开。
主流做法是把``可变的类型''降级成``不变的索引''：枚举加句柄加变体表，
运行期只做整数比较。需要灵活性时，把灵活性放进数据而不是放进代码。

\textbf{无界锁。} 互斥锁的持有时长由持锁那方的逻辑决定，而提交线程需要
``任何时刻都能在微秒内拿到''。一旦后台线程在锁里做了磁盘访问或大块分配，
帧就丢了；优先级反转会让情况更糟。可行的边界是：录制阶段按线程切分命令池
（Vulkan 要求同一 \Tp{VkCommandPool} 同一时刻只服务一个线程的录制，
这恰好把竞争挡在门外），共享数据以帧为单位做成只读快照，
跨线程传递只走有界队列或原子索引，让等待时长在构造期就被上限约束住。

\section{着色器是第三种语言：两侧类型系统与编译期特化}

着色器语言既不是 \CL{} 也不是 \CPP{}，它跨在两侧中间，最好当成
``第三种语言''来学。差异在四个地方。

\textbf{类型系统不一致。} \GSL{} 的 \Tp{mat4} 是列主序的四个 \Tp{vec4}，
\Tp{vec3} 在 \Tp{std140} 布局里要占满 16 字节；\CPP{} 侧的结构体不会自动补这些
空洞，忘记对齐就会得到平移或缩放错乱的几何。\HSL{} 更直接：
常量缓冲区里的 \Tp{float3} 由寄存器打包规则分配，
与 \CPP{} 端的内存布局不是同一条规则。这类错误不崩，只画错，
所以两侧声明必须逐字段对照，并让结构体尺寸成为断言。

\begin{lstlisting}[style=cppstyle,caption={与着色器侧逐字段对应的顶点与常量缓冲声明},label={lst:ch01-cpu-side}]
// location 编号与顶点输入描述符一致，着色器按同序声明
struct alignas(16) VertexGpu {
    float pos[3];       // location 0：xyz，w 由着色器补 1.0
    float nrm[3];       // location 1
    float uv[2];        // location 2
};

// std140：vec3 也要占满 16 字节，mat4 按列展开为 64 字节
struct alignas(16) CameraBlock {
    float view_proj[16];
    float eye_world[4];   // 第 4 分量仅用于对齐
    float exposure;
    float tail_[3];       // 尾部补齐到 16 的倍数
};
static_assert(sizeof(CameraBlock) == 96, "layout must match GLSL");
\end{lstlisting}

\begin{lstlisting}[style=glslstyle,caption={与上例逐字段呼应的最小顶点与片元着色器},label={lst:ch01-glsl-pair}]
#version 450
// ---- 顶点阶段：属性 location 与 CPU 侧声明一一对应 ----
layout(location = 0) in vec3 a_position;
layout(location = 1) in vec3 a_normal;
layout(location = 2) in vec2 a_uv;

layout(std140, binding = 0) uniform Camera {
    mat4  view_proj;
    vec4  eye_world;
    float exposure;
};

layout(location = 0) out vec3 v_world;
layout(location = 1) out vec3 v_normal;
layout(location = 2) out vec2 v_uv;

void main() {
    v_world  = a_position;
    v_normal = a_normal;
    v_uv     = a_uv;
    gl_Position = view_proj * vec4(a_position, 1.0);
}
\end{lstlisting}

\begin{lstlisting}[style=glslstyle,caption={同一个常量块在片元侧必须原样再声明一遍},label={lst:ch01-glsl-fs}]
#version 450
layout(std140, binding = 0) uniform Camera {
    mat4  view_proj;
    vec4  eye_world;
    float exposure;
};
layout(binding = 1) uniform sampler2D albedo;

layout(location = 0) in  vec3 v_world;
layout(location = 1) in  vec3 v_normal;
layout(location = 2) in  vec2 v_uv;
layout(location = 0) out vec4 out_color;

void main() {
    vec3  n    = normalize(v_normal);
    vec3  e    = normalize(eye_world.xyz - v_world);
    float ndl  = max(dot(n, vec3(0.0, 1.0, 0.0)), 0.0);
    float fres = pow(1.0 - max(dot(n, e), 0.0), 5.0);
    vec3  base = texture(albedo, v_uv).rgb;
    out_color = vec4((base * (ndl + 0.15) + fres * 0.05) * exposure, 1.0);
}
\end{lstlisting}

\textbf{编译期特化。} 两侧的连接点在编译期就已发生：同一份源码要按光照数量、
阴影级联数、是否透明、目标 API 生成若干变体。常见手段是宏加常量特化；
\GSL{} 与 \SPIV{} 也支持 specialization constant，让驱动在管线创建期把常量代入。
无论走哪条路，结果都是``着色器数量按维度相乘''，变体表因此成为工程问题，
第 \ref{ch:ch07-spirv-toolchain} 章会把它当作构建系统来讨论。

\textbf{常量更新的代价是结构性的。} 每个 uniform 的写入都是一次跨层状态更新：
它要变成可见，就得落进 \DomGPU{} 读得到的内存。旧式逐次 uniform 上传把这件事
放在每次绘制之前，于是``每 draw call 若干次上传''成为默认成本；
现代 API 把它整理成绑定模型——一组资源打包进描述符集或根签名，
少量小数据走 push constant（Vulkan 保证的容量下限是 128 字节，超出就得换缓冲区），
大批量数据进结构化缓冲区由索引寻址。换绑定比换数据便宜，
是这类 API 设计的基本取向。

\textbf{分支与一致执行。} \DomGPU{} 按组（wave 或 subgroup）推进一批像素，
组内分支不一致时两条路都要走，这就是分歧的成本。片元着色器里的 \Fn{dFdx}
与 \Fn{fwidth} 一类求导依赖同组邻居的值，所以跨三角形边界取值不可靠，
出现 \Tp{discard} 之后更容易出问题。共享变量在计算着色器里要显式下屏障
（\Fn{memoryBarrierShared}），不能照搬 \CPP{} 侧的原子习惯。
把这些当成语言的一部分来学，比事后靠 profiler 反推便宜得多。

\begin{compare}[title={同一次改一个数在两侧的代价}]
\DomCPU{} 侧改一个成员变量是一次 store，编译器还可能把它并入寄存器；
着色器侧改一个 uniform 要落到绑定内存，可能触发描述符写入与屏障，
而且对所有并发执行的图元同时生效。于是``每物体一个常数''天然是错的形状，
正确的形状是``每帧一块数组，物体只带一个索引''。\end{compare}

\section{一次 draw call 的成本结构}

``减少 draw call''这句话能流传几十年，是因为它压中了真实成本；
它也常被用错，因为它把四种成本混成了一种。拆开看一次绘制：

\begin{enumerate}
  \item \textbf{用户态校验与录制。} API 调用先做参数与状态合法性检查，
        再往命令缓冲写机器可执行的形式；这是 \DomCPU{} 时间，随调用次数线性增长。
  \item \textbf{提交粒度。} 命令要攒到一定量或遇到同步点才真正交给驱动；
        频繁切换状态会让驱动提前冲刷，冲刷本身就是成本。
  \item \textbf{状态切换。} 换管线对象等于换着色器二进制与固定功能配置，
        换绑定等于换描述符表或顶点输入状态。驱动常常要靠重新配置甚至打补丁
        来实现切换，切换次数往往比切换本身更贵。
  \item \textbf{执行与带宽。} 顶点抓取、图元处理、片元着色与写回，
        这部分与 draw call 个数几乎无关，只与几何量和像素量有关。
\end{enumerate}

批处理省的是前三种：把 $N$ 次录制合成一次，把 $N$ 次状态切换降成一次，
同时因为顶点数据连续、索引可复用而改善缓存命中。实例化省得更彻底，一次调用
画多个副本，只把逐实例属性做成缓冲区。合并静态几何省的是切换，代价是失去
逐个剔除与逐段排序的粒度，可能反而多画像素。

于是判据很清楚：瓶颈在 \DomCPU{}（\DomGPU{} 利用率偏低、帧时间随物体数量
线性上涨、录制端贴近预算）时，批处理与 \DomGPU{} 驱动的命令生成是对症的药；
瓶颈在 \DomGPU{}（\DomCPU{} 早早结束而 \DomGPU{} 溢出预算）时，减 draw call
不会让帧变快，该动的是像素数量、着色器指令与深度写入顺序。第
\ref{ch:ch08-compute-driven} 章讨论间接绘制如何把第一类成本压到极低，第
\ref{ch:ch14-perf-tools} 章给出区分这两类的测量流程。

\section{本册路线与三种读者视角}

三侧的责任划分与``怎么算合格''值得写在最前面，后面各章的每一处优化都能
回到这张表上定位。

\begin{longtable}{>{\raggedright\arraybackslash}p{0.15\textwidth}
                  >{\raggedright\arraybackslash}p{0.23\textwidth}
                  >{\raggedright\arraybackslash}p{0.23\textwidth}
                  >{\raggedright\arraybackslash}p{0.19\textwidth}}
\caption{\DomCPU{}、\DomGPU{} 与脚本层三侧的职责与验收标准}
\label{tab:ch01-three-domains} \\
\toprule
维度 & \DomCPU{} 侧 & \DomGPU{} 侧 & 脚本层 \\
\midrule
\endfirsthead
\toprule
维度 & \DomCPU{} 侧 & \DomGPU{} 侧 & 脚本层 \\
\midrule
\endhead
\midrule
\endfoot
\bottomrule
\endlastfoot
职责 & 把剔除结果整理成可见列表，按状态排序，录制并按预算提交 &
执行命令：顶点变换、光栅化、片元求值、深度与模板测试、混合写回 &
资源导入与变体生成，把材质图求值成着色器源码 \\
每帧成本 & 与可见对象数成正比，须稳定在预算内且不出现尖峰 &
与像素数、带宽和着色器工作量成正比，须不越预算 &
不参与帧内计时，只影响加载时长与变体数量 \\
失败症状 & 帧时间长尾、\DomGPU{} 利用率偏低、录制端已经超时 &
帧时间随分辨率线性增长，降低物体数无效，片元着色器成为热点 &
冷启动慢、变体爆炸、首次命中某材质时卡顿 \\
验收标准 & 最坏帧的 \DomCPU{} 段有稳定余量，热路径每帧零分配，无阻塞读回 &
\DomGPU{} 段不越预算，状态切换与冲刷次数受控 &
变体总数可枚举，预热覆盖首帧会用到的全部材质 \\
\end{longtable}

三种读者的入口不同，本册按同一顺序编号以便交叉引用。
\textbf{引擎作者}关心接口形状与稳定性，优先读第
\ref{ch:ch09-vulkan} 与 \ref{ch:ch10-d3d12-metal} 章，
把资源生命周期与管线对象创建当主战场。
\textbf{图形程序员}关心阶段与状态的正确性，第
\ref{ch:ch11-gl-legacy} 与 \ref{ch:ch12-lighting-pbr} 章是主线。
\textbf{算法工程师}想把新的光照与求交方法落到实时预算内，
从第 \ref{ch:ch13-raytracing} 与 \ref{ch:ch08-compute-driven} 章切入，
再看 \SPIV{} 一章理解算法如何变成可分发的二进制；库与工具的取舍见附录第
\ref{ap:apC-library-choice} 节。

\begin{guideline}[title={本章的五条判据}]
一、帧预算来自刷新率，考核的是尾部而非平均，两侧曲线要分开看。
二、提交是异步的，任何同步等待都挪到帧尾并允许结果过期一帧。
三、\DomCPU{} 热路径的合格线是零分配、零反射、锁有构造期上限。
四、着色器是第三种语言：两侧类型逐字段对照，变体与绑定形状是设计决定。
五、先分辨瓶颈在哪一侧，再决定是减 draw call 还是减像素。
\end{guideline}

% ============================================================
% ---- frag_ch02.tex ----
\chapter{GPU 绘图管线解剖}
\label{ch:ch02-raster-pipeline}

把 \DomGPU{} 侧想成一条装配线是最有用的心智模型：每个工位只做一件事，
工位之间由缓冲隔开，快慢取决于最慢的一段，而插错一个工位会让整条线停下来。
本章按数据流走一遍：输入装配、顶点着色、可选的细分与几何、剔除与裁剪、
屏幕映射、光栅化与插值、片元着色、深度与模板测试、混合、写回。
每一步都问三个问题：这一段的输入输出是什么，它是可编程的还是固定约定的，
它出错时画面长什么样。

\section{分段的理由与可编程边界}

阶段边界不由语言决定，由\emph{数据交接形式}决定：顶点阶段交出裁剪空间坐标，
图元装配交出图元，光栅化交出带覆盖率的片元，片元阶段交出颜色，
输出合并段负责把它和已有内容合并。边界处没有函数调用，
只有硬件级的一次入队，所以一段里的停顿不会直接变成另一段的停顿，
但一段的产出量会放大下一段的成本。

可编程的阶段是顶点、（可选的）细分控制与细分评测、（可选的）几何、片元，
以及不属于光栅化路径但共享同一套资源绑定的计算。
固定的约定是索引解码、视锥与背面剔除、近平面裁剪、透视除法与视口映射、
三角形覆盖判定、属性插值、深度与模板测试、混合、格式转换与写回。
固定不等于无成本：正因为不可编程，这些段只能通过\emph{喂给它什么数据}
来优化，这也是为什么同一个着色器在不同 API 上表现不同。

抓帧工具是理解这条线最省事的方式：它把每次提交展开成状态与阶段的序列，
直接告诉你某一段被改了 \Tp{VkViewport} 与 \Tp{VkRect2D} 的哪一个。

\begin{lstlisting}[style=shellstyle,caption={用命令行确认设备能力与着色器二进制，再交给抓帧工具逐阶段查看},label={lst:ch02-capture}]
#!/usr/bin/env bash
# 1) 先确认设备与扩展：时间戳、校准时间戳决定能不能计时
vulkaninfo --summary > gpu.txt
grep -E 'timestampPeriod|VK_EXT_calibrated' gpu.txt

# 2) 校验与反汇编 SPIR-V：看绑定、特化常量与实际生成的指令
spirv-val demo.vert.spv
spirv-dis demo.frag.spv | grep -E 'gl_FragCoord|dFdx'

# 3) 抓一帧，交给图形调试器逐阶段回放
renderdoccmd capture ./build/demo
\end{lstlisting}

\section{输入装配、顶点阶段与可选阶段}

索引解码把一维缓冲区解释成图元序列，代价与索引数量成正比；
重复索引命中顶点缓存时可以省一次抓取，这也是``按缓存局部性重排索引''
（vertex cache optimization）的全部理由。图元类型决定重复率：
三角列表几何独立、条带省索引但状态更敏感；条带里的重启标记让一段缓冲区
画出不连续的图元，代价是驱动需要额外检查。

顶点着色器对每个顶点跑一次，产出裁剪空间坐标（\GSL{} 的 \Tp{gl\_Position}、
\HSL{} 的 \Tp{SV\_Position}）与若干待插值属性。三个容易忽略的事实：

\begin{itemize}
  \item 属性个数与宽度直接变成带宽，位置之外的每个属性都要过一遍整个管线，
        并且\emph{逐片元}求值。能延后到片元阶段算的量就不要传。
  \item 插值方式是可以指定的约定：\GSL{} 用 \Tp{flat}、\Tp{centroid}、
        \Tp{noperspective}、\Tp{sample} 四个限定符改变插值语义，
        \HSL{} 用 \Tp{CENTROID}、\Tp{noperspective} 与 \Tp{linear} 一类
        插值修饰符。传整数标识符时忘了 \Tp{flat}，
        得到的就是两个整数值之间被插出来的第三种值。
  \item 语义名不是名字。同一个 \Tp{out} 在不同阶段要靠编号对上，
        编号错位不会报错，只会把法线当颜色用。
\end{itemize}

顶点阶段之后还有两个可选阶段，它们在实时管线里正在衰落。
细分（tessellation）在曲面补丁上按需生成顶点，
适合地形、曲面与程序化细节；它的成本是评测阶段的输出会成倍放大顶点数，
并且细分级别受设备上限与图元数量约束，很容易从``补细节''变成``补带宽''。
在 \DomGPU{} 硬件里，细分阶段位于顶点之后、裁剪之前，
它让``一帧生成多少顶点''变成运行期未知，进而削弱了后续的调度可预测性。

几何着色器允许一个图元产出多个图元，历史上被用来做实例展开、
体积渲染与阴影级联分发。它在实时管线里衰落有三个原因：
一是它在光栅化之前工作，很多实现因此把该 draw 的深度测试推到较晚的位置，
early-Z 的收益随之消失；二是输出量不可预知，硬件要为最坏情况预留缓冲，
容易撞上图元处理与带宽上限；三是它与间接绘制的组合在部分实现上代价明显。
关键的一点是接口本身的选择：Vulkan 没有提供几何着色器阶段，
D3D12 也基本放弃了它，转而提供计算与网格着色器这条更可控的路。
本册的处理方式因此是：把过去用几何着色器做的事挪到计算阶段，
见第 \ref{ch:ch08-compute-driven} 章。

\section{裁剪、屏幕映射、光栅化与插值}

裁剪发生在齐次坐标上：与视锥六个半平面求交，跨越近平面的三角形要被切割并
生成新顶点，这一步是固定功能，也是为什么把物体整体推到近平面之后
比让裁剪器逐三角形处理更划算。之后是透视除法得到设备坐标，
再由视口映射换算成屏幕坐标。这里有三个跨 API 的经典坑：
\DomGPU{} 各家对设备坐标的 $Y$ 方向与深度区间约定不同，
同一段矩阵代码在 OpenGL 约定下需要翻转 $Y$，
而深度范围从 $[-1,1]$ 映射到 $[0,1]$ 的方式也直接影响
可逆裁剪与深度精度分配（第 \ref{ch:ch04-transform-space} 章展开）。

剔除分两级：图元级由三角形朝向判定背面，靠带符号面积，
因此绕序与前面的矩阵共同决定正确性，镜像缩放会把内外翻转；
遮挡剔除则通常由应用自己做（层级可见性、硬件遮挡查询），
因为硬件的保守剔除只会丢掉完全在视口外或背面朝向的图元。
剪刀矩形与小的多边形剔除是在这一段生效的粗粒度开关。

光栅化决定哪些像素被覆盖，覆盖判定依赖像素中心约定，
而各家约定历史上不同：填充规则与深度、模板采样点的位置可以不是同一个，
于是共享一套几何在不同 API 上会看到边线像素的取舍差异。
覆盖率随后决定插值权重。透视正确的插值要求属性除以 $w$ 后再插值，
这是硬件做的，但\emph{做在哪个采样模式上}由你选的 MSAA 与质心约定决定。

求导是这一段最重要的隐藏约定。\GSL{} 的 \Fn{dFdx}、\Fn{dFdy}、\Fn{fwidth}
与 \HSL{} 的 \Fn{ddx}、\Fn{ddy} 都在 quad（$2\times2$ 个像素）粒度上工作，
它们需要邻居的值才成立，因此三个后果要记住：跨三角形边缘取值不可靠，
\Tp{discard} 之后不可靠，出现在非一致控制流里结果未定义。
纹理自动选级就是这么来的，写调试着色器时也常常被这么用：

\begin{lstlisting}[style=glslstyle,caption={用屏幕导数估面法线并用 gl\_FragCoord 叠加调试图案},label={lst:ch02-glsl-deriv}]
#version 450
layout(location = 0) in  vec3 v_world;
layout(location = 1) in  vec2 v_uv;
layout(location = 0) out vec4 out_color;

layout(binding = 2) uniform sampler2D albedo;

void main() {
    // 1) 导数必须在任何 discard 之前求，且要落在顶层控制流
    vec3 dpdx = dFdx(v_world);
    vec3 dpdy = dFdy(v_world);
    vec3 face_n = normalize(cross(dpdx, dpdy));

    // 2) gl_FragCoord 的原点与方向由 API 决定，调试图案要按此翻转
    ivec2 pc = ivec2(gl_FragCoord.xy);
    float checker = mod(float(pc.x >> 3) + float(pc.y >> 3), 2.0);

    // 3) 用 UV 的变化量反推可见的纹理密度，决定用哪一级 mipmap
    float density = length(vec2(dFdx(v_uv.x), dFdy(v_uv.y)))
                  * float(textureSize(albedo, 0).x);
    float lod = clamp(log2(max(density, 1.0)), 0.0, 8.0);
    vec3  base = textureLod(albedo, v_uv, lod).rgb;

    out_color = vec4(mix(base, face_n * 0.5 + 0.5, 0.25 * checker), 1.0);
}
\end{lstlisting}

同一件事在 \HSL{} 里写出来如下。注意 \Tp{SV\_Position} 在片元阶段拿到的是
像素坐标而非裁剪坐标，而 \Tp{SV\_Target} 只是渲染目标的编号约定：

\begin{lstlisting}[style=hlslstyle,caption={用 ddx、ddy 现场构造切平空间的等价写法},label={lst:ch02-hlsl-tbn}]
Texture2D    g_albedo  : register(t0);
Texture2D    g_normal  : register(t1);
SamplerState g_linear  : register(s0);

struct PSIn {
    float4 pos   : SV_Position;   // 已是屏幕像素坐标
    float3 world : TEXCOORD0;
    float3 nrm   : TEXCOORD1;
    float2 uv    : TEXCOORD2;
};

float4 PSMain(PSIn i) : SV_Target {
    float4 base = g_albedo.Sample(g_linear, i.uv);

    // 导数只在顶层求，clip 放在后面，否则邻居被丢弃时结果未定义
    float3 dPdu = ddx(i.world);
    float3 dPdv = ddy(i.world);
    float2 dUdu = ddx(i.uv);
    float2 dUdv = ddy(i.uv);

    float det = dUdu.x * dUdv.y - dUdv.x * dUdu.y;
    float r = (abs(det) > 1e-8) ? (1.0 / det) : 0.0;
    float3 T = r * ( dUdv.y * dPdu - dUdu.y * dPdv);
    float3 B = r * (-dUdv.x * dPdu + dUdu.x * dPdv);

    float3 n = normalize(i.nrm);
    T = normalize(T - n * dot(T, n));      // 施密特正交化，保住法线方向
    B = cross(n, T);

    float3 map = normalize(g_normal.Sample(g_linear, i.uv).xyz * 2.0 - 1.0);
    float3 world_n = normalize(map.x * T + map.y * B + map.z * n);

    clip(base.a - 0.5);                     // 等价于 discard
    return float4(world_n * 0.5 + 0.5, 1.0);
}
\end{lstlisting}

\section{片元测试顺序：深度、模板、混合与写回}

片元阶段的输出要过三道检查才被写入：模板测试、深度测试、混合。
它们的顺序与可编程性构成性能的主战场。

\textbf{early-Z。} 很多实现能在跑片元着色器\emph{之前}完成深度测试，
把注定被遮住的片元直接丢掉，省下整段着色与采样开销。
但这个捷径有条件：着色器里写了深度输出、用了求导、
或者需要执行 \Tp{discard}，都可能让它退化成晚期测试。
所以工程惯例是：不透明几何按\emph{由前到后}排序、深度写入打开、
深度比较用``不大于''；透明几何反过来，由后到前且关闭深度写入。
这两条规则不是玄学，它们是让 early-Z 与混合的读取依赖同时成立。

\textbf{深度写入顺序。} 先写深度再画大批量片元，等于给自己发了一张
后续片元的免费通行证；反过来（例如先画半透明地面）会让后面每一层都完整着色。
延迟渲染里几何通道（G-buffer）之所以能吃满带宽，
正是因为它必须写满四到八个渲染目标且不能靠深度剔除。

\textbf{discard 与 alpha-to-coverage。} \Tp{discard} 让片元完全不写入，
适合叶片与栏杆一类的镂空，但它破坏求导、常常关掉 early-Z，
于是每像素只丢一点点的贴图为它付出的代价远大于收益。
更好的形状是 alpha-to-coverage：把覆盖概率换算成 MSAA 采样点的覆盖率，
镂空边缘随多重采样自然抗锯齿。它在混合阶段之前生效，
Vulkan 的多重采样状态创建信息里的 \Tp{alphaToCoverageEnable} 与
D3D12 混合描述里的 \Tp{AlphaToCoverageEnable} 表示同一件事。
用之前请确保渲染目标确实是多重采样的，单采样目标上它没有意义。

\textbf{混合。} 固定功能，读目标、算、写回，因此天然要求把目标读回来，
在带压缩与 tile 架构上都是带宽成本。混合关闭时写回最容易，
所以``能不混合就别混合''是真正的优化抓手；必须混合时，
先想能不能用 OIT 之外的形状（预乘 alpha、分离双混合、
把透明改成一个覆盖缓冲）绕开顺序依赖。

\textbf{多目标与写掩码。} 一次片元输出到多个渲染目标是最贵的输出形状，
通道数直接乘在带宽上；颜色写掩码可以让某几通道不写，
省流量但不省图元处理。输出格式（是否 \Tp{sRGB}、是否 \Tp{UNORM}）
在管线末端做转换，这类转换是固定的，改格式等于改硬件读写宽度。

\textbf{模板。} 模板测试的价值在于它能表达``几何自己声明的可见区域''，
常用于轮廓、遮罩与分块重建；它的读写位掩码约定和深度共用一个附件，
所以把深度与模板的分量分配清楚，比多加一个通道更省。

\section{两条实现路线：即时模式与基于 tile}

同一份规范可以由两类硬件实现。即时模式渲染（IMR）在片元一算完就直接读写
外部显存里的颜色与深度缓冲，逻辑简单，代价是每个片元两次往返内存，
于是 overdraw 与 MSAA 的开销几乎线性落到带宽上。
基于 tile 的渲染先把一帧按屏幕分块（bin），用一个前置的图元处理过程把
每个图元登记到它覆盖的块，然后逐块把颜色与深度放在片上内存里画完，
最后一次性写回（TBDR 还会在片上做隐藏面移除）。
两类实现解释了移动平台上三条看起来武断的建议：

\begin{itemize}
  \item tile 大小固定，频繁跨块与超大图元会增加片上溢出到内存的次数，
        因此大面积的低多边形覆盖很贵，这被称为 overdraw 之外的
        ``\emph{带宽式}''成本；
  \item 一次渲染过程中途改变渲染目标尺寸或反复切换通道，
        会迫使硬件提前 flush 当前块，flush 就是带宽；
  \item MSAA 与深度测试在片上很便宜，因为比较与覆盖率可以留在寄存器里，
        只有最终 resolve 才走内存，这也是移动端偏爱 MSAA 的原因。
\end{itemize}

需要澄清的是，这条界线不是黑白：桌面 \DomGPU{} 同样做分块与带宽压缩，
而 tile 架构也允许应用按即时模式的方式提交，只是付带宽的账。

MSAA 的两种成本要分清：\emph{覆盖}用多重采样（哪些像素被写入），
\emph{着色}默认仍是每像素一次；只有显式按采样求值
（\GSL{} 的 \Tp{sample} 限定与 \Tp{gl\_SampleID}、\HSL{} 的
\Tp{SV\_SampleIndex}）才变成逐采样开销，那是成倍的成本。
resolve 把多样本压成单样本：要么由硬件的解析附件在渲染过程结束时完成，
要么显式一次逐像素平均。在 tile 架构上前者几乎免费、后者一次带宽，
在 IMR 上两者都贵。这就是同一段后期处理代码在两类设备上
差距可达数倍的结构性原因。

\section{阶段小结与排障入口}

\begin{longtable}{>{\raggedright\arraybackslash}p{0.19\textwidth}
                  >{\raggedright\arraybackslash}p{0.12\textwidth}
                  >{\raggedright\arraybackslash}p{0.26\textwidth}
                  >{\raggedright\arraybackslash}p{0.23\textwidth}}
\caption{管线阶段、可编程性、常见失败模式与优化抓手}
\label{tab:ch02-stages} \\
\toprule
阶段 & 可编程 & 常见失败模式 & 优化抓手 \\
\midrule
\endfirsthead
\toprule
阶段 & 可编程 & 常见失败模式 & 优化抓手 \\
\midrule
\endhead
\midrule
\endfoot
\bottomrule
\endlastfoot
输入装配 & 固定 & 索引重复率高、重启标记用错导致图元串位 &
按缓存局部性重排索引，改用列表或更少的条带 \\
顶点着色 & 可编程 & 属性传多了、忘 \Tp{flat} 导致整数被插值 &
把可延迟的计算挪到片元，减少逐顶点属性宽度 \\
细分与几何 & 可编程但可选 & 输出量运行期未知，放大顶点与图元带宽 &
能用实例化或计算解决就别开这两个阶段 \\
剔除与裁剪 & 固定 & 镜像缩放翻转绕序，近平面切割产生大量新顶点 &
提前做层级可见性，整体推出视锥而不是逐图元裁剪 \\
屏幕映射 & 固定 & 跨 API 的 $Y$ 方向与深度区间约定不一致 &
在一处集中翻转，投影矩阵按目标约定生成 \\
光栅化与插值 & 固定 & 像素中心约定差异，超小图元被丢弃 &
控制三角形在屏幕上的尺寸分布，避免亚像素几何 \\
片元着色 & 可编程 & \Tp{discard} 与求导冲突，
逐片元纹理过多 & 让 early-Z 成立，先做深度预通道再画细节 \\
深度与模板 & 固定 & 深度写入顺序错，透明几何关掉 early-Z &
不透明由前到后写深度，透明由后到前关写入 \\
混合与写回 & 固定 & 混合读回目标造成带宽峰值，通道数翻倍 &
能不混合就不混合，减少渲染目标数量与位宽 \\
\end{longtable}

\begin{bestpractice}[title={按阶段定位问题的四步顺序}]
一、先看是不是根本没到这个阶段：剔除与可见性列表决定了后面一切的量。
二、再看像素数与通道数：片元与写回的开销和分辨率成正比，与物体数无关。
三、然后查绑定与状态切换次数，它落在提交与驱动一侧，表现为利用率低而非慢。
四、最后才优化着色器指令，因为指令级优化的收益上限由前三步决定。
\end{bestpractice}

% ============================================================
% ---- frag_ch03.tex ----
\chapter{顶点、属性与内存布局}
\label{ch:ch03-vertex-input}

\DomGPU\ 的顶点级不认识``对象''或``网格''，它只认三件事：一段字节缓冲、每多少个字节算一个
顶点、每个属性在这段字节里的偏移与类型。本章把这三件事讲透，因为它们同时决定另外两件容易
被低估的事：\DomCPU\ 侧的资产格式设计，以及着色器 uniform 块的二进制布局。
阶段划分见 \ref{ch:ch02-raster-pipeline}，矩阵与空间的约定见 \ref{ch:ch04-transform-space}。

\section{顶点缓冲与索引缓冲}

顶点缓冲存放属性，索引缓冲存放``顶点的下标''。
两者的分工决定两类完全不同的开销：属性缓冲的尺寸影响带宽与驻留，
索引缓冲的顺序影响 \DomGPU\ 上顶点着色器的执行次数。
一次索引绘制取第 \Tp{firstIndex} 起的若干下标，每个下标加上 \Tp{vertexOffset}，
再去属性缓冲里取对应那一行的字节。
\Fn{vkCmdDrawIndexed} 的 \Tp{vertexOffset} 是加在索引之上的有符号偏移，作用是让同一份索引
复用属性缓冲的不同区段，不必为每个物件重排索引。
索引宽度只有两种通用取值：\Tp{VK\_INDEX\_\Bk TYPE\_UINT16}（上限 65535 个顶点）与
\Tp{VK\_INDEX\_\Bk TYPE\_UINT32}，更窄的形式要靠扩展。
\Fn{glDrawElements} 走同一套逻辑，只是索引缓冲来自
\Tp{GL\_ELEMENT\_\Bk ARRAY\_BUFFER} 这个隐式绑定点，而非每次调用显式传入。
绑定索引缓冲时偏移必须是索引宽度的整数倍，否则是未定义行为。

索引的价值来自\emph{变换后缓存}（post-transform cache）：光栅化按三角形顺序消费顶点，
相邻三角形大量共享顶点，索引命中缓存时就不重跑顶点着色器。一个 $N \times N$ 的四边形网格有
$(N+1)^2$ 个顶点、$6N^2$ 个索引，无索引绘制让顶点着色器执行 $6N^2$ 次；有索引时实际次数介于
$(N+1)^2$ 与 $6N^2$ 之间，取决于顶点顺序与缓存容量。所以``索引省内存''是次要的，
``索引省顶点着色器''才是主要的。

\begin{warning}[title={索引顺序同样是布局问题}]
  索引缓冲的字节布局是平的，但它的\emph{顺序}决定变换后缓存的命中率。
  按行序而不是按三角形序排列索引，常能让同一份属性缓冲少跑三成顶点着色器。
  顶点压缩与索引重排必须一起做，见 \ref{sec:ch03-compress}
  与 \ref{ch:ch08-compute-driven} 的 \DomGPU\ 侧重排。
\end{warning}

\section{步进、实例化除数与无缓冲路径}

步进（stride）是``下一个顶点在字节上距离多远''。
\Tp{VkVertexInputBinding\Bk Description} 里的 \Tp{stride} 以字节为单位且是 32 位宽度，
而缓冲偏移与尺寸用 \Tp{VkDeviceSize}（64 位），两者不要混用。同一结构里的 \Tp{inputRate}
决定步进节奏是``每顶点''还是``每实例''，后者再配 \Tp{vertexBindingDivisor} 表示多少个实例
重取一次数据：除数为 1 就是每实例前进一行，除数为 $k$ 就是每 $k$ 个实例前进一行，除数为 0
在实例速率下等价于每个实例都读同一行，用来做广播。Vulkan 明确要求 \Tp{inputRate}
为顶点速率时除数必须为 0。\CL\ 侧对应 \Fn{glVertexAttribDivisor}，\HSL\ 侧靠
\Tp{SV\_InstanceID} 手工索引。

\begin{compare}[title={硬件步进与手工索引}]
  硬件步进由输入装配器完成，着色器不做整数运算，也不占输出分量，适合每个实例一行数据的
  规整场景。手工索引（着色器里用 \Tp{gl\_InstanceID} 或 \Tp{SV\_InstanceID} 去数组取）灵活
  得多，可以做间接索引、分组与细节层次选择，代价是把地址计算搬进着色器，并在移动
  \DomGPU\ 上多一次可能的缓存未命中。
\end{compare}

另一条路径是根本不绑属性缓冲：顶点着色器只用顶点下标自己生成数据。核心配置的 OpenGL 里
\Tp{gl\_VertexID} 是内建输入，从 \Tp{firstVertex} 起编号，一个没有启用任何数组的顶点数组
对象（VAO）配合 \Fn{glDrawArrays} 就能画出全屏三角形。

\begin{lstlisting}[style=glslstyle,caption={无缓冲顶点输入：用内建下标生成全屏三角形},label={lst:ch03-vertexid}]
#version 450
// 三个顶点即可覆盖整个裁剪空间，不需要任何顶点缓冲
void main()
{
    // gl_VertexID = 0,1,2 得到 (-1,-1) (3,-1) (-1,3)
    vec2 p = vec2(float((gl_VertexID << 1) & 2), float(gl_VertexID & 2));
    gl_Position = vec4(p * 2.0 - 1.0, 0.0, 1.0);
}
\end{lstlisting}

\GSL\ 编译进 \SPIV\ 后，\Tp{gl\_VertexID} 成为 \Tp{VertexIndex} 内建，
\Tp{gl\_InstanceID} 成为 \Tp{InstanceIndex}，语义一致。
但要注意：Vulkan 管线只要声明了任何属性，那些属性所在 binding 就必须真的绑上缓冲；
``只用下标''要求把 \Tp{vertexAttributeDescriptionCount} 置 0，这本身是合法状态。
\HSL\ 侧的对应语义是 \Tp{SV\_VertexID}；实践中多数引擎仍会绑一个几字节的占位缓冲，
因为不同代驱动对零输入布局的宽容度不一致。

\section{交错布局与分离布局}

交错（interleaved）把所有属性排成一条定长记录，分离（separate，也叫 SoA）按属性切成多段。
Vulkan 的结构体替你做了一半决定：一个 binding 只描述一段等步进的连续内存，两个数组就是
两个 binding，因此分离布局在绑定描述数组里天然表达，交错布局则是``一个 binding、多个
attribute、不同 offset''。判断标准只有一个：这一趟顶点着色器实际读了多少属性。

\begin{itemize}
  \item 全读（前向主通道）：交错更优。一条 48 字节的记录在 128 字节的访存粒度里
        能装两到三条完整顶点，属性取指合并成少数几次宽事务，带宽利用率高。
  \item 只读一部分（深度预通道只要位置）：交错浪费 $(S - P)/S$ 的带宽，
        $S$ 是记录字节数、$P$ 是位置字节数。位置占 12 字节而记录含法线、切线、UV、
        颜色、权重时 $S$ 常在 48 到 64，浪费七成。
  \item 只写一部分（把位置与法线导出给后续通道）：分离布局让写入落在各自缓冲里，
        对齐与步长更好控制，也便于压缩成 16 位宽度。
\end{itemize}

混合策略是工业界常态：位置单独一段（可被剔除通道读），其余属性交错一段（主通道读），
实例数据再单独一段（实例速率步进）。代价是 binding 与描述符数量上升，而 Vulkan 对驱动强制
的下限是 16 个 binding、16 个属性、单条步进至少 2048 字节、顶点输入分量总数至少 60。

\begin{lstlisting}[style=cstyle,caption={交错布局的顶点输入描述、缓冲绑定与上传写入},label={lst:ch03-vin-desc}]
// 一条顶点记录 48 字节：位置 12 + UV 8 + 法线 4 + 切线 4 + 颜色 4 + 权重 16
static const VkVertexInputBindingDescription kBind[2] = {
  { .binding = 0, .stride = 48,
    .inputRate = VK_VERTEX_INPUT_RATE_VERTEX },
  { .binding = 1, .stride = 16,
    .inputRate = VK_VERTEX_INPUT_RATE_INSTANCE,
    .vertexBindingDivisor = 1 },   // 每个实例前进一行
};
static const VkVertexInputAttributeDescription kAttr[5] = {
  { .location = 0, .binding = 0,
    .format = VK_FORMAT_R32G32B32_SFLOAT, .offset = 0 },
  { .location = 1, .binding = 0,
    .format = VK_FORMAT_R32G32_SFLOAT,    .offset = 12 },
  { .location = 2, .binding = 0,
    .format = VK_FORMAT_R8G8B8A8_SNORM,   .offset = 20 },
  { .location = 3, .binding = 0,
    .format = VK_FORMAT_A2B10G10R10_UNORM_PACK32, .offset = 24 },
  { .location = 4, .binding = 1,   // 实例数据在另一个 binding
    .format = VK_FORMAT_R32G32B32A32_SFLOAT, .offset = 0 },
};
const VkPipelineVertexInputStateCreateInfo vin = {
  .sType = VK_STRUCTURE_TYPE_PIPELINE_VERTEX_INPUT_STATE_CREATE_INFO,
  .vertexBindingDescriptionCount   = 2,
  .pVertexBindingDescriptions      = kBind,
  .vertexAttributeDescriptionCount = 5,
  .pVertexAttributeDescriptions    = kAttr,
};

// 布局与 CPU 端数组一一对应时，上传就是整块拷贝
Vertex *dst = (Vertex *)mappedPointer;
memcpy(dst, meshData, mesh->vertexCount * sizeof(Vertex));
vkCmdBindVertexBuffers(cmd, 0, 2, buffers, offsets);
vkCmdBindIndexBuffer(cmd, indexBuf, 0, VK_INDEX_TYPE_UINT32);
vkCmdDrawIndexed(cmd, indexCount, instanceCount, 0, 0, 0);
\end{lstlisting}

\section{三套统一内存布局规则}

同一个``块''概念，\GSL\ 给了两套规则，\HSL\ 给了一套，Vulkan 的标量布局是第四套。
混用的后果不是编译错误，而是数据错位：着色器读到隔壁那列的尾巴。

\textbf{其一，}\Tp{std140}\textbf{ 是给 uniform 缓冲设计的保守规则。}
成员偏移一律向上取整到自己的对齐值；\Tp{vec2}、\Tp{vec3}、\Tp{vec4} 都按 16 字节对齐；
标量数组的每个元素跨度也向上取整到 16 字节；矩阵按列存放、每列都当 \Tp{vec4} 处理，于是
\Tp{mat3}\ 占 48 而不是 36 字节；结构体对齐至少 16，大小补齐到 16 的倍数。

\textbf{其二，}\Tp{std430}\textbf{ 是存储缓冲的紧凑规则。}
标量按自身宽度对齐，\Tp{vec2}\ 数组元素跨度是 8，标量数组元素跨度是 4，\Tp{mat3}\ 变成
36 字节；但 \Tp{vec3}\ 仍按 16 字节对齐，这一条不变。结构体对齐取最大成员对齐值，
不再强制凑到 16。

\textbf{其三，}\HSL\ \Tp{cbuffer}\textbf{ 是十六字节一行的打包规则。}
向量成员必须从行首开始，任何成员不得跨行；标量可以填进向量剩下的那几列，所以
\Tp{float3}\ 紧跟一个 \Tp{float}\ 恰好合成 16 字节；标量数组按 4 字节逐个填满一行，
这是它与 \Tp{std140}\ 最大的差别；块大小补齐到 16。

\Tp{vec3}\textbf{ 是跨距陷阱的主角。}三套规则里它都占 16 字节的\emph{槽位}，但只有
\Tp{cbuffer}\ 允许用标量填掉第 16 个字节。如果 \DomCPU\ 侧照抄一个 12 字节的三维向量
（\CPP\ 生态里最常见的 \Tp{glm::vec3}\ 正是 12 字节），从第 2 个元素起偏移全部前移 4 字节。

\begin{longtable}{>{\raggedright\arraybackslash}p{0.19\textwidth}
                  >{\raggedright\arraybackslash}p{0.21\textwidth}
                  >{\raggedright\arraybackslash}p{0.21\textwidth}
                  >{\raggedright\arraybackslash}p{0.21\textwidth}}
\caption{三套统一内存布局规则对照：数组、向量、结构体与放松手段}\label{tab:ch03-layout-rules} \\
\toprule
项目 & \GSL\ \Tp{std140} & \GSL\ \Tp{std430} & \HSL\ \Tp{cbuffer} \\
\midrule
\endfirsthead
\toprule
项目 & \GSL\ \Tp{std140} & \GSL\ \Tp{std430} & \HSL\ \Tp{cbuffer} \\
\midrule
\endhead
\bottomrule
\endfoot
典型用途 & uniform 缓冲的 \Tp{uniform} 块 & 存储缓冲的 \Tp{buffer} 块 & \Tp{cbuffer}\ 常量寄存器组 \\
标量数组元素跨度 & 16 字节，向上取整到 \Tp{vec4} & 4 字节，按标量自身宽度 & 4 字节，四个填满一行 \\
\Tp{vec2}\ 对齐 & 16 字节 & 8 字节 & 16 字节，必须行首 \\
\Tp{vec3}\ 对齐与占用 & 16 字节对齐，占 16 & 16 字节对齐，占 16 & 16 字节行首，尾部可接标量 \\
\Tp{mat3}\ 占用（列优先） & 48 字节，每列补成 \Tp{vec4} & 36 字节，每列就是 \Tp{vec3} & 48 字节，三行各 16 \\
结构体对齐 & 至少 16，且不小于最大成员 & 等于最大成员对齐 & 从行首开始，成员不得跨行 \\
结构体数组跨度 & 补齐到 16 的倍数 & 补齐到结构体自身对齐 & 补齐到 16 的倍数 \\
块尾补齐 & 16 的倍数 & 最大成员对齐的倍数 & 16 的倍数 \\
能否显式放松 & \Tp{packed}、\Tp{shared}\ 只在 OpenGL 侧存在，Vulkan \GSL\ 不接受，驱动支持也不一致 &
  本就紧凑，无需放松 & \Tp{packoffset}\ 逐成员指定 \\
第四套（标量布局） & \Tp{VK\_EXT\_\Bk scalar\_block\_\Bk layout} 让成员按自身宽度对齐，需能力位 &
  同左 & 概念上等价，本来就能填尾 \\
\end{longtable}

\begin{lstlisting}[style=glslstyle,caption={顶点属性声明与一个 std140 uniform 块：注释里的偏移量就是内存镜像},label={lst:ch03-glsl-layout}]
#version 450
layout(location = 0) in vec3  inPos;      // 归一化整数由 API 转换
layout(location = 1) in vec2  inUV;
layout(location = 2) in vec4  inNrmSign;  // xyz 法线 SNORM，w 存切线符号
layout(location = 3) in vec4  inColor;    // 10/10/10/2 归一化

layout(binding = 0) uniform SceneParams { // 默认按 std140 解释
    mat4  viewProj;       // 偏移   0，64 字节（4 列 x vec4）
    vec3  cameraPos;      // 偏移  64，槽位 16 字节
    float time;           // 偏移  76，与 cameraPos 同槽
    vec3  lightDirs[4];   // 偏移  80，每个元素 16 字节
    float lodBias;        // 偏移 144
} scene;                  // 块大小补齐到 160 字节

layout(std430, binding = 1) readonly buffer SkinMatrices {
    mat4 boneMatrices[];  // 每列 16 字节，每个 mat4 占 64 字节
};

void main()
{
    gl_Position = scene.viewProj * vec4(inPos, 1.0);
}
\end{lstlisting}

上面这块 \Tp{SceneParams}\ 在三套规则下的尾部大小分别是 160、148 与 160 字节，差异全部来自
数组跨度与尾部补齐。不要靠记忆，要按规则现场算：偏移永远等于``当前偏移向上取整到成员对齐''。

\section{\DomCPU\ 结构体到着色器结构体的镜像}

镜像的目标不是省字节，而是让 \CPP\ 侧的 \Tp{sizeof}\ 与着色器侧的块大小一致。三条纪律：
不要用 12 字节的三维向量，不要用位域，不要用 \Tp{\#pragma\ pack}\ 压小结构体；正确做法是把
填充写进结构体本身，并用 \Tp{alignas}\ 钉住槽位。

\begin{lstlisting}[style=cppstyle,caption={与 std140 块逐字节对齐的 C++ 镜像，含显式填充与静态校验},label={lst:ch03-cpp-mirror}]
struct alignas(16) Float3Slot {   // 一个 vec3：16 字节槽位
    float x, y, z, pad;           // 填充显式写出，不留给编译器猜
};
struct alignas(16) Matrix4 {      // 列优先，与 GLSL mat4 一致
    float m[4][4];                // 4 列 x 16 字节 = 64 字节
};

struct SceneParams {              // 镜像上面那个 std140 uniform 块
    Matrix4    viewProj;          //   0
    Float3Slot cameraPos;         //  64
    float      time;              //  76  与上一槽的尾部共用
    Float3Slot lightDirs[4];      //  80  数组元素跨度 16
    float      lodBias;           // 144
    float      tailPad[3];        // 148  把块补到 160
};
static_assert(sizeof(SceneParams) == 160, "std140 镜像必须逐字节一致");
static_assert(sizeof(Float3Slot) == 16, "vec3 槽位不能是 12");

constexpr uint32_t RoundUp(uint32_t value, uint32_t align) {
    return (value + align - 1u) / align * align;
}
\end{lstlisting}

\begin{note}[title={三种镜像策略的取舍}]
  第一种是显式填充，即上例：可读、可校验，代价是浪费的字节要自己维护。
  第二种是整体升宽：把所有向量都写成 \Tp{vec4}\ 与 \Tp{float4}，块里只出现 4 的倍数，
  三套规则自然合一，这是跨 API 引擎最常见的选择。第三种是反向生成：用脚本从着色器块生成
  \CPP\ 头，保证两边不漂移，代价是多一条代码生成链路，见 \ref{ch:ch07-spirv-toolchain}。
  \Tp{std430}\ 存储缓冲可以用紧凑写法，但同一结构一旦挪进 uniform 缓冲就会错位。
\end{note}

\HSL\ 侧的镜像相对自由，因为打包规则允许标量填进向量尾部；但 \Tp{cbuffer}\ 与
\Tp{std140}\ 在数组上的分歧是实打实的：\Tp{float weights[8]}\ 在 \Tp{cbuffer}\ 里是
32 字节，在 \Tp{std140}\ 里是 128 字节。同一份数据要喂两个前端时，优先把标量数组换掉。

\section{顶点压缩与蒙皮的矩阵预算}
\label{sec:ch03-compress}

顶点数据的字节数直接乘在带宽上，压缩优先级永远高于着色器优化。
可用手段基本都落在``归一化整数''这一类：把浮点写成分数，让 API 在取指时还原。

\begin{itemize}
  \item 位置：\Tp{R16G16B16\_SFLOAT}，或 \Tp{SNORM}\ 三通道加缩放基。
        半精度在 $[-1,1]$ 网格里约三位十进制有效数字，配合物件级缩放与偏移，
        重建误差通常小于半个像素。
  \item 法线与切线：\Tp{R8G8B8A8\_SNORM}\ 的分量步长是 $1/127$，约 $0.008$，
        还原后角度误差量级在 $0.5^\circ$；切线第 4 分量常只存符号 $\pm 1$，8 位足够。
  \item 颜色：\Tp{A2B10G10R10\_UNORM}\ 用 4 字节存三通道 10 位色，
        是顶点输入里被强制支持的打包格式。
  \item UV：\Tp{R16G16\_UNORM}\ 在 4K 贴图上步长约 $1/65535$，比全浮点省一半，
        但重复平铺的坐标不能用，因为超出 $[0,1]$ 会被截断。
\end{itemize}

\begin{longtable}{>{\raggedright\arraybackslash}p{0.15\textwidth}
                  >{\raggedright\arraybackslash\small}p{0.26\textwidth}
                  >{\raggedright\arraybackslash}p{0.12\textwidth}
                  >{\raggedright\arraybackslash}p{0.28\textwidth}}
\caption{常见顶点属性压缩选择与代价}\label{tab:ch03-compress} \\
\toprule
语义 & 格式 & 字节 & 主要风险 \\
\midrule
\endfirsthead
\toprule
语义 & 格式 & 字节 & 主要风险 \\
\midrule
\endhead
\bottomrule
\endfoot
位置 & \Tp{R16G16B16\_SFLOAT} & 12 & 超出半精度范围变 \Tp{inf}；需要物件级缩放基 \\
位置 & 三通道 \Tp{SNORM}\ 加缩放 & 12 以内 & 共享顶点跨物件时会露出接缝 \\
法线 & \Tp{R8G8B8A8\_SNORM} & 4 & 还原后必须再归一化，否则插值长度不足 1 \\
切线 & \Tp{R8G8B8A8\_SNORM}\ 存符号 & 4 & $w$ 只能表达 $\pm 1$，四元数形式要另存 \\
UV & \Tp{R16G16\_UNORM} & 4 & 平铺坐标被截断到 $[0,1]$ 之内 \\
颜色 & \Tp{A2B10G10R10\_UNORM} & 4 & 第 4 分量只有 2 位精度 \\
蒙皮索引 & \Tp{R8G8B8A8\_UINT} & 4 & 无归一化，直接当整数用 \\
蒙皮权重 & \Tp{R8G8B8A8\_UNORM} & 4 & 四分量之和未必恰好为 1，需在着色器里补 \\
\end{longtable}

归一化整数的共性坑是转换发生在 \DomGPU\ 取指阶段。\Tp{vec3}\ 属性配三通道 \Tp{SNORM}\
读到浮点，\Tp{ivec3}\ 配 \Tp{SINT}\ 读到整数；声明类型的符号性与格式不一致属于未定义行为，
调试时表现为随机噪声，而编译器不会报警。

蒙皮是布局问题而不是动画问题，因为预算全在字节上。影响矩阵可以是 \Tp{mat4}（64 字节）、
去掉最后一行的 12 个浮点、或一双四元数（32 字节）。骨骼数与每矩阵字节的乘积决定它放得进
uniform 缓冲还是必须走存储缓冲：Vulkan 对驱动强制的 uniform 缓冲可绑定范围下限是 16384
字节，\Tp{std140}\ 的 \Tp{mat4}\ 数组每个元素跨度 64 字节，于是一个块最多装 256 根骨骼。
想省到 48 字节并不简单：\Tp{mat4x3}\ 的列是 \Tp{vec3}，即使在 \Tp{std430}\ 下每列仍被补到
16 字节，四列还是 64；真要 48 字节得存成三列 \Tp{vec4}（\Tp{mat3x4}）再手工拼装。

\begin{lstlisting}[style=glslstyle,caption={从 std430 存储缓冲取四路影响并加权混合顶点},label={lst:ch03-skin-glsl}]
#version 450
layout(std430, binding = 1) readonly buffer SkinMatrices {
    mat4 bone[128];        // std430 下每个 mat4 仍占 64 字节
};
layout(location = 4) in vec4  inBoneWeights;
layout(location = 5) in ivec4 inBoneIndices;

vec3 Skinned(vec3 p)
{
    vec3  acc = vec3(0.0);
    float sum = 0.0;
    for (int i = 0; i < 4; ++i) {
        mat4 m = bone[inBoneIndices[i]];
        acc += (m * vec4(p, 1.0)).xyz * inBoneWeights[i];
        sum  += inBoneWeights[i];
    }
    return acc / sum;      // 归一化整数权重必须补这一刀
}
\end{lstlisting}

\begin{warning}[title={蒙皮布局的两个隐蔽错误}]
  第一，权重读成 \Tp{vec4}\ 而格式是无符号归一化时，四个分量落在 $[0,1]$，原始和常常是
  $255/255$ 的近似而不是 1；忘记除以权重和会让网格整体缩小。第二，把骨骼数组声明成运行时
  长度（\Tp{mat4 bone[]}）就把它变成了存储缓冲；uniform 块里的数组长度必须是编译期常量，
  所以可变骨骼数只能走存储缓冲。
\end{warning}

\begin{guideline}[title={本章要点}]
  \DomGPU\ 只认字节：先决定哪一趟通道读哪些属性，再决定交错还是分离。
  步进与实例化除数是输入装配器的能力，能用就别把地址计算搬进着色器；
  无缓冲路径是合法选项，但要把属性数量真的清零。
  \Tp{std140}、\Tp{std430}、\Tp{cbuffer}\ 是三套不同算法，
  共同的不变量只有``偏移向上取整到成员对齐''与``\Tp{vec3}\ 占 16 字节槽位''两条；
  分歧集中在标量数组跨度、\Tp{mat3}\ 大小与结构体尾部补齐。
  \DomCPU\ 镜像要把填充写进结构体并用 \Tp{static\_assert}\ 钉死 \Tp{sizeof}，
  宁可浪费字节，也不要让两边靠默契对齐。
  压缩优先级高于一切着色器优化：位置、法线、UV 各降一半宽度，带宽就少一半。
\end{guideline}

% ============================================================
% ---- frag_ch04.tex ----
\chapter{空间、变换与投影}
\label{ch:ch04-transform-space}

从三角形到屏幕像素之间只有一件事在发生：一串矩阵乘法，加一次除法，再加一次仿射映射。
但每个 API 在这三步的细节上都不肯完全对齐，于是移植的成本几乎全部落在约定上：
谁在哪个位置写 $\pm 1$、深度落在 $[0,1]$ 还是 $[-1,1]$、$y$ 轴朝哪边、纹素中心在哪里。
本章把这些约定算清楚。属性与内存布局见 \ref{ch:ch03-vertex-input}，
语言侧的向量与矩阵类型见 \ref{ch:ch05-glsl} 与 \ref{ch:ch06-hlsl}。

\section{五级空间与一次绘制的旅程}

工程上惯用的划分是模型、世界、视图、裁剪、屏幕五级，其中裁剪空间与屏幕空间之间还夹着
一个归一化设备坐标（NDC）空间。各自的性质差别很大：

\begin{itemize}
  \item 模型空间：资产自己的坐标系，单位与朝向由导出工具决定，通常与 \DomGPU\ 无关。
  \item 世界空间：场景里唯一的公共坐标系，所有物件的层次变换最终合成到这里。
  \item 视图空间：相机在原点、看向某条轴。法线、光照、雾与深度精度分析在这一级最清楚。
  \item 裁剪空间：\Tp{gl\_Position} 或 \Tp{SV\_Position} 所在的空间，是四维齐次空间。
        它的 $x$、$y$、$z$ 还没除以 $w$，范围是 $[-w, w]$（或 $[0, w]$），
        所以裁剪空间里的``单位立方体''随 $w$ 缩放，不是一个立方体。
  \item 屏幕空间：视口映射之后的坐标，$z$ 变成窗口深度，参与深度比较。
\end{itemize}

关键点在于这五级里只有前四级是\emph{矩阵}，最后一步是固定功能的仿射映射加一次透视除法。
除法不可关闭，但它的影响必须自己承担：顶点着色器输出的任何量，硬件都按 $1/w$ 加权插值，
因此在屏幕上呈线性的东西必须在顶点侧先除以 $w$，再交给插值器。

\begin{note}[title={谁在做除法}]
  硬件对插值按 $1/w$ 加权，正好等价于三维空间里的线性插值。于是视图空间位置、
  世界空间位置这类\emph{三维}量直接当输出变量传给片段着色器就是正确的；而屏幕坐标、
  纹素步进量这类\emph{屏幕线性}量必须显式除以 $w$ 或改用 \Tp{noperspective}。
  搞反一次的后果是贴图在斜面上被压缩或拉伸，调半天才发现是插值模式而不是 UV 错了。
\end{note}

\section{矩阵约定：行列优先、向量侧别与手性}

同一份变换有四个独立的选择，任何一个都可以反过来，组合起来才形成 API 的方言。

\textbf{行优先还是列优先}是\emph{内存}问题：\GSL\ 的 \Tp{mat4}\ 按列连续存放，
每列是一个 \Tp{vec4}；\HSL\ 的 \Tp{float4x4}\ 默认按行连续，可以用 \Tp{pack\_\Bk matrix}
显式改。两边内存不同，所以 \DomCPU\ 侧把同一个数组喂给两个前端时，必须知道自己在写
哪一侧，见 \ref{ch:ch03-vertex-input} 的镜像写法。

\textbf{行向量还是列向量}是\emph{代数}问题：列向量写 $v' = Mv$，行向量写 $v' = vM$。
两者互为转置，与行列优先无关。DirectXMath 习惯行向量后乘，\GSL\ 习惯列向量前乘，
于是同一份乘法顺序在两边是反的。

\textbf{左手还是右手}是\emph{几何}问题：右手系相机看向 $-z$，左手系看向 $+z$。
差别不止一个符号：手性改变叉乘方向、三角形绕序的可见面判定，以及 $z$ 到深度的映射符号。
左手系的资产配右手相机需要一次 $z$ 取反与绕序翻转；把左手系的透视矩阵直接搬到右手系，
最典型的表现是整个场景落在近平面之后，画面全空。

\begin{lstlisting}[style=hlslstyle,caption={HLSL 的行优先与行向量写法：与 GLSL 的列优先前乘互为转置},label={lst:ch04-hlsl-rowmajor}]
#pragma pack_matrix(row_major)

struct VSIn  { float3 pos : POSITION;   float2 uv : TEXCOORD0; };
struct VSOut { float4 pos : SV_Position; float2 uv : TEXCOORD0; };

cbuffer Camera : register(b0)
{
    float4x4 viewProj;    // 行优先，与 GLSL 的 mat4 内存布局互为转置
};

VSOut VSMain(VSIn vin)
{
    VSOut o;
    float4 p = float4(vin.pos, 1.0);
    o.pos = mul(p, viewProj);   // 行向量后乘；GLSL 侧是 viewProj * p
    o.uv  = vin.uv;
    return o;
}
\end{lstlisting}

\section{四个 API 的 NDC 差异与翻面的三种做法}

NDC 的定义在所有 API 里都是``裁剪坐标除以 $w$''，分歧只在除法之后的盒子形状与屏幕映射。
把差异一次列清：

\begin{longtable}{>{\raggedright\arraybackslash}p{0.12\textwidth}
                  >{\raggedright\arraybackslash}p{0.15\textwidth}
                  >{\raggedright\arraybackslash}p{0.13\textwidth}
                  >{\raggedright\arraybackslash\small}p{0.14\textwidth}
                  >{\raggedright\arraybackslash\small}p{0.26\textwidth}}
\caption{API 与 NDC 深度、$y$ 方向、纹素原点与常见坑}\label{tab:ch04-ndc} \\
\toprule
API & NDC 深度 & $y$ 方向 & 纹素原点 & 最常见的坑 \\
\midrule
\endfirsthead
\toprule
API & NDC 深度 & $y$ 方向 & 纹素原点 & 最常见的坑 \\
\midrule
\endhead
\bottomrule
\endfoot
OpenGL 3.3 到 4.6 & $x,y,z$ 都落在 $[-1,1]$ & NDC 与窗口都 $y$ 向上 &
  纹理左下，\Tp{gl\_\Bk FragCoord} 同在左下 &
  把纹素当角点采样产生的半纹素偏移；深度要换成 $[0,1]$ 得靠 \Fn{glClipControl}（4.5 进核心） \\
Vulkan & $x,y$ 为 $[-1,1]$，$z$ 为 $[0,1]$ & NDC 正 $y$ 朝下（帧缓冲原点左上） &
  图像原点左上，\Tp{gl\_\Bk FragCoord} 也在左上 &
  从 \GSL\ 或 \HSL\ 搬来的矩阵必须翻 $y$，或用 \Tp{VK\_\Bk KHR\_\Bk maintenance1}（1.1 进核心）允许的负视口高度 \\
Direct3D 11/12 & $x,y$ 为 $[-1,1]$，$z$ 为 $[0,1]$ & NDC $y$ 向上，屏幕 $y$ 向下 &
  纹理左上，中心在 $i + 0.5$ &
  与 D3D9 的两处差异（$z$ 曾是 $[-1,1]$、采样点在角上）是遗留资产上下翻面的根因 \\
Metal & $x,y$ 为 $[-1,1]$，$z$ 为 $[0,1]$ & NDC $y$ 向上，帧缓冲原点左上 &
  纹理左上 &
  从 \GSL\ 转写片段程序时行号方向最容易错，翻转要连着 $0.5$ 一起算 \\
\end{longtable}

翻 $y$ 有三条路，各自有代价。第一条是在投影矩阵的 $y$ 行上取负，最省事，但同一个矩阵
就不能再用于立方图捕获这类需要朝向一致的场合。第二条是负视口高度（Vulkan 允许，语义上
等价于把 NDC 到窗口的映射翻过来），矩阵保持干净，但任何读窗口坐标的代码都要跟着翻。
第三条是在顶点着色器里对 \Tp{gl\_Position} 的 $y$ 分量取负，作用范围最小，
代价是多一处只有着色器知道的隐式约定。

\begin{warning}[title={深度范围不是只有 0 与 1 的差别}]
  $z$ 的映射符号与深度比较的默认方向是绑在一起的：多数 API 默认 \Tp{LESS}，即更小者胜。
  远平面映射到 $1$ 时清屏值取 $1$；换成 $[-1,1]$ 时清屏值仍是 $1$；一旦改用 reverse-Z，
  远平面变成 $0$，于是\emph{清屏值与比较函数必须同时改}，漏改一侧的表现是整屏被最远的
  物体覆盖。想把深度边界推到 $[0,1]$ 之外，Vulkan 需要
  \Tp{VK\_\Bk EXT\_\Bk depth\_\Bk range\_\Bk unrestricted}；想让物体跨过近平面而不被裁掉，
  需要显式关闭深度裁剪（Direct3D 与 Metal 有对应开关，Vulkan 侧靠扩展），
  见 \ref{ch:ch09-vulkan}。
\end{warning}

\section{透视投影的推导}

推导只需要一个相似三角形：相机在原点看向 $-z$，像平面在 $z = -1$，
距离 $d$ 处的视口高度正好是 $2h$，其中 $h = 1/\tan(\theta/2)$、$\theta$ 是垂直视场角。
于是 $y_{ndc} = h \cdot y_{view} / d$，$x_{ndc} = h \cdot x_{view} / (d \cdot a)$，
$a$ 是宽高比。因为 $d = -z_{view}$，``除以距离''这件事必须由 $w = -z_{view}$ 提供，
这就是透视矩阵里那个 $\pm 1$ 的全部来历。深度这一维没有几何约束，只有区间映射：
设 $z_{clip} = A \cdot z_{view} + B$，则 $z_{ndc} = -A + B/d$，
把 $d = n$ 映到 $z_n$、$d = f$ 映到 $z_f$，解二元一次方程得
$B = (z_n - z_f) \cdot nf/(f - n)$ 与 $A = (z_n - z_f) \cdot f/(f - n) - z_n$。

\begin{lstlisting}[style=cppstyle,caption={透视与正交矩阵生成：一次覆盖深度范围、Y 翻转、reverse-Z 与无穷远平面四个约定},label={lst:ch04-matrix-gen}]
// 视图空间约定：右手系，相机看向 -Z，前方距离 d 满足 z_view = -d
// 存储为列优先 c[col][row]，与 GLSL 的 mat4 内存布局一致
enum class DepthOut { kZeroToOne, kMinusOneToOne, kReversed };

struct Mat4 { float c[4][4]; };

static void Endpoints(DepthOut out, float &zn, float &zf) {
    zn = (out == DepthOut::kMinusOneToOne) ? -1.0f
       : (out == DepthOut::kReversed)       ?  1.0f : 0.0f;
    zf = (out == DepthOut::kReversed) ? 0.0f : 1.0f;
}

Mat4 Perspective(float fovY, float aspect, float n, float f,
                 DepthOut out, bool flipY, bool infiniteFar) {
    float zn, zf; Endpoints(out, zn, zf);
    const float h = 1.0f / std::tan(0.5f * fovY);
    const float s = zn - zf;
    float a, b;                       // clip.z = a * z_view + b
    if (infiniteFar) {                // f 趋于无穷时的极限形式
        a = -zf;      b = s * n;
    } else {
        a = s * f / (f - n) - zn;   b = s * n * f / (f - n);
    }
    Mat4 m{};
    m.c[0][0] = h / aspect;
    m.c[1][1] = flipY ? -h : h;       // 目标帧缓冲原点在左上时取负
    m.c[2][2] = a;
    m.c[2][3] = -1.0f;                // clip.w = -z_view：除法与插值的根
    m.c[3][2] = b;
    return m;
}

Mat4 Orthographic(float l, float r, float b0, float t0,
                  float n, float f, DepthOut out, bool flipY) {
    float zn, zf; Endpoints(out, zn, zf);
    const float s = zn - zf;
    Mat4 m{};
    m.c[0][0] = 2.0f / (r - l);
    m.c[1][1] = (flipY ? -2.0f : 2.0f) / (t0 - b0);
    m.c[2][2] = s / (f - n);
    m.c[3][0] = -(l + r) / (r - l);
    m.c[3][1] = -(b0 + t0) / (t0 - b0);
    m.c[3][2] = zn + s * n / (f - n);
    m.c[3][3] = 1.0f;                 // 正交没有透视除法
    return m;
}
\end{lstlisting}

用这个函数生成四种方言只需要两个参数的组合：\GSL\ 面向 OpenGL 时传
\Tp{kMinusOneToOne} 且不翻 $y$；Vulkan 传 \Tp{kZeroToOne} 并翻 $y$；Direct3D 与 Metal
传 \Tp{kZeroToOne} 不翻 $y$（前提是相机与矩阵都按同一手性处理）。
正交投影的 $n$ 可以取 $0$：它的深度是距离的线性函数，$n$ 只出现在减法里。透视不行，
因为 $z_{ndc} = -A + B/d$ 的全部变化量都由 $B$ 承担，而 $B$ 含因子 $n$：把 $n$ 推向 $0$
等于把所有深度压向同一个常数，编码就此浪费，这正是近平面不能推到无穷远的理由。
反过来，把远平面推到无穷是良构的，极限就是代码里 \Tp{infiniteFar} 那一支，
此时 $A$ 只剩 $-z_f$、$B$ 只剩 $(z_n - z_f) n$。

\section{深度精度、reverse-Z 与无穷远平面}

把 $[0,1]$ 前向映射写成距离的函数：$z(d) = (f/(f - n))(1 - n/d)$，
对 $d$ 求导正比于 $1/d^{2}$。定点深度把 $[0,1]$ 等分成 $2^{k}$ 层，
于是分辨率在近处极密、远处极疏。取 $n = 0.1$、$f = 1000$，$z = 0.5$ 出现在
$d \approx 2n = 0.2$ 处：一半编码只服务于相机前 20 厘米，剩下的都要挤进后 999 米。
这是\emph{分布}问题，与位宽无关：24 位与 32 位定点只是把同一条曲线切得更细而已。

两条出路。第一条是把 $[n, f]$ 收窄到当前帧真正需要的范围：级联阴影各自用自己的区间，
后处理只覆盖可见几何，这是最有效的做法。第二条是换用浮点深度并配合 reverse-Z：
取 \Tp{kReversed} 加 \Tp{infiniteFar}，映射退化成 $z(d) = n/d$，
近处落在 $1$ 附近、远处趋向 $0$；浮点格式在 $0$ 附近的编码密度远高于 $1$ 附近，
于是精细的比较能力落在几何最密集的地方，远处的共面闪烁随之减少。

\begin{bestpractice}[title={reverse-Z 的三件配套工作}]
  一、清屏深度与比较函数一起改：远平面是 $0$，\Tp{LESS} 换成 \Tp{GREATER}，
  两者只改一侧会直接让早深度测试失效或整屏翻黑。
  二、把远平面从矩阵里彻底拿掉，$f$ 不再出现在常量里，剔除、雾与淡出都要改按距离算，
  而不是按深度值算。
  三、任何显式写深度的着色器路径（\Tp{gl\_FragDepth}\ 一类）都要跟着反转，
  否则两个来源的深度会打架，见 \ref{ch:ch13-raytracing}。
\end{bestpractice}

\section{视锥剔除与纹素对齐}
\label{sec:ch04-cull}

视锥是投影矩阵定义的六个半空间的交。不必构造平面顶点，直接从 \Tp{viewProj} 的行做加减：
记第 $i$ 行为 $r_i$（列优先存储下它的分量是 \Tp{c[0][i]} 到 \Tp{c[3][i]}），则
左面 $= r_3 + r_0$、右面 $= r_3 - r_0$、下面 $= r_3 + r_1$、上面 $= r_3 - r_1$、
远面 $= r_3 - r_2$；近平面依赖深度约定，$[-1,1]$ 时是 $r_3 + r_2$，
而 $[0,1]$ 的条件是 $z \ge 0$，于是近平面就是 $r_2$ 本身。
每个面除以法线长度后，球体测试是 $\Tp{dot}(N, C) + d \ge -r$：每个面三次乘法、
三次加法，六个面合计约 $36$ 次浮点操作（逐面测试时大多数球在一两个面上就判负，可提前退出）。

预算要算在总开销上：一万个包围球是 $3.6 \times 10^{5}$ 次浮点操作，在任何现代 \DomCPU\
上都远低于 0.1 毫秒，真正的成本是\emph{被剔除}的物件省下的绘制调用与顶点处理。
所以实践顺序是先做层次剔除（节点级包围体一次排除整棵子树），再做叶级精确测试，
最后才考虑把剔除搬到 \DomGPU\ 上做压缩与间接绘制，见 \ref{ch:ch08-compute-driven}。

\begin{lstlisting}[style=glslstyle,caption={顶点变换：gl\_Position 与插值模式的显式选择},label={lst:ch04-glsl-xform}]
#version 450
layout(location = 0) in vec3 inPos;
layout(location = 1) in vec2 inUV;

layout(binding = 0) uniform Camera { mat4 viewProj; mat4 view; } cam;

out vec2 vUV;                     // 默认按 1/w 加权插值，透视正确
out vec3 vViewPos;                // 三维线性量，交给硬件即可
noperspective out vec2 vScreen;   // 屏幕线性量：自己除以 w
flat out int vMaterialId;         // 整数不能插值

void main()
{
    vec4 viewPos = cam.view * vec4(inPos, 1.0);
    vec4 clip    = cam.viewProj * viewPos;

    gl_Position = clip;             // 裁剪空间；透视除法由固定功能完成
    vUV         = inUV;
    vViewPos    = viewPos.xyz;
    vScreen     = (clip.xy / clip.w) * 0.5 + 0.5;
    vMaterialId = 7;
}
\end{lstlisting}

纹素对齐（texel perfect）要求两件事：一个屏幕像素恰好落一个纹素，且采样点落在纹素中心。
第一条由尺寸决定：窗口 $W \times H$、纹理 $T_w \times T_h$，比例 $T_w/W = 1$；
第二条由偏移决定，把窗口坐标换成 UV 时要减掉那个半像素：$u = (x - 0.5)/W$。
翻 $y$ 时连续坐标是 $y' = H - y$，而纹素\emph{索引}是 $i' = N - 1 - i$，
两者差一个 $0.5$，混用就会出现那条经典的``偏移半格''。

\begin{itemize}
  \item 视口与纹理尺寸不等时，放大用双线性即可，缩小需要 mipmap；mip 级别由 UV 的
        屏幕导数决定，所以在 $2 \times 2$ 像素块里用条件分支算出的 UV 会让导数失真，
        要改用显式梯度的采样接口。
  \item 全屏通道里位置由下标生成、UV 由 $0.5$ 偏移决定，比绑顶点缓冲更省，
        见 \ref{ch:ch03-vertex-input} 的无缓冲路径。
  \item 多重采样下``中心''换成样本位置；需要保证插值点落在三角形内部时用 \Tp{centroid}，
        它可以和 \Tp{noperspective} 叠加。
\end{itemize}

\begin{guideline}[title={本章要点}]
  五级空间里只有前三级是纯几何，第四级靠矩阵、第五级靠固定功能的除法与视口映射；
  插值模式属于第五级，必须和矩阵约定一起决定。
  行优先与行向量、左手与右手是彼此独立的选择，任何一处反了都会转置或镜像整个场景。
  深度范围与 $y$ 方向是每个 API 的两个方言参数：把翻转放进矩阵、放进视口还是放进着色器，
  取决于还有谁在读窗口坐标。
  透视矩阵里那个 $\pm 1$ 不可省略；$n$ 不能推到 $0$，$f$ 可以推到无穷，
  reverse-Z 的价值来自浮点编码密度的分布，而不是视觉偏好。
  剔除的数学便宜、结果值钱；纹素对齐的坑几乎总是那半个像素。
\end{guideline}

% ============================================================
% ---- frag_ch05.tex ----
\part{着色语言：GLSL、HLSL 与 SPIR-V}
\chapter{GLSL：语法、限定符与内建}
\label{ch:ch05-glsl}

\DomCPU{} 侧的语言把内存、类型与生命周期都交给程序员，\DomGPU{} 侧的着色语言则把
这些交还给硬件与驱动。本章讲 \GSL{}：它的版本指令如何在编译期切分语法，它的类型系统
在哪些地方与 \CPP{} 对不上，限定符与存储类各自的更新代价，以及控制流写法怎样被硬件的
并行粒度反过来约束。第~\ref{ch:ch06-hlsl} 章给出 \HSL{} 的镜像视角，
第~\ref{ch:ch07-spirv-toolchain} 章讲两者共同的后端 \SPIV{}。

\section{版本与 profile：首行版本指令决定可用语法}

\GSL{} 源码的第一条有效指令通常是版本声明，形式是 \Tp{\#version} 加一个十进制号：
\Tp{110}、\Tp{120} 对应 \Lib{OpenGL} 2.x 那一代；\Tp{130}、\Tp{140}、\Tp{150}
对应 3.0/3.1/3.2；\Tp{330} 对应 3.3；\Tp{400} 到 \Tp{460} 对应 4.x 各台阶。号后可以
跟 profile 关键字：\Tp{core} 或 \Tp{compatibility}，都不写时默认落在
\Tp{compatibility}。这条指令之所以关键，是因为它决定的是前端解析行为，而不是运行时
能力探测：写 \Tp{in}/{\Bk}\Tp{out} 却声明 \Tp{120}，编译直接失败；反过来
\Tp{compatibility} 里保留的 \Tp{attribute}、\Tp{varying}、\Tp{gl\_FrontColor}
这类固定功能遗产在 \Tp{core} 下就是未定义标识符。跨版本移植没有隐式降级，只有显式
变体：要么用预处理（\GSL{} 有 \Tp{\#define}、\Tp{\#if}、\Tp{\#ifdef}）裁出分支，
要么在构建时准备多份源文件，见第~\ref{ch:ch07-spirv-toolchain} 章。

移动端是另一条平行谱系：\Tp{\#version 100 es}（\Lang{OpenGL ES} 2.0）、
\Tp{300 es}、\Tp{310 es}、\Tp{320 es} 分别对应 ES 2.0/3.0/3.1/3.2。ES 分支强制
精度声明：全局或每个变量前必须给出 \Tp{precision} 语句（\Tp{highp}、\Tp{mediump}、
\Tp{lowp}），否则片元着色器连编译都过不去。ES 1.00 还带着附录 A 的一整套限制
（循环边界必须是编译期常量、不许递归），到 \Tp{300 es} 才松开。

面向 \Lib{Vulkan} 的 \GSL{} 又是一层方言。工具链普遍接受 \Tp{\#version 450 core}，
在此之上提供 \Tp{layout(push\_constant)}、\Tp{layout(constant\_id = n)} 这类
\DomGPU{} API 专属声明；它们不属于 \Lib{OpenGL} 4.5 的 \GSL{} 4.50，把这类源码喂给
\DomCPU{} 侧的 \Lib{OpenGL} 上下文只会拿到语法错误。反过来，\Tp{layout(location =)}
之外的那些 OpenGL 专属状态查询也不会在 Vulkan 里出现。

\begin{note}[title={版本声明是编译期契约，不是运行期开关}]
同一份 \GSL{} 源码在两个版本下的语法树可以完全不同：\Tp{150} 之前没有
\Tp{flat} 插值限定，\Tp{330} 之前没有隐式 LOD 的分离采样器接口，\Tp{430} 之前
既没有 \Tp{buffer} 块也没有 compute 阶段。写跨引擎的着色器时，先确定最低
\Tp{\#version}，再决定能用哪些限定符，顺序不能反。
\end{note}

\section{类型系统：与 \DomCPU{} 侧 \CPP{} 对不上的地方}

向量与矩阵是首要差异源。\DomCPU{} 侧一个 \Tp{float3} 结构体可以按 12 字节排布并
随意取址；\GSL{} 里 \Tp{vec3} 是独立类型，没有指针也没有引用，参数传递只有
\Tp{in}、\Tp{out}、\Tp{inout} 三种显式形式。真正影响字节流的是跨距：\Tp{std140}
布局下 \Tp{vec3} 必须对齐到 16 字节，\Tp{mat3} 被当成三列、每列补到 \Tp{vec4} 的
宽度，因此 \cpp{11} 侧手写结构体时最稳的做法是把 \Tp{vec3} 一律写成
\Tp{vec4}（\Tp{xyzw}）并让 \DomCPU{} 侧也放一个占位的 4 分量量。
\Tp{std430} 允许更紧凑的打包（\Tp{vec3} 不必补满 16 字节、结构体成员按其基础类型
对齐），代价是跨编译器实现细节更敏感，历史项目里常见两种布局混用导致的偏移错位。

没有指针还意味着没有动态内存。数组长度必须是编译期可求的常量表达式；确实需要运行时
尺寸时，唯一的合法出口是 SSBO 末尾的不定长数组 \Tp{data[]}，靠 \Tp{buffer} 块的
实际绑定长度界定有效范围；顶点数量、workgroup 数量这类运行时参数交给
\Tp{gl\_VertexID}、\Tp{gl\_NumWorkGroups} 之类的内建输入去循环，而不是靠变长数组。
\GSL{} 也没有 \Tp{template}、没有重载以外的泛型手段，公共代码只能靠头文件包含
（\Tp{\#include} 由前端处理，Vulkan 工具链支持）或预处理宏复用。

类型检查比 \CPP{} 更严：不存在 \Tp{int} 到 \Tp{float} 的隐式转换，
\Tp{float x = 1;} 是错误，必须写 \Tp{float x = 1.0;}；整型与浮点型混算同样报错；
函数重载按实参类型精确匹配，\Tp{min(int, float)} 找不到候选。位运算只作用于
\Tp{int}/{\Bk}\Tp{uint} 家族，\Tp{float} 需要用 \Fn{packHalf2x16}、\Fn{unpackHalf2x16}
这类打包内建转成整数位模式。矩阵是列主序存储，\Tp{mat4 * vec4} 里向量在右侧；
这与 \HSL{} 默认的 \Tp{mul(v, M)} 写法方向相反，见
第~\ref{ch:ch06-hlsl} 章的对照表。

\begin{warning}[title={采样器不能放进 \Tp{push\_constant}，也不能放进 \Tp{buffer}}]
\Tp{sampler2D} 这类类型是纹理与采样状态的合体：它可以直接进入
\Tp{uniform} 块，但在 Vulkan 语义下不能出现在 \Tp{buffer} 块里，更不能出现在
\Tp{push\_constant} 里，因为描述符句柄必须在绑定阶段就位。要按逐帧数据切换纹理，
把采样器留在 \Tp{uniform} 块、把索引或矩阵放进 push constant 是唯一稳定的分工。
\end{warning}

\section{限定符与存储类：可用阶段、更新代价与对齐要求}

限定符分两族。一族决定数据从哪来、活多久、改动要付什么代价，即存储类：
\Tp{uniform}（块外个体变量或 \Tp{uniform} 块）、\Tp{layout(push\_constant)}、
\Tp{buffer}（可读写 SSBO）、\Tp{shared}（workgroup 内共享内存）、
\Tp{image2D} 一类可写纹理，以及跨阶段的 \Tp{in}/{\Bk}\Tp{out}。另一族只影响存储与数值
行为而不改变生命周期：\Tp{flat}、\Tp{smooth}、\Tp{noperspective} 决定插值方式，
\Tp{centroid} 限定多重采样取点，\Tp{invariant} 保证同一表达式在两个着色器里得到
逐位一致的结果，\Tp{precise} 禁止跨表达式重结合带来的差异，
\Tp{readonly}/{\Bk}\Tp{writeonly}/{\Bk}\Tp{coherent}/{\Bk}\Tp{volatile} 是给 \Tp{buffer} 与
\Tp{image} 的内存访问提示。表~\ref{tab:ch05-storage} 按更新代价把主要存储类排开。

\begin{longtable}{>{\raggedright\arraybackslash}p{0.17\textwidth}
                  >{\raggedright\arraybackslash}p{0.12\textwidth}
                  >{\raggedright\arraybackslash}p{0.24\textwidth}
                  >{\raggedright\arraybackslash}p{0.21\textwidth}}
\caption{GLSL 存储类与限定符：阶段、更新代价与对齐}\label{tab:ch05-storage} \\
\toprule
\textbf{存储类} & \textbf{可用阶段} & \textbf{更新代价与生命周期} & \textbf{对齐与布局} \\
\midrule
\endfirsthead
\toprule
\textbf{存储类} & \textbf{可用阶段} & \textbf{更新代价与生命周期} & \textbf{对齐与布局} \\
\midrule
\endhead
\midrule
\endfoot
\bottomrule
\endlastfoot
块外 \Tp{uniform} & 全部 & 逐个设置，驱动要校验类型与状态；每帧大量改动会打断提交批次 & 无块布局可言，由实现决定内部打包 \\
\Tp{uniform} 块 & 全部（含 compute） & 整块一次上传，适合逐帧矩阵与光照参数；块内容变化会令已绑定描述符失效 & \Tp{std140}：\Tp{vec3} 补到 16 字节，\Tp{mat3} 按三列 \Tp{vec4} \\
\Tp{push\_constant} & 全部 & 不建缓冲、不发描述符，字节直接进命令流，单次代价最低但有总量上限 & 无绑定号；按 \Tp{std430} 风格紧凑排布 \\
\Tp{buffer}（SSBO） & 全部，且可写 & 需分配缓冲、绑定范围与显式屏障；跨绘制复用时最灵活也最容易写错 & \Tp{std430}；不定长数组 \Tp{data[]} 只能是末成员 \\
\Tp{shared} & 仅 compute & 硬件在 workgroup 存活期自动分配，无上传代价 & 尺寸必须是编译期常量，容量受设备上限 \\
\Tp{image2D} 族 & 全部，读写皆可 & 每次读写都是显式内存操作，写完需 \Tp{memory\Bk Barrier\Bk Image()} & 必须声明 \Tp{rgba8} 一类格式；没有自动过滤 \\
\Tp{in}/{\Bk}\Tp{out} & 阶段之间 & 由插值器承载，无运行时更新概念 & \Tp{flat}/{\Bk}\Tp{noper\Bk spective} 改变插值，\Tp{location} 决定槽位 \\
\end{longtable}

代价排序里有一条经验法则：越靠近命令流的数据越便宜，越需要跨阶段可见性的数据越贵。
逐 draw call 都变的少量参数放 \Tp{push\_constant}；整个场景周期不变的矩阵与雾参数放
\Tp{uniform} 块；需要着色器写回、由 CPU 或别的着色器再读的数据放 \Tp{buffer}；
同一 workgroup 内部的中转放 \Tp{shared}。混用会让驱动退化成逐次刷新，\DomCPU{} 侧
看不出来，第~\ref{ch:ch14-perf-tools} 章的性能剖析会立刻暴露。

\section{控制流与并行粒度：写法被硬件约束的部分}

\GSL{} 的控制流语法看着与 \CPP{} 一样，执行模型却不同。片元阶段以 2 乘 2 的像素组
为单位调度，组内分支不是跳转而是掩码执行：两条路径都可能被执行，只是写回被屏蔽。
由此产生两个直接后果。第一，分支体里的纹理取样与导数运算在另一条路径上仍会占用
发射带宽；第二，导数内建依赖同一组内的邻居提供参照值，所以在非一致控制流里调用
\Fn{dFdx}、\Fn{dFdy}、\Fn{fwidth} 得到的是未定义结果——这正是第~\ref{ch:ch02-raster-pipeline}
章里 quad 概念在语言侧的回声。需要逐像素判断后再算导数时，得把判断条件本身保持成
组内一致，或改用显式 LOD 的 \Fn{textureLod} 绕开导数。

compute 阶段的一致性问题落在另一个粒度上。\Fn{barrier} 只同步本 workgroup 内的
invocation，跨 group 的次序必须靠 \Tp{buffer} 上的原子操作与
\Fn{memoryBarrier}、\Fn{bufferMemoryBarrier}、\Fn{groupMemoryBarrier} 显式建立。
循环次数在不同 invocation 之间差异越大，长尾越明显：早退出的 invocation 会占着
同一个执行组的槽位等最慢的那一个。下面这个内核把三种存储类放在一处对照：
\Tp{push\_constant} 提供数量，\Tp{buffer} 承载粒子，\Tp{shared} 做组内归约。

\begin{lstlisting}[style=glslstyle,caption={计算着色器：push constant 读入数量，SSBO 读写数据，shared 做 workgroup 内归约},label={lst:ch05-compute}]
#version 450 core
layout(local_size_x = 64) in;

layout(push_constant) uniform Push {
    uint count;
} push;

layout(std430, binding = 0) buffer Particles {
    vec4 data[];
};

shared float lum[64];

void main()
{
    uint i = gl_GlobalInvocationID.x;
    if (i >= push.count) {
        return;
    }
    vec4 p = data[i];
    lum[gl_LocalInvocationIndex] = dot(p.xyz, vec3(0.2126, 0.7152, 0.0722));
    barrier();

    if (gl_LocalInvocationIndex == 0u) {
        float s = 0.0;
        for (uint k = 0u; k < gl_WorkGroupSize.x; ++k) {
            s += lum[k];
        }
        data[i].w = s / float(gl_WorkGroupSize.x);
    }
}
\end{lstlisting}

读这段代码要注意三处。\Tp{layout(local\_size\_x = 64) in;} 声明的是导入的只读
workgroup 尺寸，写成 \Tp{in} 而没有变量名是 compute 的合法特例。数组
\Tp{data[]} 的长度在编译期未知，越界访问不会有任何检查。\Tp{if (i >= push.count)}
的提前退出保护了共享数组的写入，但归约循环仍按整组 64 项累加，
未初始化项的值来自上一次 kernel——这是真实缺陷，工程上要么先写零，要么把上界改成
实际命中的 invocation 数。

\section{内建变量与常用函数族}

内建变量按阶段划分，全是只读或以 \Tp{gl\_} 前缀保留的可写输出：顶点阶段有
\Tp{gl\_{\Bk}Position}、\Tp{gl\_{\Bk}VertexID}、\Tp{gl\_{\Bk}InstanceID}；片元阶段有
\Tp{gl\_{\Bk}Frag{\Bk}Coord}（窗口坐标，左下原点、\Tp{z} 已透视除完）、
\Tp{gl\_{\Bk}Point{\Bk}Coord}、\Tp{gl\_{\Bk}Front{\Bk}Facing}、可写的 \Tp{gl\_{\Bk}Frag{\Bk}Depth}；几何与
细分阶段有 \Tp{gl\_{\Bk}PrimitiveID} 与 patch 相关量；compute 阶段是一组
\Tp{uvec3}：\Tp{gl\_{\Bk}Work{\Bk}GroupID}、\Tp{gl\_{\Bk}Local{\Bk}InvocationID}、
\Tp{gl\_{\Bk}Global{\Bk}InvocationID}、\Tp{gl\_{\Bk}Num{\Bk}Work{\Bk}Groups}。

函数族大致分六类。几何一族有 \Fn{dot}、\Fn{cross}、\Fn{normalize}、\Fn{length}、
\Fn{distance}、\Fn{faceforward}、\Fn{reflect} 与 \Fn{refract}；矩阵一族有
\Fn{matrix{\Bk}Comp{\Bk}Mult}、\Fn{transpose}、\Fn{determinant} 与 \Fn{inverse}；
通用数值一族有 \Fn{clamp}、\Fn{mix}、\Fn{step}、\Fn{smoothstep}、\Fn{fract} 与
\Fn{modf}；再就是指数对数与三角一族、打包一族（\Fn{pack{\Bk}Half2x16} 为首），
以及原子一族（\Fn{atomicAdd}、\Fn{atomic{\Bk}Comp{\Bk}Swap} 为首，只能作用于 \Tp{buffer} 或
\Tp{atomic\_{\Bk}uint} 计数器）。纹理族的分工最值得记住：\Fn{texture} 隐式求 LOD，
因此依赖片元阶段的导数；\Fn{textureLod} 显式给层级，可以在任何阶段用；
\Fn{textureGrad} 直接给出 \Tp{uv} 方向上的梯度，用于在分支内部手工喂导数；
\Fn{texelFetch} 绕开过滤与归一化坐标，按整数纹素取点；\Fn{textureSize} 只做查询。

下面这个片元着色器把 \Tp{uniform} 块、纹理取样与导数用法放在一起，是本章的
语法总检查：所有声明都合法于 \Tp{450 core}，且刻意只用组内一致的导数调用。搬到
ES 变体时要补两件事：版本行换成 \Tp{300 es}，并给浮点类型加默认的
\Tp{precision} 声明。

\begin{lstlisting}[style=glslstyle,caption={结构完整的片元着色器：UBO 提供参数，textureLod 与 fwidth 决定取样层级},label={lst:ch05-frag}]
#version 450 core

layout(binding = 0) uniform Scene {
    mat4 viewProj;
    vec3 lightDir;
    float exposure;
};

layout(binding = 1) uniform sampler2D albedoTex;
layout(binding = 2) uniform sampler2D detailTex;

layout(location = 0) in vec3 vNormal;
layout(location = 1) in vec2 vUV;
layout(location = 2) flat in uint vMaterialID;

layout(location = 0) out vec4 fragColor;

float roughness(vec2 uv)
{
    return texture(detailTex, uv).r;
}

void main()
{
    if (vMaterialID == 0u) {
        discard;
    }
    vec3 n = normalize(vNormal);
    vec2 uv = vUV;

    float edge = fwidth(uv.x) + fwidth(uv.y);
    float lod = clamp(log2(edge * 256.0), 0.0, 8.0);
    vec4 base = textureLod(albedoTex, uv, lod);

    float ndl = max(dot(n, lightDir), 0.0);
    vec3 rgb = base.rgb * ndl * roughness(uv) * exposure;
    fragColor = vec4(rgb, base.a);
}
\end{lstlisting}

\Fn{fwidth} 与 \Fn{textureLod} 的这对组合是纹理过滤的手写替代：\Tp{edge} 给出
\Tp{uv} 每像素变化量，取对数换成层级，再交给 \Fn{textureLod}。这段之所以安全，
是因为两个 \Fn{fwidth} 调用都处在函数最外层的一致控制流里；把同样的调用挪进
\Tp{if} 分支的某一条路径，结果就不再可预测。第~\ref{ch:ch12-lighting-pbr} 章里
各向异性过滤与 mip 选择的讨论会回到这个式子。

\begin{bestpractice}[title={写 \GSL{} 之前的四条落地约定}]
\begin{itemize}
\item 源码首行固定 \Tp{\#version 450 core}，需要兼容移动端时另开一份
\Tp{300 es} 变体，而不是在同一文件里用宏裁语法。
\item 所有跨 \DomCPU{}/{\Bk}\DomGPU{} 边界的块显式写 \Tp{std140} 或 \Tp{std430}，
并在 \DomCPU{} 侧用 16 字节对齐的镜像结构体，禁止依赖默认布局。
\item 只用 \Tp{flat} 传整数与索引进片元阶段，避免非插值语义在别处被误读。
\item 需要 LOD 控制的取样一律写成 \Fn{textureLod} 或 \Fn{textureGrad}，
把隐式导数的使用点留在最外层控制流。
\end{itemize}
\end{bestpractice}

% ============================================================
% ---- frag_ch06.tex ----
\chapter{HLSL 与着色模型}
\label{ch:ch06-hlsl}

\HSL{} 是 \Lib{Direct3D} 侧的着色语言，语法上更接近 \CL{} 家族：显式类型、显式寄存器
槽、把接口写进变量名后面的冒号里。它与 \GSL{} 的核心分歧不在算子，而在约定的位置：
\GSL{} 把约定藏在限定符与 \Tp{\#version} 里，\HSL{} 把约定摊开在语义（semantic）与
着色模型（shader model）里。本章按这条线索展开，最后与第~\ref{ch:ch05-glsl} 章
逐条对照。

\section{着色模型：能力台阶与寄存器空间}

着色模型是 \HSL{} 的能力版本号，写成 \Tp{vs\_5\_0}、\Tp{ps\_6\_0}、\Tp{cs\_6\_5}
这样的阶段前缀加数字：\Tp{vs} 顶点、\Tp{ps} 片元、\Tp{gs} 几何、\Tp{hs} 与 \Tp{ds}
细分的两级、\Tp{cs} 计算。这个数字在编译时决定两件事：能用哪些语法与内建，以及能占
多少硬件资源。台阶是硬性的——写 \Tp{double} 需要 \Tp{ps\_5\_0} 起，用 wave 内建需要
\Tp{cs\_6\_0} 起，\Tp{min16float} 一族要 \Tp{ps\_5\_1} 或更高。

寄存器空间是最容易被低估的一层。\HSL{} 的每个资源声明都落在一个带字母的槽位里：
\Tp{b} 是常量缓冲、\Tp{t} 是纹理、\Tp{s} 是采样器、\Tp{u} 是可读写视图、\Tp{c}
是老的立即常数寄存器。台阶之间的差别很大：\Tp{ps\_4\_0} 与 \Tp{ps\_4\_1} 每个阶段
只有 256 个 16 字节常数寄存器，\Tp{ps\_5\_0} 起涨到 4096；\Tp{b0} 到 \Tp{b13}、
\Tp{t0} 到 \Tp{t127}、\Tp{s0} 到 \Tp{s15} 是长期稳定的上限，\Tp{u} 槽在
\Tp{ps\_5\_0} 只有 8 个，\Tp{ps\_5\_1} 与描述符绑定时代扩到 64。超过上限不是运行时
报错，而是编译阶段直接拒绝。

\begin{note}[title={台阶同时决定工具与运行时}]
\Tp{ps\_5\_0} 及以下的目标串由老编译器处理，输出 \Lib{DXBC}；\Tp{ps\_6\_0} 及以上
只能交给新编译器，输出 \Lib{DXIL}，并且在设备侧还要求对应的 wave 运算、位宽与描述符
能力。同一份源码换目标串就可能编译不过，这与 \GSL{} 换 \Tp{\#version} 的情况同构，
见本章第~\ref{sec:ch06-bytecode} 节。
\end{note}

\section{\Tp{cbuffer} 的打包规则与 \GSL{} 的 \Tp{uniform} 块}

\HSL{} 没有 \Tp{layout(std140)} 这种东西，它的对齐由一条规则统辖：一切按 16 字节的
寄存器为单位。具体后果有五条。标量按 \Tp{x}、\Tp{y}、\Tp{z}、\Tp{w} 顺序往一个
寄存器里填，填满才进下一个；向量与不可拆分标量的组合，一个 \Tp{float3} 会独占整个
寄存器（它不能跨寄存器边界）；数组的每个元素都从新寄存器的 \Tp{x} 开始，因此
\Tp{float w[4]} 占 64 字节而不是 16 字节；结构体成员同样对齐到寄存器边界；块的整体
大小向上取整到 16 的倍数。

\GSL{} 的 \Tp{std140} 在 \Tp{vec3} 与标量数组这两点上与 \HSL{} 结论相同，但
\Tp{std430} 允许紧凑打包，于是同一份数据在 \Tp{cbuffer}、\Tp{std140}、\Tp{std430}
三种约定下会有三种偏移。两边都没有语言级的偏移查询内建，现实做法有两条：把每个字段
都写成四分量量，让三种约定的结果重合；或者在 \DomCPU{} 侧用序列化函数按规则逐字段算
偏移，并对块的总大小做静态断言。老 \HSL{} 提供 \Tp{packoffset(c1.x)} 直接指定位置，
这在后来的工具链里已不再被支持，遗留代码读到时要当成分层负担看待。

\section{语义系统：名字即接口}

\HSL{} 的变量声明后面可以跟一个冒号加标识符，这就是语义。它同时承担三个角色：
顶点输入布局的名字、阶段之间传输通道的标识、以及系统级输入输出的入口。普通语义
\Tp{POSITION}、\Tp{NORMAL}、\Tp{COLOR}、\Tp{TEXCOORD0} 到 \Tp{TEXCOORD15} 属于
前两类；带 \Tp{SV\_} 前缀的系统语义属于第三类，其中 \Tp{SV\_Position} 是 VS 输出的
裁剪空间坐标（写它而不是 \Tp{POSITION} 是新台阶的规矩），\Tp{SV\_Target} 或
\Tp{SV\_Target0} 到 \Tp{SV\_Target7} 是渲染目标，\Tp{SV\_Depth} 覆写深度，
\Tp{SV\_Coverage} 覆写逐位覆盖遮罩（与保守光栅配合），计算阶段读
\Tp{SV\_DispatchThreadID}、\Tp{SV\_GroupID}、\Tp{SV\_GroupThreadID}，几何与细分
阶段读 \Tp{SV\_VertexID}、\Tp{SV\_InstanceID}、\Tp{SV\_PrimitiveID}。

关键点在于阶段间的匹配是按名字与语义做的，而不是按槽位编号。\GSL{} 用
\Tp{layout(location = 0)} 把插值通道钉在编号上，\HSL{} 则要求 VS 的输出结构与 PS
的输入结构在字段名与语义上一致；不一致时编译器只给警告，运行时读到的是未定义值。
这套设计让同一份 PS 可以配多个 VS，代价是重构字段名时没有任何静态保证。第~\ref{ch:ch03-vertex-input}
章讲顶点输入布局时，语义与顶点流就是靠名字串起来的。

下面这对 VS 与 PS 是本章的语法总样本：\Tp{cbuffer} 提供逐帧参数，\Tp{Texture2D}
与 \Tp{SamplerState} 分离声明，输出靠语义而不是位置编号。

\begin{lstlisting}[style=hlslstyle,caption={完整的 VS 与 PS：cbuffer、Texture2D 取样、语义化的输入输出},label={lst:ch06-vs-ps}]
struct VSInput {
    float3 pos : POSITION;
    float3 nrm : NORMAL;
    float2 uv  : TEXCOORD0;
};

struct VSOutput {
    float4 pos : SV_Position;
    float3 nrm : NORMAL;
    float2 uv  : TEXCOORD0;
};

cbuffer Frame : register(b0) {
    float4x4 viewProj;
    float3   lightDir;
    float    exposure;
};

Texture2D    albedo  : register(t0);
Texture2D    detail  : register(t1);
SamplerState samp    : register(s0);

VSOutput VSMain(VSInput input)
{
    VSOutput o;
    o.pos = mul(float4(input.pos, 1.0), viewProj);
    o.nrm = normalize(input.nrm);
    o.uv  = input.uv;
    return o;
}

float4 PSMain(VSOutput input) : SV_Target
{
    float4 base = albedo.Sample(samp, input.uv);
    float rough = detail.Sample(samp, input.uv).r;
    float ndl = saturate(dot(input.nrm, lightDir));
    float3 rgb = base.rgb * ndl * rough * exposure;
    return float4(rgb, base.a);
}
\end{lstlisting}

\begin{warning}[title={\Tp{register} 号与 API 侧的绑定必须一致}]
\Tp{register(b0)} 只是源码里的槽位声明：\Lib{Direct3D} 12 的实际绑定由根签名与
描述符表决定，两边编号语义不同。最常见的错误是源码写 \Tp{register(t3)} 而根签名只
留了三个纹理槽，编译通过、运行时读到空描述符。迁移到 \Tp{ps\_6\_0} 之后，槽号更
只是占位，真正的位置来自绑定时的描述符堆。
\end{warning}

\section{资源类型族：纹理、缓冲与采样器}

\HSL{} 把资源分成三类。
\begin{itemize}
  \item 只读纹理：\Tp{Texture1D}、\Tp{Texture2D}、\Tp{Texture3D}、\Tp{TextureCube}、
        \Tp{Texture2D{\Bk}Array} 与 \Tp{Texture2DMS}。
  \item 取样方法：\Tp{Sample}、\Tp{SampleLevel}、\Tp{Sample{\Bk}Bias}、\Tp{Sample{\Bk}Grad}、
        \Tp{SampleCmp}、\Tp{Gather} 与 \Tp{Load}；元信息靠 \Tp{Get{\Bk}Dimensions}，
        自身可算层级用 \Tp{Calculate{\Bk}Level{\Bk}Of{\Bk}Detail}。
  \item 缓冲：\Tp{Buffer<float4>} 这类带类型的视图、\Tp{Structured{\Bk}Buffer} 与它的可写形态
        \Tp{RW{\Bk}Structured{\Bk}Buffer}、按字节寻址的 \Tp{Byte{\Bk}Address{\Bk}Buffer} 与
        \Tp{RWByte{\Bk}Address{\Bk}Buffer}，以及配合流式输出的 \Tp{Append{\Bk}Structured{\Bk}Buffer}
        与 \Tp{Consume{\Bk}Structured{\Bk}Buffer}。
  \item 采样状态：\Tp{Sampler{\Bk}State} 与 \Tp{Sampler{\Bk}Comparison{\Bk}State}。
\end{itemize}

采样器与纹理分离是 \HSL{} 早于 \GSL{} 落地的设计（老式 \Tp{sampler2D} 组合声明在
新台阶里已不使用），它带来两个后果：同一个纹理可以在一次绘制里用不同采样器多次取样；
\DomCPU{} 侧必须显式创建并绑定采样器对象。这与 \GSL{} 的 \Tp{sampler2D} 正好相反——
那边纹理与过滤状态被合成一个句柄类型，而 \SPIV{} 里则要拆成
\Tp{OpType\Bk SampledImage}，见第~\ref{ch:ch07-spirv-toolchain} 章。

\section{DXBC 与 DXIL：两代字节码}
\label{sec:ch06-bytecode}

目标串 \Tp{vs\_4\_0} 到 \Tp{ps\_5\_1} 这一段产出 \Lib{DXBC}：一个带容器头的字节码
流，指令以 \Tp{dcl\_}、\Tp{mov}、\Tp{ld\_} 这类前缀形式出现，反汇编工具与调试器都
围绕它设计。切到 \Tp{ps\_6\_0} 之后产出换成 \Lib{DXIL}：底层是 \Lib{LLVM} 位码，
模块里带版本元数据，加载前必须通过运行库的校验，因此调试与反射接口也跟着换了一代。
对 \DomCPU{} 侧程序员的可见差别有三条：\Lib{DXIL} 支持资源数组的动态索引与更灵活的
控制流；\Lib{DXIL} 的编译单元里可以放多个入口，靠 \Tp{-E} 选择；\Lib{DXIL} 需要随包
分发校验用的运行时组件，部署时多一个版本匹配问题。

命令行层面，老接口是 \Tp{fxc /T ps\_5\_0 /E PSMain /Fo ps.obj ps.hlsl}，新接口是
\Tp{dxc -T ps\_6\_0 -E PSMain -Fo ps.dxil ps.\Bk hlsl}，加上 \Tp{-Zi} 生成调试信息、
\Tp{-O3} 提高优化级别、\Tp{-D LOWPRE=1} 注入宏。想让新编译器也吐 \GSL{} 侧可读的中
间形式，就走 \SPIV{} 目标，见第~\ref{ch:ch07-spirv-toolchain} 章。

\section{\GSL{} 与 \HSL{}：同一算法逐条对照}

把第~\ref{ch:ch05-glsl} 章那个 \Tp{shared} 归约内核改写成 \HSL{}，最直观的收益是
wave 内建：一条 \Tp{Wave\Bk Active\Bk Sum} 取代了共享数组加屏障的整段样板。

\begin{lstlisting}[style=hlslstyle,caption={计算内核：wave 归约与组内归约两条路径，由宏切换},label={lst:ch06-cs}]
StructuredBuffer<float4>    src   : register(t0);
RWStructuredBuffer<float>   total : register(u0);
groupshared float gsum[64];

#define USE_WAVE 1

[numthreads(64, 1, 1)]
void CSMain(uint3 dtid : SV_DispatchThreadID,
            uint3 tid  : SV_GroupThreadID)
{
    float v = src[dtid.x].w;

#if USE_WAVE
    float laneSum = WaveActiveSum(v);
    if (WaveIsFirstLane()) {
        InterlockedAdd(total[0], laneSum);
    }
#else
    gsum[tid.x] = v;
    AllMemoryBarrierWithGroupSync();
    if (tid.x == 0) {
        float s = 0.0;
        [unroll]
        for (uint i = 0u; i < 64u; ++i) {
            s += gsum[i];
        }
        InterlockedAdd(total[0], s);
    }
#endif
}
\end{lstlisting}

\GSL{} 一侧到 \Tp{450} 为止都没有 wave 内建，等价写法只有共享内存加分段归约，再
用 \Tp{atomicAdd} 提交结果：

\begin{lstlisting}[style=glslstyle,caption={同一个归约的 \GSL{} 写法：shared 分段相加后原子提交},label={lst:ch06-cs-glsl}]
#version 450 core
layout(local_size_x = 64) in;

layout(binding = 0) readonly buffer Src {
    vec4 data[];
} src;

layout(binding = 1) buffer Total {
    float sum;
} acc;

shared float lane[64];

void main()
{
    lane[gl_LocalInvocationIndex] = src.data[gl_GlobalInvocationID.x].w;
    barrier();

    for (uint s = 32u; s > 0u; s /= 2u) {
        uint k = gl_LocalInvocationIndex;
        if (k < s) {
            lane[k] += lane[k + s];
        }
        barrier();
    }

    if (gl_LocalInvocationIndex == 0u) {
        atomicAdd(acc.sum, lane[0]);
    }
}
\end{lstlisting}

两份代码的行为差别不只在行数：\HSL{} 的 \Tp{WaveActiveSum} 作用域是硬件执行组，
宽度由设备决定，与 \Tp{[numthreads(64, 1, 1)]} 无关；\GSL{} 的 \Tp{shared} 数组
作用域是 workgroup，必须自己保证所有 invocation 都越过了写入点。
表~\ref{tab:ch06-glsl-hlsl} 把两侧最常踩的差异列出来。

\begin{longtable}{>{\raggedright\arraybackslash}p{0.12\textwidth}
                  >{\raggedright\arraybackslash}p{0.25\textwidth}
                  >{\raggedright\arraybackslash}p{0.25\textwidth}
                  >{\raggedright\arraybackslash}p{0.10\textwidth}}
\caption{\GSL{} 与 \HSL{} 语法与语义差异速查}\label{tab:ch06-glsl-hlsl} \\
\toprule
\textbf{主题} & \textbf{\GSL{}} & \textbf{\HSL{}} & \textbf{陷阱} \\
\midrule
\endfirsthead
\toprule
\textbf{主题} & \textbf{\GSL{}} & \textbf{\HSL{}} & \textbf{陷阱} \\
\midrule
\endhead
\midrule
\endfoot
\bottomrule
\endlastfoot
入口函数 & \Tp{void main()}，阶段由编译选项给 & 任意名字，如 \Tp{VSMain}，配 \Tp{-E} & 名字冲突只在链接期发现 \\
顶点索引 & \Tp{gl\_VertexID} & \Tp{SV\_VertexID} & 一边是变量，一边是语义 \\
裁剪坐标 & 写 \Tp{gl\_Position} & 写 \Tp{SV\_Position} & 普通语义不参与裁剪 \\
片元输出 & \Tp{out} 加 \Tp{location} & 返回值语义 \Tp{SV\_Target} & \GSL{} 可多个 \Tp{out}，编号要匹配混合配置 \\
深度写入 & \Tp{gl\_FragDepth} & \Tp{SV\_Depth} & 一旦赋值就完全由着色器负责 \\
丢弃 & \Tp{discard} & \Tp{discard} 或 \Tp{clip(x)} & 导数在丢弃后的组内不可靠 \\
常量块 & \Tp{uniform} 块 & \Tp{cbuffer} & 打包规则见 §2 \\
绑定方式 & \Tp{binding} 编号 & \Tp{register} 槽号 & 新台阶里槽号只是占位 \\
精度控制 & \Tp{highp}、\Tp{mediump} & \Tp{min16float} 一族 & 只在移动端与低精度台阶有实义 \\
矩阵乘 & \Tp{M * v}，列主序 & \Tp{mul(v, M)}，默认行主序 & 换编译器要重设主序开关 \\
纹理取样 & \Tp{texture(t, uv)} & \Tp{t.Sample(s, uv)} & 采样状态是否合体 \\
导数 & \Tp{dFdx}、\Tp{fwidth} & \Tp{ddx}、\Tp{ddy}、\Tp{width} & 两边都要求组内一致 \\
\end{longtable}

\begin{compare}[title={同一个乘法的四种写法}]
\GSL{} 写 \Tp{gl\_Position = proj * view * vec4(p, 1.0);}，向量在右、矩阵列主序。
\HSL{} 默认写 \Tp{mul(float4(p, 1.0), WVP)}，向量在左、矩阵行主序；两边在
\SPIV{} 里都是 \Tp{Op\Bk Matrix\Bk Times\Bk Vector} 这条指令，但元素排列取决于编译
开关。跨引擎代码里
最稳的约定是：把矩阵统一存成行主序，\GSL{} 侧用 \Tp{transpose} 一次，或者在
\DomCPU{} 侧上传前就转置好。
\end{compare}

\begin{bestpractice}[title={从 \GSL{} 迁到 \HSL{} 的三条检查}]
\begin{itemize}
\item 目标串先定：\Tp{ps\_5\_0} 用老编译器与 \Lib{DXBC}，\Tp{ps\_6\_0} 起必须换
\Lib{DXIL} 并确认设备支持所用台阶。
\item 常量块的每个字段按 16 字节起址画一遍表，标出被补出来的洞，再让
\DomCPU{} 侧的镜像结构体对齐这些偏移。
\item 阶段之间的每个字段都要有语义：普通数据用 \Tp{TEXCOORD} 编号，系统量用
\Tp{SV\_} 前缀，两边结构与名字逐字一致，别依赖编译器帮你检查。
\end{itemize}
\end{bestpractice}

% ============================================================
% ---- frag_ch07.tex ----
\chapter{SPIR-V 与跨语言编译链}
\label{ch:ch07-spirv-toolchain}

\SPIV{} 是 \Lib{Khronos} 的着色器与计算内核中间表示：定长字流的字节码，加上一套
文本反汇编形式。它在 \DomGPU{} 侧的地位类似 \Lib{LLVM} IR 在 \DomCPU{} 侧的位置——
不是给程序员写的，而是让多种前端语言与多个后端 API 在一处会合。本章讲它的模块结构、
围绕它的工具链、把它做进构建系统的方式，以及跨语言方案实际做不到什么。

\section{\SPIV{} 的模块结构：编号、类型声明与存储类}

一个 \SPIV{} 模块是一串 32 位字。每个指令第一个字的高 16 位是操作码、低 16 位是字
数，后面跟操作数。模块顺序是硬性的：先是五个字的头（魔数 \Tp{0x07230203}、版本、
生成器号、ID 上界、schema 保留字），然后一次性声明所有能力（\Tp{OpCapability}，
例如 \Tp{Shader}、\Tp{VulkanMemoryModel}、\Tp{Int64}），再是内存模型、入口点声明与
执行模式，最后才是类型、常量、变量与函数体。除函数可以先声明后定义之外，任何引用
都必须指向已经出现过的 ID，所以类型一定先于使用该类型的函数出现。

ID 系统是最需要先适应的部分：每个类型、常量、变量、函数、基本块与指令结果都占一个
数字 ID，文本形式里写成 \Tp{\%17}、\Tp{\%main} 这样的名字。函数由 \Tp{OpFunction}
开始、若干 \Tp{OpLabel} 分隔的基本块组成、以 \Tp{OpReturn} 或
\Tp{OpBranch} 结束，参数是 \Tp{OpFunctionParameter}，块内汇聚点用 \Tp{OpPhi}。
控制流必须结构化：不允许多入口块，也不允许任意跳转，这与 \GSL{}/{\Bk}\HSL{} 前端的
循环展开策略直接相关。

变量声明 \Tp{OpVariable} 带一个存储类参数，它就是第~\ref{ch:ch05-glsl} 章那些限定符
的落地形态：\Tp{UniformConstant} 承载纹理与采样器，\Tp{Uniform} 承载常量缓冲，
\Tp{PushConstant} 承载推送常量，\Tp{StorageBuffer} 承载 \Tp{buffer} 块，
\Tp{Workgroup} 对应 \Tp{shared}，\Tp{Private} 与 \Tp{Function} 是全局和函数局部
临时量，\Tp{Input} 与 \Tp{Output} 覆盖跨阶段变量。跨阶段与跨 API 的可见性
差异，在 \SPIV{} 层面就压缩成这几个类别加上一串装饰（\Tp{OpDecorate} 给出的
\Tp{Binding}、\Tp{DescriptorSet}、\Tp{Location}、\Tp{BuiltIn}）。

把模块读成文本靠 \Tp{spirv\Bk-dis}，反向靠 \Tp{spirv\Bk-as}：

\begin{lstlisting}[style=shellstyle,caption={反汇编片段：入口、执行模式与两种存储类的变量声明},label={lst:ch07-dis}]
; SPIR-V
; Version: 1.5
; Generator: Khronos Glslang Reference Front End
       OpCapability Shader
       OpMemoryModel Logical GLSL450
 OpEntryPoint Fragment %main "main" %fragColor
       OpExecutionMode %main OriginUpperLeft
  %v4f = OpTypeVector %float 4
   %ps = OpTypePointer PushConstant %v4f
 %push = OpVariable %ps PushConstant
   %fn = OpFunction %void None %voidfn
  %lbl = OpLabel
       OpReturn
       OpFunctionEnd
\end{lstlisting}

这段里有三处工程上真正会去查的东西：\Tp{Op\Bk EntryPoint} 的名字决定 API 侧加载时
找的入口符号；\Tp{Op\Bk ExecutionMode} 决定 \Tp{gl\_FragCoord} 的原点在左上还是左下
（\Tp{Origin\Bk UpperLeft} 与 \Lib{OpenGL} 习惯相反）；\Tp{OpVariable} 的存储类决定
描述符布局该把这条变量放进哪个集合。调试三角形朝上、顶点语义错位、缓冲偏移偏一
个字节，三类问题的现场都在这里。

\section{编译链全景：每种源语言的前端}

\SPIV{} 的上游有四个现实前端。\GSL{} 走 \Lib{glslang}（官方参考实现）或包了一层的
\Lib{shaderc}；\HSL{} 走 \Lib{DXC}（新编译器）或 \Lib{slang}；\MSL{} 走 Apple 自带
编译器，它产出平台私有的 \Tp{AIR}，需要打包步骤；\Lib{Naga} 作为 \Lib{Rust}
实现，接受 \Tp{WGSL} 并把 \SPIV{} 作为主要输出之一。\Lib{Direct3D} 老编译器只出
\Lib{DXBC}，不会碰 \SPIV{}。工具之间的分工见表~\ref{tab:ch07-tools}，命令行层面
的典型序列是下面这段。

\begin{longtable}{>{\raggedright\arraybackslash}p{0.21\textwidth}
                  >{\raggedright\arraybackslash}p{0.24\textwidth}
                  >{\raggedright\arraybackslash}p{0.14\textwidth}
                  >{\raggedright\arraybackslash}p{0.13\textwidth}}
\caption{\SPIV{} 编译链上的工具对比}\label{tab:ch07-tools} \\
\toprule
\textbf{工具} & \textbf{输入到输出} & \textbf{许可} & \textbf{典型用途} \\
\midrule
\endfirsthead
\toprule
\textbf{工具} & \textbf{输入到输出} & \textbf{许可} & \textbf{典型用途} \\
\midrule
\endhead
\midrule
\endfoot
\bottomrule
\endlastfoot
\Tp{glslang\Bk Validator} & \GSL{} 到 \SPIV{} & BSD 风格 & 官方参考前端 \\
\Tp{glslc}（\Lib{shaderc}） & \GSL{} 到 \SPIV{} & Apache 2.0 & 构建期批处理 \\
\Tp{slangc}（\Lib{slang}） & 三种语言到 \SPIV{} 或 DXIL & NCSA 风格 & 统一前端与模块化编译 \\
\Tp{dxc} & \HSL{} 到 DXIL 或 \SPIV{} & NCSA 风格 & \Tp{ps\_6\_0} 起的编译器 \\
\Tp{fxc} & \HSL{} 到 DXBC & 闭源（SDK） & 老台阶产物 \\
\Tp{spirv\Bk-cross} & \SPIV{} 到 \GSL{}、\HSL{}、\MSL{} & Apache 2.0 或 MIT & 反向翻译 \\
\Tp{spirv\Bk-opt} & \SPIV{} 到 \SPIV{} & Apache 2.0 & 发布前瘦身 \\
\Tp{spirv\Bk-val} & 只读校验 & Apache 2.0 & 提交前的合法性检查 \\
\Tp{spirv\Bk-dis}、\Tp{spirv\Bk-as} & 文本到 \SPIV{} & Apache 2.0 & 复现缺陷与手写测试 \\
\Tp{spirv\Bk-link}、\Tp{spirv\Bk-remap} & 多模块合并 & Apache 2.0 & 把库拆成多文件再合 \\
\Tp{metal} 编译器 & \MSL{} 到 \Tp{AIR} & 闭源（Apple） & Apple 平台离线编译 \\
\Tp{naga} & \Tp{WGSL} 到 \SPIV{} 等 & 随仓库（\Lib{Rust}） & \Lib{wgpu} 与 \Lib{Rust} 管线 \\
\end{longtable}

\begin{lstlisting}[style=shellstyle,caption={一次完整的往返：编译、反汇编、再翻译、再编译、优化与校验},label={lst:ch07-roundtrip}]
# 1) 前端：GLSL 源进，SPIR-V 出（-V 是 Vulkan 方言模式）
glslangValidator -V -S frag -o pbr.frag.spv pbr.frag
glslc --target-env=vulkan1.1 -O0 -o blur.spv blur.comp

# 2) 反汇编成文本，检查入口与存储类，再汇编回去验证可逆
spirv-dis --no-header pbr.frag.spv -o pbr.frag.dis
spirv-as pbr.frag.dis -o pbr.frag.back.spv

# 3) 反向翻译：同一份模块喂给 HLSL 与 MSL 两个后端
spirv-cross --hlsl --shader-model 50 pbr.frag.spv \
    --output pbr.hlsl --entry-point main
spirv-cross --msl blur.spv --msl-version 200 --output blur.msl

# 4) 再编译：HLSL 侧直接出 SPIR-V，回到中间表示
dxc -T ps_6_0 -E PSMain -spirv -Fo ps.spv ps.hlsl
slangc ps.hlsl -target spirv -stage fragment \
    -entry PSMain -o ps.slang.spv

# 5) 发布前的优化与校验
spirv-opt -O --legalize-hlsl ps.spv -o ps.opt.spv
spirv-val --target-env vulkan1.1 ps.opt.spv
\end{lstlisting}

\section{反向翻译与跨语言层的现实边界}

\Lib{SPIRV-Cross} 是这条链上唯一的反向桥梁：拿 \SPIV{} 模块，吐 \GSL{}、\HSL{} 或
\MSL{} 源码，也可以只吐反射信息给 \DomCPU{} 侧生成绑定表。它的价值在于把编写语言与
目标 API 解耦——只维护一份 \GSL{}，再自动得到 \HSL{} 与 \MSL{} 变体。

边界也很清楚。第一，\SPIV{} 里没有源语言信息，反翻译出来的代码是可编译的机器风格
源码，不是人写的样子，宏、头文件结构与命名全部丢失，只适合进构建产物而不适合回手写
树。第二，坐标系与深度范围是平台约定，翻译器不会替你改：\Lib{Vulkan} 的裁剪空间
\Tp{z} 区间与 \Lib{OpenGL} 不同，屏幕原点又与 \HSL{} 习惯相反，这类差异要人工在
投影矩阵或 \Tp{Op\Bk ExecutionMode} 层面处理，第~\ref{ch:ch04-transform-space} 章
的矩阵推导就是给这一步用的。第三，语义要重新对齐：\GSL{} 的 \Tp{layout(location)}
到了 \HSL{} 侧会变成 \Tp{SV\_Target} 之类的语义，采样器与纹理的分合方式也要靠
选项显式声明，否则 \Tp{Op\Bk Type\Bk SampledImage} 无法落进目标语言的类型系统。

\begin{warning}[title={跨语言层吞不掉的三件事}]
\begin{itemize}
\item 一致性要求：\HSL{} 的动态资源索引必须带 \Tp{nonuniform\Bk resource} 属性，
翻译到 \SPIV{} 时要开启 \Tp{Runtime\Bk DescriptorArray} 与
\Tp{NonUniform\Bk Resource} 两个能力；\GSL{} 侧没有对应语法糖，只能改写成一致的
索引或分成多次取样。
\item 浮点位数：\Tp{mediump}、\Tp{min16float} 这类低精度承诺是各自语言的，
\SPIV{} 只记录类型宽度，跨后端映射时可能悄悄变宽，数值结果随之改变。
\item 调试信息：名称、行号在优化与翻译中被反复重写，出问题时以
\Tp{spirv\Bk-dis} 文本为准，别相信任何一层的源码行号。
\end{itemize}
\end{warning}

\section{优化、特化常量与打包校验}

\Tp{spirv\Bk-opt} 的入口是 \Tp{-O}（整体策略）与 \Tp{-Os}（体积优先），也可以给
配置文件或一串单 pass：\Tp{-\/-eliminate-dead-branches}、\Tp{-\/-merge-blocks}、
\Tp{-\/-unroll-loops}、\Tp{-\/-legalize-hlsl}（把 \HSL{} 前端产物整形成
\Lib{Direct3D} 12 能接受的形式）、\Tp{-\/-strip-debug}。它做的是模块内的等价变换，
跨模块的符号合并不是它的职责，那是 \Tp{spirv\Bk-link} 的活：把多个模块并成一个，
重命名冲突 ID，并按选项合并描述符集合。

真正影响架构决策的是特化常量：值推迟到创建管线时才给。Vulkan 侧的 \GSL{} 方言写法
是 \Tp{layout(constant\_id = n)}，声明仍然是 \Tp{const}，但编译器不会替它折叠后续
循环，于是同一个模块能配出多个质量的变体。

\begin{lstlisting}[style=glslstyle,caption={用特化常量控制卷积核半径：编译期保留循环结构，管线创建期定值},label={lst:ch07-spec}]
#version 450 core

// radius is fixed when the pipeline is created, not at compile time
layout(constant_id = 0) const int radius = 3;

layout(local_size_x = 8, local_size_y = 8) in;

layout(binding = 0) uniform sampler2D srcTex;

layout(binding = 1) writeonly buffer Dst {
    vec4 data[];
} dst;

void main()
{
    ivec2 c = ivec2(gl_GlobalInvocationID.xy);
    int r = radius;
    vec3 acc = vec3(0.0);

    for (int y = -r; y <= r; ++y) {
        for (int x = -r; x <= r; ++x) {
            acc += texelFetch(srcTex, c + ivec2(x, y), 0).rgb;
        }
    }

    float n = float((2 * r + 1) * (2 * r + 1));
    dst.data[c.y * 512 + c.x] = vec4(acc / n, 1.0);
}
\end{lstlisting}

\DomCPU{} 侧把这些值填进 \Tp{Vk\Bk SpecializationInfo}（一组 \Tp{const\_id} 到字节偏移
的映射），管线创建时后端会先做常量折叠再选寄存器分配。代价是模块体积不能提前缩减，
而且 \Tp{spirv\Bk-opt} 在常量未知时不会展开这些循环。另一侧，\HSL{} 没有等价语言
构造：同样的效果要么在 \Tp{cbuffer} 里传参数并保留动态循环，要么编译期用宏切变体。

打包与校验的规则集中在三点。版本兼容：模块头的版本字段必须不高于目标环境支持的
\SPIV{} 版本，核心 API 每提升一代都会放宽这个上界，具体对应关系写在规范的一致性
附录里；产物里用了更高版本才有的指令，模块创建时就会失败。合法性：\Tp{spirv\Bk-val}
加 \Opt{target-env} 检查能力声明是否齐全、存储类用法是否越界、结构控制流是否成立；
开发期再叠一层运行时的验证层，两者覆盖面不同，不能互相替代。复用：同一份
\SPIV{} 可以在多个描述符布局之间搬运，靠 \Tp{spirv\Bk-remap} 或
\Tp{Op\Bk Decorate} 重写 \Tp{Binding} 与 \Tp{DescriptorSet}，而不必回源语言重编。

\begin{guideline}[title={模块发布的三条硬规则}]
\begin{itemize}
\item 进包的只有优化后的 \SPIV{} 或目标 API 的字节码，源码与中间文本留在仓库，
产物里不出现可反推业务逻辑的字符串。
\item 每个模块记录四元组：源语言版本、前端工具与版本、目标环境、\Opt{target-env}
校验结果，缺一项就不允许被打包脚本收进资产表。
\item 入口名保持稳定：入口符号在 \Tp{Op\Bk EntryPoint} 里是字符串，改名等于换
管线，别让它跟着文件重命名一起漂移。
\end{itemize}
\end{guideline}

\section{把着色器编译做进构建与 CI}

着色器一旦数量超过十几个，手工编译就会变成事故源：改了 \Tp{.comp} 忘了重编、CI
上的产物与本地不一致。做法是把每个源文件变成一个构建产物，依赖关系交给工具。

\begin{lstlisting}[style=cmakestyle,caption={CMake 里把 GLSL 逐文件编成 SPIR-V：源文件变更只触发对应产物},label={lst:ch07-cmake}]
find_program(GLSLANG glslangValidator REQUIRED)
find_program(SPIRV_OPT spirv-opt REQUIRED)

set(SHADER_SOURCES pbr.frag pbr.vert blur.comp)
set(SHADER_BINS "")

foreach(src IN LISTS SHADER_SOURCES)
    get_filename_component(ext "${src}" LAST_EXT)
    set(out "${CMAKE_BINARY_DIR}/shaders/${src}.spv")
    add_custom_command(
        OUTPUT "${out}"
        COMMAND ${GLSLANG} -V -S ${ext} -o "${out}" "${src}"
        COMMAND ${SPIRV_OPT} -O "${out}" -o "${out}"
        DEPENDS "${src}"
        COMMENT "spirv for ${src}"
        VERBATIM
    )
    list(APPEND SHADER_BINS "${out}")
endforeach()

add_custom_target(shaders ALL DEPENDS ${SHADER_BINS})
target_compile_definitions(render PRIVATE
    SHADER_DIR="${CMAKE_BINARY_DIR}/shaders")
\end{lstlisting}

这段有四个要点。\Tp{foreach} 让增量粒度落在单个文件上，一个内核改坏不会连带全部
重编。\Tp{LAST\_EXT} 的后缀直接喂给 \Tp{-S}，前提是命名约定统一（\Tp{.frag}、
\Tp{.vert}、\Tp{.comp}）。\Tp{COMMAND} 链里把优化接在编译之后，失败即中断该目标。
\Tp{VERBATIM} 保证路径里的空格不被二次解析。

\begin{note}[title={CI 里让着色器步骤可判定}]
把校验做成独立步骤，失败信息才有定位价值：先跑前端编译，再跑
\Tp{spirv\Bk-val}，最后比对产物哈希是否变化。缓存键用源文件与前端版本的联合哈希，
否则升级编译器后 CI 会静默复用旧产物。
\end{note}

变体爆炸是同一问题的另一半。按宏组合切变体时，\Tp{quality}、\Tp{shadow}、
\Tp{alpha} 三个开关就能出八份模块，再加上目标环境就是十六份。可执行的做法是把变体
表外置成数据文件、让构建脚本读它，并且规定：能用特化常量表达的维度，不生成第二份
模块；只有影响指令选择的维度才允许切成独立变体。

\begin{bestpractice}[title={编译链的维护约定}]
\begin{itemize}
\item 前端工具与版本写进仓库（子模块或包清单），产物头部记录的生成器号要能对上。
\item 每次提交都跑 \Tp{spirv\Bk-opt} 前后对比，体积变化超过阈值就人工复核，
防止一次限定符改动把整个块推成动态索引。
\item 反向翻译的产物只进构建目录，永远不回写源目录，避免两套语言互相覆盖。
\item 跨平台的坐标与深度差异用一处矩阵约定解决，并把它写进第~\ref{ch:ch09-vulkan}
章与第~\ref{ch:ch10-d3d12-metal} 章共同引用的检查表。
\end{itemize}
\end{bestpractice}

% ============================================================
% ---- frag_ch08.tex ----
\chapter{计算着色器与 GPU 驱动管线}
\label{ch:ch08-compute-driven}

计算着色器（compute shader）把 \DomGPU{} 从固定的顶点到片元管线里解放出来：
它没有光栅化、没有图元装配，只有一张可读写的大内存和一层线程层级。
一旦摆脱管线阶段的概念束缚，很多原本必须在 \DomCPU{} 上做的簿记工作——
剔除、压缩、排序、构建八叉树、生成间接绘制参数——都可以整段搬进显卡，
让 \DomCPU{} 只负责提交一次派发（dispatch），由 \DomGPU{} 自己决定这一帧到底画多少。
本章讨论这套机制的真实边界：线程层级能给你多少并行、同步要做对哪几件事、
间接绘制如何让实例数量在运行期才确定，以及哪些任务其实并不适合搬下来。

\section{线程层级的三个尺度}

\GSL{} 的计算派发是三维网格：你给出工作组的排布，运行时把每个工作组
（workgroup）切成若干线程（invocation）。线程有三个天然层级，性能和正确性都挂靠在这三层上：

\begin{itemize}
  \item \textbf{线程}（invocation）：最小执行单位，靠 \Tp{gl\_GlobalInvocationID}
        定位自己处理哪条数据。
  \item \textbf{子组}（subgroup）：一组同步推进、能互换了寄存器的线程，
        宽度 \Tp{gl\_{\Bk}Subgroup{\Bk}Size} 由实现决定，常见为 32。子组内用
        \Fn{subgroup{\Bk}Shuffle{\Bk}Up}、\Fn{subgroup{\Bk}Ballot}、\Fn{subgroupAdd} 做无内存的横向归约。
  \item \textbf{工作组}（workgroup）：容量由 \Tp{local\_size} 固定，是能够共享
        \Tp{shared} 内存并以 \Fn{barrier} 同步的最大团块。
\end{itemize}

共享内存（shared memory）是这一层的核心资源，也是第一个硬上限：
每工作组的 \Tp{shared} 容量按字节计，通常只有几十 KB 量级，
而且它和纹理缓存是同一块片上存储的不同用途。想在一个工作组里塞一张
大的查找表或一整个图像块，往往在编译期或派发期就撞墙。因此实践中
把工作组大小压到能同时驻留尽量多工作组，让 \DomGPU{} 有用不完的可调度线程来
掩盖内存延迟，这比把单个工作组做大更重要。

第二个坑是 \textbf{bank 冲突}。共享内存被划分为若干 bank，
同一工作组内的线程若在同一拍访问落在同一 bank 的地址，访问会被串行化，
归约会明显变慢。步长与 bank 数互素、或按 \Tp{shared} 数组维度做交错布局，
是消除冲突的常规手法；具体 bank 数属于设备属性，不能假设，
要按 \DomCPU{} 侧查询到的限制来摆布局。

这三个上限都不能靠记忆写死，必须在初始化时从设备属性读出来：
可分发的三维工作组数 \Tp{max\Bk Compute\Bk Work\Bk Group\Bk Count}、
单工作组每轴线程上限 \Tp{max\Bk Compute\Bk Work\Bk Group\Bk Size}、
单工作组线程总数 \Tp{max\Bk Compute\Bk Work\Bk Group\Bk Invocations}，
再加每工作组可用的共享内存字节 \Tp{max\Bk Compute\Bk Shared\Bk Memory\Bk Size}。
它们共同决定了你能声明多大的 \Tp{local\_size}，也决定了把整批数据
切成多大一块才不会因为驻留不足而排队。

下面的核是一个可见性打包（cull/compact）计算核：它同时用上了 \Tp{shared}
内存、\Fn{barrier} 与全局 \Fn{atomicAdd}，是本章后面几个话题的公共骨架——
它产出的紧凑输出下标，正好喂给后面的间接绘制参数缓冲。

\begin{lstlisting}[style=glslstyle,caption={带共享内存与块内单次原子的打包核},label={lst:ch08-cull-compact}]
#version 450
layout(local_size_x = 64) in;
layout(binding = 0) readonly buffer In  { uint vis[]; } ins;
layout(binding = 1)          buffer Out { uint cnt; uint slot[]; } outb;
shared uint sKeep[64];
shared uint sBase;

void main() {
  uint lid = gl_LocalInvocationID.x;
  uint gid = gl_GlobalInvocationID.x;
  bool keep = (gid < uint(ins.vis.length())) && (ins.vis[gid] == 1u);
  sKeep[lid] = keep ? 1u : 0u;      // shared write
  if (lid == 0u) sBase = 0u;
  barrier();                        // shared writes visible workgroup-wide
  if (lid == 0u) {
    uint sum = 0u;
    for (uint i = 0u; i < gl_WorkGroupSize.x; ++i) sum += sKeep[i];
    sBase = atomicAdd(outb.cnt, sum);   // ONE global atomic per workgroup
  }
  barrier();                        // sBase ready for every thread
  if (keep) {
    uint local = 0u;
    for (uint i = 0u; i < lid; ++i) local += sKeep[i];
    outb.slot[gid] = sBase + local; // compacted destination index
  }
}
\end{lstlisting}

注意这里每个工作组只发一次 \Fn{atomicAdd} 去领取本组的写入起点，
组内的具体位置改由 \Tp{shared} 数组上的线性前缀求得。把全局原子从
每线程一次降到每工作组一次，是这类打包核最常见的提速手段：
原子的成本随争用超线性上升，而线程数往往是工作组长度的百倍。

\section{同步与所有权：barrier 不等于内存布局转换}

计算着色器里的同步有两个正交的方向，混淆它们是数据竞争的根源。

\textbf{方向一：执行与内存的顺序，作用域在工作组内。}
\Fn{barrier} 同时做到两件事——让所有线程在此对齐，并保证 barrier 之前的
内存写对 barrier 之后的读可见。若只需要内存顺序而不需要控制流对齐，
用 \Fn{memoryBarrierShared} 只冲刷 \Tp{shared} 通路的可见性，
或用 \Fn{groupMemoryBarrier} 限定到某类内存。用错作用域的后果通常不崩，
只是在某些设备上偶发读到旧值，因而极难排查。

\textbf{方向二：跨阶段与跨派发的所有权转移。}
上面这个核写完 \Tp{Out} 缓冲后，紧接着的图形或下一次派发要读它，
工作组内的 \Fn{barrier} 对此毫无帮助——它管不到别的派发，更管不到图形阶段。
跨阶段必须在 API 层显式声明依赖：在 Vulkan 里是一次
\Tp{vkCmdPipeline\Bk Barrier}，指定源阶段、目标阶段、访问掩码；
如果这条内存是以图像（image）形式存在，还要顺带完成
\textbf{布局转换}（layout transition），否则采样到的就是未定义内容。

\begin{warning}[title={barrier 与 pipeline barrier 常被混为一谈}]
计算核里的 \Fn{barrier} 只保证同一工作组、同一次派发内的顺序。
把计算产出的数据交给图形阶段读，靠的是 API 侧的 pipeline barrier；
而图像从 \Tp{VK\_IMAGE\Bk LAYOUT\Bk GENERAL} 之类转到可采样布局，靠的是同一条 barrier 里的
layout 字段。漏掉前者会得到竞争读写，漏掉后者会得到画面闪烁或花屏，
两者的现象完全不同，修法也不同。
\end{warning}

一个务实的心智模型：\Tp{shared} 与 \Tp{storage} 缓冲之间没有自动一致性，
\Tp{shared} 只在工作组生命周期内有效，出了工作组就得先落到
\Tp{storage} 或图像，再由 pipeline barrier 交给下游。任何``上一个 pass 算的、
下一个 pass 直接用''的想法，都要拆成写、屏障、读三步。

把这三步固定下来还有个副作用值得利用：既然所有权靠 barrier 显式转移，
你就能在同一批命令里把``计算写、图形读、再计算写回''编排成一条中途不回
\DomCPU{} 的链。这正是 GPU 驱动管线省掉一次 \DomCPU{} 往返的根本原因——
不是计算核本身跑得快，而是这批数据的\textbf{所有权}从头到尾没离开过显存。

\section{间接绘制：让 GPU 决定这一帧画多少}

传统绘制里，实例数量是 \DomCPU{} 在录制命令时就写死的常量。间接绘制
（indirect draw）把这几个参数搬到显存里的一块缓冲，由 \DomGPU{} 在
执行到该绘制命令时才去读，于是剔除核可以边跑边填这份参数表。
在 \GSL{} 侧对应 \Fn{glDrawArraysIndirect} 与 \Fn{glDrawElementsIndirect}；
在 Vulkan 侧对应 \Tp{vkCmdDraw\Bk Indexed\Bk Indirect}，它从一个以
\Tp{VkDraw\Bk Indexed\Bk Indirect\Bk Command} 布局的缓冲读取
索引数、实例数、首索引、顶点偏移与基础实例。

更关键的是``多 draw 间接''：一次提交里可以携带一整段这样的结构数组，
由 \DomGPU{} 逐条消费，配合 count 缓冲还能让\textbf{到底执行几条}也在运行期决定——
Vulkan 的 \Tp{vkCmdDraw\Bk Indexed\Bk Indirect\Bk Count} 会去读一个存放
绘制条数的显存地址，而这个计数正是由前面那个打包核用 \Fn{atomicAdd}
累加出来的。整条链是这样闭合的：

\begin{enumerate}
  \item 剔除核判定每个实例可见与否，写回可见标记；
  \item 打包核把可见实例的紧凑下标与总条数写入缓冲（\Tp{cnt}）；
  \item 间接绘制从同一缓冲读条数，只画活下来的实例；
  \item \DomCPU{} 全程不知道、也不需要知道最终画了多少。
\end{enumerate}

\begin{lstlisting}[style=cppstyle,caption={间接参数缓冲的填充与提交},label={lst:ch08-indirect-submit}]
// layout written by the cull/compact kernel above
struct VkDrawIndexedIndirectCommand {
  uint32_t indexCount, instanceCount,
           firstIndex, vertexOffset, baseInstance;
};
void submit_indirect(VkCommandBuffer cmd, VkBuffer args,
                     VkBuffer count, uint32_t maxDraws) {
  // args: VK_BUFFER_USAGE_INDIRECT_BUFFER_BIT
  // count: same SSBO holds the draw count filled by atomicAdd
  vkCmdDrawIndexedIndirectCount(cmd, args, /*offset*/ 0,
                                count, /*offset*/ 0,
                                maxDraws,
                                sizeof(VkDrawIndexedIndirectCommand));
}
\end{lstlisting}

\DomCPU{} 侧只需把 \Tp{maxDraws} 设成一个上界（总实例数），
真实条数由 \Tp{count} 缓冲给出。上界设得过大不会多画，只会让驱动
预留稍多的调度槽，代价通常可以忽略；设得过小则截断，属于逻辑错误。

\begin{note}[title={count 缓冲的清零时机}]
每帧开始要把这个由 \Fn{atomicAdd} 累加出来的条数复位为零，否则它会跨帧堆积，
把上一帧的实例也算进来。常见做法是用一个单线程派发把它写回零，或在
ping-pong 的两块缓冲间交替，读侧永远指向本帧刚清零的那一块。清零与下游读取
之间要保留 pipeline barrier，别让清零和累加竞争同一地址。
\end{note}

\section{持久化线程与 work queue}

默认派发模型是``一份数据一个线程''：网格大小等于数据量，线程干完一件活就退出。
当任务数量或耗时高度不均时，短任务线程早早结束、长任务线程还在跑，
尾部的空闲浪费随之而来。\textbf{持久化线程}（persistent threads）翻转了这个模型：
只派发恰好能占满 \DomGPU{} 的工作组数量，让每个工作组进一个循环，
每次从全局队列里用 \Fn{atomicAdd} 领一个任务编号，处理完再领下一个。

\begin{compare}[title={一次性派发与持久化线程}]
一次性派发网格大、每线程干一件活，负载均衡靠操作系统式的抢占，
在 \DomGPU{} 上其实没有抢占，只能靠任务本身均匀；优点是简单、无需全局同步。
持久化线程网格固定、任务动态领取，天然做工作窃取（work stealing），
适合任务数远超线程数、且单任务耗时抖动大的场景；代价是需要一个原子的
任务计数器，并且一旦需要``整个网格都完成某步再往下''，就必须靠多个派发
或协作组网格同步来表达，而不能在一个派发里做全局 barrier。
\end{compare}

这里的取舍核心是：\DomGPU{} 上没有跨工作组的全局 barrier（工作组之间默认不可见、
也不保证同时驻留）。持久化模式能成立，靠的是``任务彼此独立、领取顺序无关''。
若算法确需全局同步点（例如多轮扫描），要么把它拆成多次派发、在派发之间用
pipeline barrier 衔接，要么改用子组内归约加原子的分层方案，别指望一次派发里同步所有工作组。

队列的组织方式也影响收益。若任务编号连续分配、且相邻任务的数据也相邻，
持久化线程领到的相邻任务会命中同一批显存页，访存局部性好；反之若任务顺序
被打乱，每次领取都近似随机访问，带宽立刻见底。所以 work queue 不只是
一个原子计数器，它的消费顺序本身也是一个性能参数。

\section{哪些任务真能从计算获益}

计算着色器不是万能加速器，能不能获益取决于并行度和分支发散两件事。
并行度不足的活（数据量小到填不满 \DomGPU{} 的流多处理器），搬到显卡
反而被派发与同步开销拖垮；分支发散的活（同工作组内线程各走各的控制流），
会让整条子组按最坏路径串行执行，实测吞吐接近标量循环。

\begin{longtable}{>{\raggedright\arraybackslash}p{0.20\textwidth}
                 >{\raggedright\arraybackslash}p{0.13\textwidth}
                 >{\raggedright\arraybackslash}p{0.24\textwidth}
                 >{\raggedright\arraybackslash}p{0.22\textwidth}}
\caption{任务族、并行粒度、收益来源与失败模式}\label{tab:ch08-taskmap} \\
\toprule
任务族 & 并行粒度 & 收益来源 & 失败模式 \\ \midrule
\endfirsthead
\toprule
任务族 & 并行粒度 & 收益来源 & 失败模式 \\ \midrule
\endhead
\bottomrule
\endfoot
遮挡与视锥剔除 & 线程 & 剔除结果直接喂间接绘制，省回传 & 包围盒测试分支发散 \\
粒子与模拟更新 & 线程 & 位置速度就地读写，无 \DomCPU{} 往返 & 数据量小填不满核 \\
Mip 与压缩 & 工作组 & 块内归约、\Tp{shared} 存邻域 & 跨行采样 bank 冲突 \\
计数与基数排序 & 子组 & 子组归约加分层原子 & 全局原子每线程一次 \\
图像后处理 & 线程 & 逐像素独立、访存规则 & 核太短派发开销占比高 \\
八叉树构建 & 工作组 & 分层归约与队列 & 需全局同步却单派发 \\
光线实例装配 & 线程 & 产出 \Tp{cnt} 供 count 读 & 上界设小被截断 \\
\end{longtable}

读这张表的正确方式是把它当\textbf{先验}而非结论：同一任务在不同数据规模下
可能落在完全不同的行列。剔了在几万个实例时是计算核的甜点，在几十个实例时
则是纯粹的负担。判断依据永远是``并行度够不够、分支发散不发散、
同步点能不能落在工作组内''这三问。

\DomGPU{} 侧算得再快，也要有 \DomCPU{} 侧的任务图来编排：谁先派发、谁等谁、
哪些能并发提交。成熟做法是在 \DomCPU{} 维护一个显式的依赖图，把每个计算派发
当成图上的节点，节点间的边就是提交序列里的 pipeline barrier 与信号量，
而间接绘制把``这一节点要执行多少次''从图里彻底移到了显存。
换句话说，任务图负责拓扑，计算核负责执行，间接参数负责把数量交给 \DomGPU{}。

\subsection{分轮派发与全局同步的折中}

当一个算法确实需要``所有工作组都算完这一轮再进下一轮''，在一个派发里是做不到的——
工作组之间没有全局 barrier。可行的折中是把它拆成多轮派发，轮与轮之间靠一次
提交里的 pipeline barrier 衔接：每一轮的中间结果落在 \Tp{storage} 缓冲里，
下一轮再读回来。代价是每次跨轮都要付一次内存可见性同步，好处是每轮内部
仍可自由用子组归约与 \Tp{shared}。分几轮、在哪条边界切，是纯经验问题，
只能在目标设备上打点量出来，别凭直觉设定。

\begin{bestpractice}[title={把计算搬下 \DomGPU{} 前的三条自检}]
其一，估并行度：任务实例数是否远大于目标设备的工作组驻留上限，否则不划算。
其二，估分支：同工作组内线程的控制流是否大致一致，发散就考虑按分支预排序。
其三，估同步：需要的同步点能否都收进工作组内；需要全局同步的，先想好
拆成多派发还是分层归约，再动手。三条里任何一条含糊，先用最小核打点测吞吐。
\end{bestpractice}

% ============================================================
% ---- frag_ch09.tex ----
\part{图形 API 与渲染工程}
\chapter{Vulkan：把控制权交给应用}
\label{ch:ch09-vulkan}

Vulkan 与 \Lang{Direct3D 12} 这一代显式（explicit）\DomGPU{} API 共享同一句哲学：
驱动不再替你做决定，代价与自由被一起交到应用手上。\DomCPU{} 侧要自己维护每个对象
的创建与销毁顺序，要自己安排显存的分页与复用，要在每次提交前把依赖关系声明清楚；
换来的回报是开销可预测、命令可多线程并行录制、以及一份能被驱动近乎直接执行的命令流。
\textbf{验证层}（validation layer）是这份契约的看门人：它把原本只在真机上偶发的错误
（用了已释放的内存、漏了一道 barrier、图像布局不对）提前到开发机上，
以 VUID 违规的形式打印出来。理解 Vulkan，本质上就是理解这套显式契约的每一条。

\section{显式 API 的代价与收益}

先说代价。OpenGL 时代那种``绑定一个对象、画一下、忘了它''的写法在这里不成立：
每一个 \Tp{VkBuffer}、\Tp{VkImage}、\Tp{VkPipeline} 都是显式创建、显式销毁的句柄，
销毁顺序必须严格逆着依赖方向——先毁用到它的管线，再毁它引用的布局与描述符，
最后毁设备与实例。顺序错了未必当场崩，但验证层会报出 use-after-destroy。
内存更是要自己从 \Tp{VkDevice\Bk Memory} 里划，并把缓冲或图像绑定上去，
驱动不会替你分配或搬动。

再说收益，而且收益是实打实的。第一是可预测：一帧要提交多少命令、触碰哪些内存、
做哪些同步，全都是应用写死的，不存在驱动``背地里''补做的隐式工作，
帧时间因此平稳。第二是可并行：命令录制可以在多个线程上各自进行，
只要每个线程用独立的命令缓冲；提交仍然串行到一个队列，但录制这一最耗 \DomCPU{}
的阶段被摊到了多核。第三是贴近硬件：现代 \DomGPU{} 的内部结构本就接近这套
命令流加显式同步的模型，显式 API 只是不再让驱动去反推你的意图。

\begin{note}[title={验证层在 instance 阶段就该挂上}]
验证层是 instance 级的组件：在 \Fn{vkCreate\Bk Instance} 时通过
创建信息里的层名列表启用它，越早越好，因为它能捕获几乎所有后续误用。
发布构建里再关掉，代价是运行时校验的开销。别等到出现疑难杂症才想起启用，
那时很多结构性错误（对象生命周期、布局）早已扩散。
\end{note}

\section{对象层级与创建顺序}

Vulkan 的对象是一棵严格的树，创建顺序自上而下、销毁顺序自下而上：

\begin{itemize}
  \item \textbf{实例} \Tp{VkInstance}：整个 Vulkan 使用的根，携带应用信息、
        要启用的实例层与实例扩展、以及要请求的设备扩展清单的初步协商。
  \item \textbf{物理设备} \Tp{VkPhysical\Bk Device}：对一块 \DomGPU{} 的只读描述，
        通过 \Fn{vkEnumerate\Bk Physical\Bk Devices} 得到，用查询函数读出它的
        能力、内存类型与限制；它不是一个可提交的执行体。
  \item \textbf{逻辑设备} \Tp{VkDevice}：对某块物理设备的会话，几乎所有实际
        API 调用都作用在它上面。创建时声明你要用哪些队列族及其优先级。
  \item \textbf{队列} \Tp{VkQueue}：从逻辑设备取得，是提交的入口；
        图形、计算、传输是不同队列族，可以重叠。
  \item \textbf{交换链} \Tp{VkSwapchain\Bk KHR}：一组待呈现的图像，绑定到窗口表面，
        是呈现（present）的载体。
\end{itemize}

创建顺序就是这棵树的深度优先展开：先实例，再枚举物理设备并从中挑选，
再用选中的物理设备建逻辑设备，从逻辑设备取队列，最后用窗口表面与逻辑设备建交换链、
取出交换链图像并为其建 \Tp{Vk\Bk Image\Bk View}。任何一环的创建信息都必须把
\Tp{sType} 填对、把 \Tp{pNext} 链正确串好，否则 \Fn{vkCreate\Bk Device} 会直接
返回失败——Vulkan 大量依赖这种``结构体自描述''的约定，因为它没有 C++ 的虚表可用。

\section{命令录制与提交模型}

工作发生在命令缓冲（command buffer）里：你向 \Fn{vkBegin\Bk Command\Bk Buffer}
申请一条可录制的 \Tp{VkCommand\Bk Buffer}，往里塞 \Fn{vkCmd\Bk Draw\Bk Indexed}、
\Fn{vkCmd\Bk Bind\Bk Pipeline}、\Fn{vkCmd\Bk Pipeline\Bk Barrier} 一类的记录命令，
再 \Fn{vkEnd\Bk Command\Bk Buffer} 收尾。命令缓冲从一个命令池
（command pool，\Tp{VkCommand\Bk Pool}）里分配，池绑定到一个队列族，
并为它的命令缓冲提供后端内存。

这里的关键约束有三条。其一，一条命令缓冲同一时刻只能被一个线程录制，
不能并发写同一条；但不同线程各持一条并行录制是官方支持的，这正是多核收益的来源。
其二，命令池本身不是线程安全的：要么给每个录制线程各建一个池，要么在分配时加锁，
否则两个线程从同一个池 \Fn{vkAllocate\Bk Command\Bk Buffers} 会竞争内部分配器。
其三，重录需要先 \Fn{vkReset\Bk Command\Bk Buffer}，或在建池时带上允许重置的
标志 \Tp{VK\_COMMAND\_POOL\Bk CREATE\_RESET\Bk COMMAND\_BUFFER\Bk BIT}；
用完即弃的每帧命令缓冲通常标成一次性提交
\Tp{VK\_COMMAND\_BUFFER\Bk USAGE\_ONE\_TIME\Bk SUBMIT\_BIT}，让驱动按一次性优化。

\section{同步三件套：fence、semaphore、pipeline barrier}

显式同步是 Vulkan 最容易被误解的部分，因为它把三种完全不同的机制
并排放在你面前，各自解决一个不同的问题。

\textbf{fence（\Tp{VkFence}）解决 ``\DomCPU{} 等 \DomGPU{}''。} 提交时把一个 fence
交给 \Fn{vkQueue\Bk Submit}，\DomGPU{} 执行完这批命令后把 fence 置为 signaled，
\DomCPU{} 用 \Fn{vkWait\Bk For\Bk Fences} 阻塞在它上面，再用 \Fn{vkReset\Bk Fences}
复位以便下一轮。典型用途是``在飞帧数''环形缓冲：重录一条命令缓冲前，先等它上一帧
对应的 fence 落定，避免覆写还在被 \DomGPU{} 消费的命令。

\textbf{semaphore 解决 ``\DomGPU{} 侧的先后''。} 二元信号量（binary semaphore，
\Tp{Vk\Bk Semaphore}）只能在某次提交里 signal 一次、在另一次提交里 wait 一次，
用于跨队列或跨提交的顺序，比如 `` acquire 到交换链图像'' 到 `` 呈现前的等待''。
时间线信号量（timeline semaphore，Vulkan 1.2 起并入核心）对象仍是 \Tp{Vk\Bk Semaphore}，
只是用 \Tp{Vk\Bk Semaphore\Bk Type\Bk Create\Bk Info} 把它建成时间线类型，携带一个
64 位单调值，可反复等待不同阈值，比二元信号量更适合表达 \DomCPU{} 与 \DomGPU{}
都参与的复杂依赖，也能被 \DomCPU{} 直接用 \Fn{vkWait\Bk Semaphores} 等待。

\textbf{pipeline barrier 解决 ``阶段间与内存可见性''。}
\Fn{vkCmd\Bk Pipeline\Bk Barrier} 在一次提交流内部，声明哪些执行阶段必须等哪些
访问（\Tp{Vk\Bk Access\Bk Flag\Bk Bits}）完成，并规定目标阶段；如果附带
\Tp{Vk\Bk Image\Bk Memory\Bk Barrier}，还能顺带完成图像布局转换。
它不是把 \DomCPU{} 或整条流水线卡住，而是在命令流里插入一个 \DomGPU{} 内部的依赖点。

\begin{warning}[title={三者语义别混}]
fence 是 \DomCPU{} 阻塞等 \DomGPU{}；binary 与 timeline semaphore 是 \DomGPU{} 侧
的执行顺序；pipeline barrier 是提交流内部、阶段与访问粒度上的内存与执行依赖。
用 fence 去做本该是 barrier 的事（比如让写缓冲对后续绘制可见）根本不生效，
fence 只在提交边界起作用；反过来拿 semaphore 去表达阶段内的可见性也不对。
三者不是彼此的替代品，把语义搞混是 Vulkan 新手最普遍的一类 bug。
\end{warning}

\section{内存模型：选内存类型、对齐与复用}

创建一个 \Tp{VkBuffer} 或 \Tp{VkImage} 只定义了句柄，它还没有显存。
真正分配是三步：向对象查询 \Tp{Vk\Bk Memory\Bk Requirements}
（大小、对齐、以及一份合法内存类型位掩码），据此从物理设备的内存类型表里挑一个
满足需求的类型，再 \Fn{vkAllocate\Bk Memory} 得到一块 \Tp{VkDevice\Bk Memory}，
最后 \Fn{vkBind\Bk Buffer\Bk Memory} 或 \Fn{vkBind\Bk Image\Bk Memory} 把它绑上。

选择的核心是 \Tp{Vk\Bk Memory\Bk Property\Bk Flags} 里的三种位如何组合。
\Tp{VK\_{\Bk}MEMORY\Bk PROPERTY\Bk DEVICE\Bk LOCAL\Bk BIT} 表示显存在 \DomGPU{} 侧，
快，但 \DomCPU{} 摸不到，适合常驻资源；
\Tp{VK\_{\Bk}MEMORY\Bk PROPERTY\Bk HOST\Bk VISIBLE\Bk BIT} 才允许映射指针读写，
适合上传用的中转（staging）缓冲；
\Tp{VK\_{\Bk}MEMORY\Bk PROPERTY\Bk HOST\Bk COHERENT\Bk BIT} 保证 \DomCPU{} 写入对
\DomGPU{} 立即可见，不必手动 flush/{\Bk}invalidate；非 coherent 时必须自己调
\Fn{vk{\Bk}Flush\Bk Mapped\Bk Memory\Bk Ranges}。
实践中常见配方是：\DomGPU{} 资源用 device-local（必要时配一段 host-visible 中转），
频繁更新的 uniform 用 host-visible 且 coherent。

对齐是第二个隐坑。绑定时的偏移必须是 requirements 里 alignment 的倍数；把多个小对象
打包进一块大 \Tp{VkDevice\Bk Memory} 时还要照顾 \Tp{non\Bk Coherent\Bk Atom\Bk Size}，
让每段映射区对齐，否则 coherent 保证会在区界处失效。正因为手动管理麻烦，
几乎所有引擎都套一层分配器（如 VMA）或复用大块内存做子分配，
再配合 dedicated allocation 扩展在驱动层面告诉它这块内存只服务某一个对象。

\section{描述符、pipeline layout 与两条渲染路径}

着色器绑定的资源通过描述符（descriptor）表达：\Tp{Vk\Bk Descriptor\Bk Set\Bk Layout}
声明一组 binding 的类型与数量，若干 set layout 加上 push constant 范围组装成
\Tp{Vk\Bk Pipeline\Bk Layout}，管线创建时引用这个 layout，绘制时把具体
描述符集合绑上去。描述符集合从 \Tp{Vk\Bk Descriptor\Bk Pool} 分配、
用 \Fn{vkUpdate\Bk Descriptor\Bk Sets} 写入；池一旦重置，从它分配的所有集合立即失效，
这是``用了已释放描述符''这类报错的常见来源。所谓 bindless 的现实形态，是依赖
descriptor indexing 一族的扩展：把数组在着色器里声明成运行时长度，用非统一索引取用，
并允许在绑定之后再更新，从而在一个大数组里指向成千上万的资源——
它确实可用，但设备能力与``部分常驻''语义仍需逐一查询，不能想当然。

渲染目标这条路径有两种写法。老写法是 render pass 与 subpass：先用
\Tp{Vk\Bk Render\Bk Pass} 声明附件、格式与多趟（subpass）之间的依赖，
这些依赖由驱动翻译成 tile 内存上的隐式同步；帧绘制要绑定一个与之匹配的
\Tp{Vk\Bk Framebuffer}。新写法是 dynamic rendering（\Fn{vkCmd\Bk Begin\Bk Rendering}），
Vulkan 1.3 起并入核心：录制时直接给出附件与清除值，不预先创建 render pass 对象，
少一层要维护一致性的隐式状态，对几何缓冲（G-buffer）配置多变、或要与光线查询混排的
延迟渲染更友好。代价是你得自己显式插 \Fn{vkCmd\Bk Pipeline\Bk Barrier} 来表达
原先 subpass 依赖帮你做的顺序。

\begin{lstlisting}[style=cstyle,caption={pipeline layout 组装与图形管线创建},label={lst:ch09-pipeline}]
VkPipelineShaderStageCreateInfo stages[2] = {
  { .sType = VK_STRUCTURE_TYPE_PIPELINE_SHADER_STAGE_CREATE_INFO,
    .stage = VK_SHADER_STAGE_VERTEX_BIT, .module = vs, .pName = "main" },
  { .sType = VK_STRUCTURE_TYPE_PIPELINE_SHADER_STAGE_CREATE_INFO,
    .stage = VK_SHADER_STAGE_FRAGMENT_BIT, .module = fs, .pName = "main" } };

VkPipelineLayoutCreateInfo lci = {
  .sType = VK_STRUCTURE_TYPE_PIPELINE_LAYOUT_CREATE_INFO,
  .setLayoutCount = 1, .pSetLayouts = &set_layout };
VkPipelineLayout pl;
vkCreatePipelineLayout(dev, &lci, NULL, &pl);

VkGraphicsPipelineCreateInfo pi = {
  .sType = VK_STRUCTURE_TYPE_GRAPHICS_PIPELINE_CREATE_INFO,
  .stageCount = 2, .pStages = stages,
  .pVertexInputState = &vin, .pInputAssemblyState = &ia,
  .pViewportState = &view, .pRasterizationState = &rs,
  .pMultisampleState = &ms, .pDepthStencilState = &ds,
  .pColorBlendState = &cb, .layout = pl,
  .renderPass = rp, .subpass = 0 };
VkPipeline pipe;
vkCreateGraphicsPipelines(dev, VK_NULL_HANDLE, 1, &pi, NULL, &pipe);
\end{lstlisting}

一帧的录制、提交与呈现把前述对象串起来：先等本帧槽位的 fence 落定，
重录命令缓冲，向交换链取一张图像并让 acquire 写一个信号量，
提交时以该信号量作等待、以颜色附件输出阶段作等待目标，
呈现时再等绘制完成信号量。下面这段就是这条主线，
其中 \Fn{vkCmd\Bk Begin\Bk Rendering} 走的是 dynamic rendering 路径。

\begin{lstlisting}[style=cstyle,caption={一帧：等 fence、录制、提交、呈现},label={lst:ch09-frame}]
vkWaitForFences(dev, 1, &fence[frame], VK_TRUE, UINT64_MAX);
vkResetFences(dev, 1, &fence[frame]);
vkResetCommandBuffer(cmd[frame], /*flags*/ 0);

VkCommandBufferBeginInfo bi = {
  .sType = VK_STRUCTURE_TYPE_COMMAND_BUFFER_BEGIN_INFO,
  .flags = VK_COMMAND_BUFFER_USAGE_ONE_TIME_SUBMIT_BIT };
vkBeginCommandBuffer(cmd[frame], &bi);
vkCmdBeginRendering(cmd[frame], &rinfo);
vkCmdBindPipeline(cmd[frame], VK_PIPELINE_BIND_POINT_GRAPHICS, pipe);
vkCmdDrawIndexed(cmd[frame], 36, 1, 0, 0, 0);
vkCmdEndRendering(cmd[frame]);
vkEndCommandBuffer(cmd[frame]);

uint32_t image = 0;
vkAcquireNextImageKHR(dev, swap, UINT64_MAX, acq[frame],
                      VK_NULL_HANDLE, &image);
VkPipelineStageFlags wait_stage =
    VK_PIPELINE_STAGE_COLOR_ATTACHMENT_OUTPUT_BIT;
VkSubmitInfo si = { .sType = VK_STRUCTURE_TYPE_SUBMIT_INFO,
  .waitSemaphoreCount = 1, .pWaitSemaphores = &acq[frame],
  .pWaitDstStageMask = &wait_stage, .commandBufferCount = 1,
  .pCommandBuffers = &cmd[frame], .signalSemaphoreCount = 1,
  .pSignalSemaphores = &rdr[frame] };
vkQueueSubmit(queue, 1, &si, fence[frame]);
VkPresentInfoKHR pr = { .sType = VK_STRUCTURE_TYPE_PRESENT_INFO_KHR,
  .waitSemaphoreCount = 1, .pWaitSemaphores = &rdr[frame],
  .swapchainCount = 1, .pSwapchains = &swap, .pImageIndices = &image };
vkQueuePresentKHR(queue, &pr);
frame = (frame + 1) % MAX_FRAMES_IN_FLIGHT;
\end{lstlisting}

\begin{longtable}{>{\raggedright\arraybackslash}p{0.24\textwidth}
                 >{\raggedright\arraybackslash}p{0.13\textwidth}
                 >{\raggedright\arraybackslash}p{0.22\textwidth}
                 >{\raggedright\arraybackslash}p{0.21\textwidth}}
\caption{Vulkan 对象的创建代价、生命周期与并行可用性}\label{tab:ch09-objects} \\
\toprule
对象 & 创建代价 & 生命周期 & 可否并行使用 \\ \midrule
\endfirsthead
\toprule
对象 & 创建代价 & 生命周期 & 可否并行使用 \\ \midrule
\endhead
\bottomrule
\endfoot
\Tp{Vk\Bk Instance} & 进程一次 & 直到退出 & 仅创建期 \\
\Tp{Vk\Bk Physical\Bk Device} & 枚举 & 随实例有效 & 只读、可并发查 \\
\Tp{Vk\Bk Device} & 一次 & 直到退出 & 多数调用需外部同步 \\
\Tp{Vk\Bk Queue} & 随设备 & 随设备 & 每队列提交需串行 \\
\Tp{Vk\Bk Swapchain\Bk KHR} & 可重建 & 尺寸变化即失效 & 取图可并行 \\
\Tp{Vk\Bk Command\Bk Pool} & 每线程一个 & 随线程 & 池本身非线程安全 \\
\Tp{Vk\Bk Command\Bk Buffer} & 从池分配 & 随池 & 各持一条可并录 \\
\Tp{Vk\Bk Device\Bk Memory} & 分配并绑定 & 手动释放复用 & 分配需加锁 \\
\Tp{Vk\Bk Pipeline\Bk Layout} & 每材质一套 & 直到退出 & 建后只读并发 \\
\Tp{Vk\Bk Pipeline} & 重、创建慢 & 直到退出 & 可离线并行创建 \\
\Tp{Vk\Bk Descriptor\Bk Set} & 从池分配 & 池重置即失效 & 写入需外部同步 \\
\end{longtable}

验证层最常见的报错基本能归到几族：hazard（漏了 barrier 导致的竞争读写）、
layout 不符（图像在错误的布局下被采样或写入）、生命周期（用了或毁了仍被引用的对象）、
以及提交约定（重录一条还在飞、或呈现前没等绘制信号量）。把这四族记牢，
配合上面那帧的 fence 与信号量流程，绝大多数 Vulkan 崩溃都能提前定位。

\begin{guideline}[title={对象生命周期由你负全责}]
显式 API 里没有垃圾回收：每个句柄何时毁、按什么序毁、毁前是否还有引用它的
命令在飞，都是应用的事。把这张表和上面那帧流程当基线——先建后用的顺序、
每帧槽位的 fence、跨阶段的 barrier，三者对齐了，验证层才会安静，
\DomGPU{} 才会把这份命令流当成它乐意直接执行的东西。
\end{guideline}

% ============================================================
% ---- frag_ch10.tex ----
\chapter{Direct3D 12 与 Metal}
\label{ch:ch10-d3d12-metal}

\Lang{Direct3D 12} 与 \Lang{Metal} 是另外两套显式 \DomGPU{} API，它们与 Vulkan 站在
同一契约上——\DomCPU{} 负责对象生命周期、内存与同步的显式管理——但把复杂度放在了
不同的地方。把三家放在一起读，才能真正看清``显式 API''这四个字的边界：
哪些概念是共通的、哪些是各家独有的取舍，以及跨平台抽象层要付出什么代价去抹平差异。
本章按这个思路展开：先给 D3D12 的命令模型与描述符绑定，再看它最有名的资源状态转换；
然后跳到 Metal，说它把录制与提交合并到编码器的选择；最后一张概念表收口。

\section{Direct3D 12 的命令模型}

\DomCPU{} 侧的工作从一个 \Tp{ID3D12\Bk Device} 开始，它同时扮演 Vulkan 里
\Tp{Vk\Bk Instance} 与 \Tp{Vk\Bk Device} 的角色——既创建队列与堆，也直接创建资源、
管线与描述符。执行入口是 \Tp{ID3D12\Bk Command\Bk Queue}：它按类型（Direct、
Compute、Copy）分配，接受 \Fn{Execute\Bk Command\Bk Lists} 传入的一段命令列表数组，
并在自己这一条流上串行执行这些列表。

录制发生在 \Tp{ID3D12\Bk Graphics\Bk Command\Bk List} 上，但它只是``录制器''，
它的后端内存来自与之配套的 \Tp{ID3D12\Bk Command\Bk Allocator}。两者都要先
\Fn{Reset} 才能重新使用；分配器与列表是一对一绑定的（Reset 列表时必须传回同一个分配器）。
这条配对决定了并行度：想多线程录制，就每线程各持一份 allocator + command list，
不要共享。提交本身仍在队列上串行。

Bundle（\Tp{D3D12\Bk COMMAND\_LIST\_TYPE\_BUNDLE}）是 D3D12 独有的第三种粒度：
离线录制一段``只能改几何与材质、不能改视口、裁剪盒、渲染目标''的命令序列，
之后在 Direct 列表里用 \Fn{Execute\Bk Command\Bk List} 反复回放。
它的收益不是让 \DomGPU{} 少做什么，而是把 \DomCPU{} 侧的驱动验证与打包摊薄到多次使用上，
很适合每个实例状态变化很少、但要绘制上千次的场景（植被、粒子）。

\Tp{ID3D12\Bk Fence} 是 D3D12 版的 \DomCPU{} 等 \DomGPU{}：提交后调用队列的
\Fn{Signal} 把某个 64 位值写入 fence，\DomCPU{} 侧 \Fn{Get\Bk Completed\Bk Value}
查已完成的值，或 \Fn{Set\Bk Event\Bk On\Bk Completion} 挂到 Windows 事件上阻塞。
这条 fence 只有 \DomCPU{} 视角；\DomGPU{} 内部的顺序靠队列的 \Fn{Wait}
（等待某个 fence 达到某值）与提交批次边界来表达，
跨队列或跨进程的 \DomGPU{} 顺序还需要显式的 event。

把这套 fence 排成环形，就得到一个常见的帧节奏：\DomCPU{} 每帧用递增的 fence 值
提交一次，并记住若干帧前的那个值；重录前对那个旧值 \Fn{Set\Bk Event\Bk On\Bk Completion}，
让线程睡到 \DomGPU{} 追上来，而不是忙等。环的大小（三到四）是在显存占用与
\DomCPU{} 不阻塞之间折中：环太浅，\DomCPU{} 被迫等 \DomGPU{}；环太深，帧延迟被拉长。
这与第~\ref{ch:ch09-vulkan} 章 in-flight fence 的思路一致，只是 D3D12 用单调递增的
64 位计数值取代了 Vulkan 那种 signaled、reset 的二态 fence。

\section{描述符堆、描述符表与 root signature}

D3D12 的绑定模型分成两半。一半叫描述符堆（\Tp{ID3D12\Bk Descriptor\Bk Heap}）：
一段连续的、按类型（CBV-SRV-UAV、Sampler、RTV、DSV）分配的槽位数组，
槽位大小由设备查询 \Fn{Get\Bk Descriptor\Bk Handle\Bk Increment\Bk Size} 给出。
\DomCPU{} 用 \Fn{Create\Bk Shader\Bk Resource\Bk View} 这类调用把资源句柄写进
某个 \DomCPU{} 描述符句柄；如果堆建的时候带了 shader-visible 标志，这些槽位
就可以被 \DomGPU{} 通过设备侧句柄访问（\Fn{Get\Bk GPU\Bk Descriptor\Bk Handle\Bk For\Bk Heap\Bk Start}）。
另一半叫 root signature（\Tp{ID3D12\Bk Root\Bk Signature}），它是一份紧凑的绑定
``合同''：告诉管线有哪些绑定槽、每个槽是常量、描述符表、还是可直接放的
shader-visible 描述符。合同由 \Tp{D3D12\_ROOT\_SIGNATURE\_DESC} 序列化成 blob 后
创建；一旦绑定（\Fn{Set\Bk Graphics\Bk Root\Bk Signature}），后续的
\Fn{Set\Bk Descriptor\Bk Heaps} 告诉命令列表``当前哪些堆是活跃的''，
\Fn{Set\Bk Graphics\Bk Root\Bk Descriptor\Bk Table} 就把堆里的一个连续段接到合同的某个槽上。

对照第~\ref{ch:ch09-vulkan} 章：Vulkan 的 \Tp{Vk\Bk Descriptor\Bk Set} 是``一次组装、
多次绑定''的显式集合，而 D3D12 的表更像一个游标——它不复制描述符，
只是让 root signature 里的某个槽指向当前活跃堆中的一段偏移。
这意味着同一份表内容变化时不需要重建集合，但也要求你在绑定之前
把堆里对应位置的描述符都填好——这是``用了未初始化描述符''这类 bug 的来源。

\begin{note}[title={root signature 与描述符堆的分工}]
root signature 描述的是``接口''（有哪些槽、各是什么类型），描述符堆装的是``数据''
（句柄放在哪段内存）。一个管线只带一个 root signature，换一套绑定合同就得换
signature 并重编管线；堆则可以有很多块、随帧增长，只是同一时刻能被 \DomGPU{}
看到的，是先用 \Fn{Set\Bk Descriptor\Bk Heaps} 登记过的那几块。把两者分开理解，
就不会在``每帧重建 signature''这类昂贵做法上踩坑。
\end{note}

\section{显式资源状态转换：ResourceBarrier}

D3D12 最著名、也最容易出错的一块是资源状态。每个资源在同一时刻处于某个
\Tp{D3D12\Bk RESOURCE\_STATES} 枚举值（\Tp{COPY\_DEST}、\Tp{RENDER\_TARGET}、
\Tp{PIXEL\_SHADER\_RESOURCE}、\Tp{DEPTH\_WRITE}、\Tp{UNORDERED\_ACCESS}……）之一，
\DomGPU{} 的缓存与压缩策略按状态不同而不同，读写必须发生在正确状态下。
命令列表用 \Fn{Resource\Bk Barrier} 显式声明状态转换：一条 barrier 说清楚
哪个资源、哪个子资源、从什么状态到什么状态。

\begin{lstlisting}[style=cppstyle,caption={D3D12 的资源转换与描述符绑定},label={lst:ch10-d3d12-barrier}]
// tex was filled by an upload copy; flip it to SRV for the shader
D3D12_RESOURCE_BARRIER tr = {};
tr.Type        = D3D12_RESOURCE_BARRIER_TYPE_TRANSITION;
tr.Flags       = D3D12_RESOURCE_BARRIER_FLAG_NONE;
tr.Transition.pResource   = tex.Get();
tr.Transition.Subresource = D3D12_RESOURCE_BARRIER_ALL_SUBRESOURCES;
tr.Transition.StateBefore = D3D12_RESOURCE_STATE_COPY_DEST;
tr.Transition.StateAfter  = D3D12_RESOURCE_STATE_PIXEL_SHADER_RESOURCE;
m_cmdList->ResourceBarrier(1, &tr);

m_cmdList->SetGraphicsRootSignature(m_rootSig.Get());
m_cmdList->SetPipelineState(m_pso.Get());
m_cmdList->SetDescriptorHeaps(1, &m_srvHeap);
m_cmdList->SetGraphicsRootDescriptorTable(0,
    m_srvHeap->GetGPUDescriptorHandleForHeapStart());
m_cmdList->DrawIndexedInstanced(36, 1, 0, 0, 0);
\end{lstlisting}

\begin{warning}[title={状态转换为什么成为头号错误源}]
一是\textbf{漏转换}：把刚复制完的图像当 SRV 采样而没转状态，得到的是
\DomGPU{} 私有布局下的随机数据，画面出现规律性噪点；调试器可能只是提示
``resource in state X but used as Y''，也可能什么都不说。
二是\textbf{多转换}：从同一状态转到同一状态（A 到 A）是合法的但无意义，
过多转换会把缓存反复冲刷，性能曲线出现奇怪的抖动。
三是\textbf{跨列表不一致}：fence 只保证执行顺序，不保证状态声明连续——
如果两批列表在同一条队列上把资源从 A 转到 B，中间那批忘了写 B 的
使用状态，验证层不会报，但采样结果已经不对。
四是\textbf{别名与 UAV barrier}：从堆里共址复用两块资源要用 aliasing barrier，
它会让其中一块内容作废；UAV barrier 则用于让一个计算 pass 写完 UAV、
下一个 pass 再读同一 UAV 时可见。两者都不是 transition，
把它们当成 transition 用是隐蔽的 bug。
\end{warning}

还有一处常被误解的``隐式状态''：交换链。DXGI 呈现要求后备缓冲在
\Tp{PRESENT} 状态；驱动在 \Fn{Present} 里会自动帮你转，但\textbf{你必须在录制
时把它当 \Tp{RENDER\_TARGET} 用}，不能在同一帧里对同一个后备缓冲既写又呈现。
多帧 in-flight 环（通常三到四张）与 DX12 交换链的组合，
就是靠 fence 值而不是靠 Vulkan 那种信号量做帧节奏的。

\section{Metal：编码器即开即用}

\Lang{Metal} 的 \Tp{MTL\Bk Device} 大致相当于 \Tp{ID3D12\Bk Device}，
\Tp{MTL\Bk Command\Bk Queue} 相当于命令队列。差别在录制这一层：Metal 没有
``先录好一个命令缓冲、再作为数组提交到队列''这种分离，
而是从 \Tp{MTL\Bk Command\Bk Buffer} 上\textbf{直接开一个编码器}
（\Tp{MTL\Bk Render\Bk Command\Bk Encoder}、\Tp{MTL\Bk Compute\Bk Command\Bk Encoder}
或 blit encoder），编码器立刻把记录命令写进命令缓冲的当前编码段，
你调用 \Fn{endEncoding} 关掉这一段，再对同一命令缓冲 \Fn{commit}，提交就完成了。
一帧里可以开多个 encoder（比如 G-buffer 与后处理各一段），也可以
在同一个 encoder 里画多次。

\begin{lstlisting}[style=cppstyle,caption={Metal 的一帧：开编码器、绘制、呈现、提交},label={lst:ch10-metal-frame}]
id<MTLCommandBuffer> cb = [m_queue commandBuffer];
MTLRenderPassDescriptor *rpd = m_passDesc;
rpd.colorAttachments[0].texture    = m_drawable.texture;
rpd.colorAttachments[0].loadAction = MTLLoadActionClear;
rpd.colorAttachments[0].storeAction = MTLStoreActionStore;

id<MTLRenderCommandEncoder> enc =
    [cb renderCommandEncoderWithDescriptor:rpd];
[enc setRenderPipelineState:m_pss];
[enc setVertexBuffer:m_vbuf offset:0 atIndex:0];
[enc setFragmentTexture:m_tex atIndex:0];
[enc drawIndexedPrimitives:MTLPrimitiveTypeTriangle
                indexCount:36 indexType:MTLUInt16IndexType
                indexBuffer:m_ibuf indexBufferOffset:0];
[enc endEncoding];
[cb presentDrawable:m_drawable];
[cb commit];
\end{lstlisting}

管线状态对象在 Metal 里叫 PSS：由 \Tp{MTL\Bk Render\Bk Pipeline\Bk Descriptor}
组装，一次 \Fn{new\Bk Render\Bk Pipeline\Bk State\Bk With\Bk Descriptor} 得到
不可变的 \Tp{MTL\Bk Render\Bk Pipeline\Bk State}。它比 Vulkan 与 D3D12 的
管线对象更``扁平''：把顶片二元体、混合、深度模板、光栅、附件像素格式
一次性打包，运行时没有 subpass、也没有独立 pipeline layout 这一层。
资源绑定走另一条路：argument buffer——由函数生成的 \Tp{MTL\Bk Argument\Bk Encoder}
把一堆资源引用编成一段 \Tp{MTL\Bk Buffer}，绘制时把这个 buffer 绑到某个索引位；
它更像 Vulkan 的 descriptor set 而不是 D3D12 的 root signature。

Metal 的显式性主要落在\textbf{对象图}与\textbf{驻留}上，而不是内存布局与 barrier 上：
资源对象自带存储模式（private、shared、managed），你不需要显式分配
\Tp{Vk\Bk Device\Bk Memory}，但需要告诉驱动哪些资源在命令缓冲期间要驻留——
或调用 \Fn{use\Bk Resource}、\Fn{use\Bk Heap}，或用把一组资源打包的 residency
set 让命令缓冲知道它们的寿命覆盖到这次提交。跨 \DomCPU{}/{\Bk}\DomGPU{} 的同步靠
\Tp{MTL\Bk Shared\Bk Event}（可跨进程，等价于 Vulkan 的时间线信号量），
\DomGPU{} 内部顺序靠 \Tp{MTL\Bk Event}（由 encoder 的 signal 与 wait
两个调用配对）。

\subsection{编码器家族与 argument buffer 的绑法}
计算与搬运走的是另两个编码器：\Tp{MTL\Bk Compute\Bk Command\Bk Encoder} 与 blit
编码器，用法与渲染编码器同构——从命令缓冲直接开、用完 \Fn{endEncoding}。
argument buffer 的绑法分三步：先从函数取一个 \Tp{MTL\Bk Argument\Bk Encoder}
（\Fn{make\Bk Argument\Bk Encoder\Bk For\Bk Buffer\Bk At\Bk Index}），用它的
\Fn{encoded\Bk length} 建一块够大的 \Tp{MTL\Bk Buffer}，把资源引用编进去；
绘制时把这块 buffer 绑到某个索引位，并对其中每个资源调用 \Fn{use\Bk Resource}
声明本次提交要驻留它。这样一次派发或绘制的成百个绑定就压缩成了一次 buffer 绑定，
这也是 Metal 实现 bindless 的主要手段。

\section{三家的最小公共抽象与复杂度落点}

三家放在一起，可以列出一张最小概念表：设备或会话、命令队列、录制单位、提交单位、
\DomCPU{} 等 \DomGPU{} 的原语、\DomGPU{} 内部顺序的原语、内存与堆对象、
资源绑定的描述、管线状态对象、渲染目标配置、以及跨阶段依赖或状态转换。
下面这张表把这些概念按三家各自的类型对齐。

\begin{longtable}{>{\raggedright\arraybackslash}p{0.14\textwidth}
                 >{\raggedright\arraybackslash}p{0.22\textwidth}
                 >{\raggedright\arraybackslash}p{0.20\textwidth}
                 >{\raggedright\arraybackslash}p{0.22\textwidth}}
\caption{三家显式 API 的概念对照}\label{tab:ch10-concepts} \\
\toprule
概念 & Vulkan & Direct3D 12 & Metal \\ \midrule
\endfirsthead
\toprule
概念 & Vulkan & Direct3D 12 & Metal \\ \midrule
\endhead
\bottomrule
\endfoot
设备与会话 & \Tp{Vk\Bk Device} & \Tp{ID3D12\Bk Device} & \Tp{MTL\Bk Device} \\
命令队列 & \Tp{Vk\Bk Queue} & \Tp{ID3D12\Bk Command\Bk Queue} & \Tp{MTL\Bk Command\Bk Queue} \\
录制单位 & \Tp{Vk\Bk Command\Bk Buffer} & \Tp{ID3D12\Bk Command\Bk List} & \Tp{MTL\Bk Render\Bk Command\Bk Encoder} \\
提交单位 & \Tp{Vk\Bk Submit\Bk Info} & \Fn{Execute\Bk Command\Bk Lists} & \Fn{commit} \\
\DomCPU{} 等 \DomGPU{} & \Tp{Vk\Bk Fence} & \Tp{ID3D12\Bk Fence} & \Tp{MTL\Bk Shared\Bk Event} \\
\DomGPU{} 内部顺序 & \Tp{Vk\Bk Semaphore} & \Fn{Queue\Bk Wait} & \Tp{MTL\Bk Event} \\
内存与堆 & \Tp{Vk\Bk Device\Bk Memory} & \Tp{ID3D12\Bk Heap} & \Tp{MTL\Bk Heap} \\
资源绑定 & \Tp{Vk\Bk Descriptor\Bk Set} & \Tp{ID3D12\Bk Descriptor\Bk Heap} & \Tp{MTL\Bk Argument\Bk Buffer} \\
管线状态 & \Tp{Vk\Bk Pipeline} & \Tp{ID3D12\Bk Pipeline\Bk State} & \Tp{MTL\Bk Render\Bk Pipeline\Bk State} \\
渲染目标配置 & \Tp{Vk\Bk Render\Bk Pass} & \Fn{OMSet\Bk Render\Bk Targets} & \Tp{MTL\Bk Render\Bk Pass\Bk Descriptor} \\
跨阶段依赖 & \Fn{vkCmd\Bk Pipeline\Bk Barrier} & \Fn{Resource\Bk Barrier} & 编码器与驻留 \\
\end{longtable}

三家\textbf{把复杂度放在了不同的地方}。Vulkan 的显式性最彻底：内存类型选择、
对齐、子分配、跨阶段 barrier 与图像布局，全都得你写；换来最细的旋钮、
也换来最长的样板。D3D12 把内存留了一部分给运行时（committed resource 可以直接
创建一个完整资源），但把状态与描述符堆的簿记全交给你：\textbf{资源状态}是它最
独特的负担，而描述符堆与 root signature 之间``合同与游标''的两段式绑定，是它最
需要建立心智模型的地方。Metal 则把内存完全放在资源对象里（storage mode 决定
是否 device-private），把绑定收在 argument buffer 与编码器 setBuffer 里，
真正的复杂度落在\textbf{对象图与驻留}：哪些资源要被这条命令缓冲持有、
什么时候 purge、哪些可以塞进 residency set。

\begin{compare}[title={跨平台抽象层的三条常见策略}]
第一种策略是``以 Vulkan 为原型''：抽象层把三家都表达成 device-local/
host-visible 内存、显式 barrier 与 descriptor set 的样子，D3D12 上的
``状态''被翻译成 barrier 的抽象，Metal 上的驻留被藏在内存分配器之后。
好处是最贴近现代 \DomGPU{}；代价是 D3D12 与 Metal 上的``多做的那一层''必须自己维护。
第二种策略是``最小公分母''：只暴露三家都真正支持的东西，比如只用
device-local 的图像与固定的三种绑定槽，放弃 bindless 之类的高级特性，
简单，但天花板很低。第三种策略是``薄壳 + 后端方言''：抽象层只做对象句柄
与状态机的翻译，每家后端保留自己的最佳实践（比如 D3D12 后端直接用
bundle，Metal 后端直接用 argument buffer）；这是现代引擎最常见的选择，
代价是要维护三套方言，但不会把任何一家拖到别人的抽象下。
\end{compare}

跨平台抽象层最难抹平的两点，是``提交/呈现模型''与``内存与同步''。
提交模型上，Vulkan 的信号量与交换链图像获取、D3D12 的 fence 值与 DX12 交换链的
隐式转换、Metal 的 \Fn{present\Bk Drawable} 语义各自不同；呈现时机、
\DomGPU{} 是否已经把这张图读完、能不能重录这一帧的命令缓冲——三个问题
在三个 API 上答案不一样。内存与同步上，barrier 与状态转换的粒度
不能一对一映射（Vulkan 是阶段乘访问掩码，D3D12 是枚举状态，Metal 是
驻留加编码器分段），抽象层要么把它们当作同一个``资源使用意图''翻译成三家，
要么让上层渲染图明确记录``这一帧用到哪些布局/状态/驻留集''，最后再下沉。

把这三家真正合到一处跑的现实样本是 \Lib{MoltenVK}：它把 Vulkan 的 API 逐调用映射到
Metal，于是 \DomGPU{} 侧的 barrier 退化成 Metal 的编码器分段与驻留声明，
描述符集合则摊平成 argument buffer。它证明``以 Vulkan 为原型''的抽象确实可行，
但也暴露代价：Metal 本没有的东西（多趟 subpass、显式布局）要在翻译层里造出来，
而这些造出来的抽象恰恰是最难跑满性能的地方。设计抽象层时把这条样本当参照，
比凭空发明一套第四语言要稳妥得多。

\begin{guideline}[title={学一套就够，另两套读概念表就够了}]
显式 API 的学习曲线并不真的需要每家重走一遍：把 Vulkan 的 barrier、
D3D12 的状态转换与 Metal 的驻留当成同一件事的三种拼法，
就足以撑起跨平台的引擎。真正需要小心的是那些``看起来像但不等价''的映射——
fence 与 semaphore、root signature 与 descriptor set、bundle 与 secondary
command buffer。这张表和前面几段就是在给这些坑打标。
\end{guideline}

% ============================================================
% ---- frag_ch11.tex ----
\chapter{OpenGL 与 EGL：状态机遗产与现代替代}
\label{ch:ch11-gl-legacy}

OpenGL 的接口形状在 1992 年就定下来了：一台带全局状态的绘图机器，命令按顺序
进入驱动的作业流，由驱动去猜这组状态对应什么硬件配置。这层抽象让上手路径极短，
也留下三个到今天仍未根治的问题：命令之间的依赖是隐式的，状态变更的代价不可见，
而一个上下文在同一时刻只能被一个线程持有。本章讲清楚这三件事在 API 上的具体
表现，以及 \Lib{Vulkan}、\Lib{Direct3D} 12 用哪些显式对象替换了它们。

\section{全局状态机为什么难做并行}

\subsection{隐式依赖}

在 OpenGL 里没有 ``管线对象'' 这种东西：一次 \Fn{glDrawElements} 的语义由几十条
绑定命令的累积结果决定------当前 VAO、活动纹理单元上的采样器、激活的程序对象、
混合方程、深度函数、光栅化状态。驱动必须把这坨状态哈希成一套硬件描述，才能在下
一次绘制时决定要不要重新配置。这个哈希是驱动内部的私有策略，应用看不到它，也
无法为它做预算。显式 API 把这件事变成 \Tp{VkPipeline}：状态组合在进入管线对象时就
被固化，绘制时只剩 ``换一个指针''。

\subsection{上下文即线程归属}

\Fn{glMakeCurrent} 把上下文绑到调用线程，代价是每线程一份完整状态。跨线程共享只有
两条路：把绘制集中在一个线程（于是上传与提交串行化），或者用共享组加外部同步
（见 \S 4 的讨论）。OpenGL 里没有队列、没有命令缓冲、没有
``提交给谁'' 的概念------所有命令都进同一个隐式作业流，因此 ``多线程渲染'' 在
OpenGL 上的正确含义只是 ``把状态变更收进更小临界区''，而不是把 GPU 时间片拆开。

\subsection{校验与 no\_error}

每一次 API 调用，驱动都要校验参数、状态组合与资源是否就绪；校验代码在
\DomCPU{} 侧，且难以并行。有了 \Lib{GL\_KHR\_no\_error}（4.6 把它收进核心），
应用可以要求实现关掉校验，代价是非法参数下的行为未定义、通常直接失去上下文，
从而把这份 CPU 开销换回来。显式 API 把同样的思路做成了默认值：校验只发生在 \Fn{vkValidatePipeline}
这类显式调用与调试图层里。

\begin{guideline}[title={状态机的三条硬约束}]
\begin{itemize}
\item 状态是 ``绑定的''，不是 ``参数的''：任何一次绘制都可能改动全局，
因此并行提交必须先复制或隔离状态。
\item 代价不可见：改变一个采样器绑定可能触发管线重配，而 API 不报错也不提示。
\item 校验与提交同路径：想省掉校验只能整个换掉上下文类型（\Lib{GL\_KHR\_no\_error}），
没有中间档。
\end{itemize}
\end{guideline}

\section{从固定管线到 4.6 core}

固定功能管线（FFP）用一套内置的顶点/纹理/合并阶段替换了着色器：\Fn{glBegin} 到
\Fn{glEnd} 的立即模式、显示列表、\Fn{glMatrixMode} 的模型视图投影栈、
\Fn{glTexEnvi} 的纹理合并方程。2002 年 \Lib{GL\_ARB\_vertex\_program} 与
\Lib{GL\_ARB\_fragment\_program} 把可编程阶段开放出来，\Lib{GL\_ARB\_vertex\_shader}
随后把 \GSL{} 语言接到这两级上，\GSL{} 1.1 随 OpenGL 2.0（2004）进入核心；
此后顶点属性用 \Fn{glVertexAttribPointer} 显式绑定，
\Fn{glClientActiveTexture} 一类遗留调用被逐步弃用。

真正的分水岭是 OpenGL 3.2（2009）引入的 profile 机制：同一个版本号下有
\Tp{GL\_CONTEXT\_CORE\_PROFILE} 与 \Tp{GL\_CONTEXT\_COMPATIBILITY\_PROFILE} 两种
上下文，兼容 profile 保留所有 1.x 时代的 API，core profile 则把立即模式、固定
纹理方程与内置属性删除。4.6（2017）在核心侧补齐了 \SPIV{} 输入
（\Lib{GL\_ARB\_gl\_spirv}）、直接状态访问与
\Lib{GL\_KHR\_parallel\_shader\_compile}。

遗留 API 今天在三个地方持续收租。第一，兼容 profile 要求驱动同时维护两套
光栅化与顶点前端路径，很多新扩展只在 core 上下文里可用。第二，立即模式的顶点
数据要经过一次驱动侧 staging 拷贝，因为 \Fn{glVertex3fv} 的参数指针随时可能失效。
第三，跨平台一致性：OpenGL ES 从未有兼容 profile，桌面与移动的代码路径因此分叉。

\section{离屏渲染、异步上传与屏障}

\Lib{GL\_{\Bk}ARB\_{\Bk}framebuffer\_{\Bk}object}（2005 年的扩展，3.0 入核心）把 ``渲染目标''
变成可命名的对象：颜色附件可以是纹理层或 renderbuffer，深度模板可以是
\Lib{GL\_{\Bk}DEPTH24\_{\Bk}STENCIL8} 这类混合附件。附件规格必须自洽------位深与采样数
不匹配会让 \Fn{gl{\Bk}Check{\Bk}Framebuffer{\Bk}Status} 返回不完成状态，而不是报错。多目标渲染
由 \Fn{gl{\Bk}Draw{\Bk}Buffers} 决定写入哪几个 \Tp{GL\_{\Bk}COLOR\_{\Bk}ATTACHMENT0} 起的附件；
\Lib{GL\_{\Bk}ARB\_{\Bk}draw\_{\Bk}buffers\_{\Bk}blend}（4.2）进一步允许每个附件用不同混合方程。

上传路径上，\Lib{GL\_ARB\_pixel\_buffer\_object}（2.1 入核心）把像素上传
的源从客户端内存改成缓冲对象：绑定 \Tp{GL\_PIXEL\_UNPACK\_BUFFER} 之后，
\Fn{glTexSubImage2D} 的最后一个参数被解释为缓冲内的字节偏移，驱动因此可以
异步 DMA 而不必同步拷贝。配合 \Lib{GL\_ARB\_buffer\_storage}（4.4）的持久映射，
CPU 写入与 GPU 消费可以真正重叠------代价是同步责任完全归应用。

\begin{lstlisting}[style=cppstyle,caption={FBO 离屏目标与双半区持久映射 PBO 上传，含屏障用法},label={lst:ch11-fbo-pbo}]
// 目标：CPU 写纹理数据时不阻塞当前帧的绘制
constexpr GLuint kW = 1024, kH = 1024;
const GLsizeiptr kBytes = (GLsizeiptr)kW * (GLsizeiptr)kH * 4;

GLuint fbo, colorTex, depthRb, pbo;

void InitOffscreen(void)
{
    glCreateFramebuffers(1, &fbo);
    glCreateTextures(GL_TEXTURE_2D, 1, &colorTex);
    glTextureStorage2D(colorTex, 10, GL_RGBA8, kW, kH);   // 10 级 mip
    glNamedFramebufferTexture(fbo, GL_COLOR_ATTACHMENT0, colorTex, 0);

    glCreateRenderbuffers(1, &depthRb);
    glNamedRenderbufferStorage(depthRb, GL_DEPTH24_STENCIL8, kW, kH);
    glNamedFramebufferRenderbuffer(fbo, GL_DEPTH_STENCIL_ATTACHMENT,
                                   GL_RENDERBUFFER, depthRb);

    if (glCheckFramebufferStatus(GL_FRAMEBUFFER) != GL_FRAMEBUFFER_COMPLETE)
        HandleIncomplete();   // 附件规格不自洽时不会自动报错

    glCreateBuffers(1, &pbo);
    // storage 一旦创建不可改尺寸：驱动可直接把这段内存挂进 GPU 页表
    glNamedBufferStorage(pbo, kBytes * 2,
        GL_MAP_WRITE_BIT | GL_MAP_PERSISTENT_BIT | GL_MAP_FLUSH_EXPLICIT_BIT);
}

void UploadNextFrame(const void* src)
{
    static unsigned slot = 0;
    const GLsizeiptr off = (slot++ & 1u) * kBytes;   // 双半区轮转
    char* dst = nullptr;
    glMapNamedBufferRange(pbo, off, kBytes,
        GL_MAP_WRITE_BIT | GL_MAP_PERSISTENT_BIT | GL_MAP_FLUSH_EXPLICIT_BIT,
        (void**)&dst);

    memcpy(dst, src, (size_t)kBytes);
    // 显式 flush：只发布这半区，避免整块缓冲的写回代价
    glFlushMappedNamedBufferRange(pbo, off, kBytes);

    // 映射内存的写，只有 flush + BufferUpdate 屏障后才对后续 GL 命令可见
    glMemoryBarrier(GL_BUFFER_UPDATE_BIT);

    glBindBuffer(GL_PIXEL_UNPACK_BUFFER, pbo);
    glBindTexture(GL_TEXTURE_2D, colorTex);
    glTexSubImage2D(GL_TEXTURE_2D, 0, 0, 0, kW, kH,
                    GL_RGBA, GL_UNSIGNED_BYTE, (const void*)(intptr_t)off);
    glBindBuffer(GL_PIXEL_UNPACK_BUFFER, 0);

    // 若这张纹理随后被 compute 以 image2D 写、被 fragment 采样，
    // 顺序屏障要换成：
    // glMemoryBarrier(GL_SHADER_IMAGE_ACCESS_BIT | GL_TEXTURE_FETCH_BIT);
}
\end{lstlisting}

\subsection{屏障到底同步了什么}

\Fn{glMemoryBarrier}（4.2 起核心，来自 \Lib{GL\_ARB\_shader\_image\_load\_store}）
的掩码是 ``位集合''，每一位代表一类访问。它的语义是：屏障之前的写入，对屏障之后
的、访问类别落在掩码内的操作可见。它不跨越 ``作业流与映射内存'' 之外的边界------
持久映射的正确组合永远是 ``写 -> \Fn{glFlushMappedNamedBufferRange} -> 屏障 ->
使用''，顺序写反就是数据撕裂而不是报错。常用组合：

\begin{itemize}
\item \Tp{GL\_BUFFER\_UPDATE\_BIT}：映射内存写入的缓冲，随后被上传、绘制或
间接绘制命令读取。
\item \Tp{GL\_SHADER\_IMAGE\_ACCESS\_BIT}：image 与 buffer 之间的着色器读写。
\item \Tp{GL\_TEXTURE\_FETCH\_BIT}：纹理内容更新之后被采样。
\item \Tp{GL\_COMMAND\_BITS}：\Lib{GL\_ARB\_indirect\_parameters} 写入的命令流。
\end{itemize}

\Fn{glMemoryBarrier} 是设备侧的，不带等待；要真的把 \DomCPU{} 停在某个 GPU 时间点
上，用的是 \Lib{GL\_ARB\_sync}（3.2）的 \Fn{glFenceSync} 加
\Fn{glClientWaitSync}，或设备侧等待的 \Fn{glWaitSync}。

\section{上下文、共享与表面：WGL、GLX、EGL}

上下文创建从来不是 GL 自己的事：Windows 用 \Lib{WGL}，先 \Fn{wgl{\Bk}Choose{\Bk}Pixel{\Bk}FormatARB}
选像素格式、再 \Fn{wgl{\Bk}Create{\Bk}Context{\Bk}AttribsARB} 定版本与 profile；X11 用
\Lib{GLX}（\Fn{glXChooseFBConfig}、\Fn{glXCreate{\Bk}Context{\Bk}AttribsARB}）；
嵌入式、Android 与 Wayland 用 \Lib{EGL}。EGL 把 ``可绘制表面'' 独立成
\Tp{EGLSurface}，于是同一份 GL 对象可以画进窗口、pbuffer，或者干脆用
\Lib{EGL\_{\Bk}KHR\_{\Bk}surfaceless\_{\Bk}context} 不要表面；Wayland 上窗口表面由合成器持有，
GL 侧要用的是导入进来的 \Tp{EGLImage}。

\subsection{共享组}

创建上下文时传入共享来源（\Fn{eglCreateContext} 的 share\_context 实参、
GLX 的 share\_list、\Fn{wglShareLists}）后，名字对象在组内共享：纹理、缓冲、
renderbuffer、framebuffer、program、pipeline、sampler、query 与 sync 对象。
绑定不共享------每个上下文各有一份完整的绑定状态，因此 ``一个线程一个上下文''
依然是共享组的标准用法。共享组内跨上下文使用同一对象必须显式同步：
\Fn{glFenceSync} 生产、\Fn{glWaitSync} 消费，或 EGL 侧的
\Lib{EGL\_KHR\_wait\_sync}。忘记同步的结果通常是随机花屏，而不是错误码。

\subsection{呈现与帧节拍}

\Fn{eglSwapBuffers} 是队列化的：驱动可以在上一帧未结束时立刻返回，节拍由
\Fn{eglSwapInterval} 决定（1.5 入核心，更早由 \Lib{EGL\_EXT\_swap\_control}
提供）。interval 设 1 即垂直同步，设 0 得到最低延迟但可能撕裂；
\Lib{EGL\_EXT\_swap\_control\_tear} 允许把负值解释成 ``推迟到下一次
合成之后再撕''，从而把撕裂挪到不易见的位置。WGL 与 GLX 的对应物是
\Fn{wglSwapIntervalEXT}、
\Fn{glXSwapIntervalEXT}（后者还能给出目标刷新周期）。EGL 侧减少拷贝的扩展是
\Lib{EGL\_EXT\_buffer\_age}（避免整屏重绘）、\Lib{EGL\_KHR\_swap\_buffers\_with\_damage}
（只呈现脏区）与 \Lib{EGL\_KHR\_gl\_colorspace}（配合
\Lib{EGL\_EXT\_gl\_colorspace\_bt2020\_pq} 之类扩展给出 HDR 传递函数）。

\section{扩展、版本与功能门控}

版本查询看 \Tp{glGetString}(\Tp{GL\_VERSION}) 与
\Tp{GL\_CONTEXT\_MAJOR\_VERSION}，profile 看
\Tp{GL\_CONTEXT\_PROFILE\_MASK}，鲁棒性看 \Tp{GL\_RESET\_ISOLATION\_BEHAVIOR}。
扩展列表在 3.0 之前是 \Tp{glGetString}(\Tp{GL\_EXTENSIONS}) 的一整串字符串（有
长度上限，因此 3.0 之后必须逐条用 \Tp{glGetStringi}(\Tp{GL\_EXTENSIONS}, i) 取，
计数用 \Tp{GL\_NUM\_EXTENSIONS}）。名字后缀说明可信度：\Lib{ARB} 是多家厂商一致
通过、最可能进入核心的路径；\Lib{KHR} 是 Khronos 评审但尚未进核心；\Lib{EXT}
是单家或少数实现，跨厂商行为可能不同；带厂商前缀的（\Lib{NV}、\Lib{AMD}、
\Lib{INTEL}）只在该厂商可用。函数指针必须按扩展逐个加载（\Lib{glad}、
\Lib{GLEW} 或手写 \Fn{wglGetProcAddress}/{\Bk}\Fn{eglGetProcAddress}），且
``版本号够'' 不等于 ``扩展在''：核心特性有版本门槛，扩展没有。

功能门控的正确形状是把每个能力族折成一个运行时标志位，绘制路径按标志位选分支，
并保留一条 ``保底但一定能跑'' 的路径。下面这段 \GSL{} 是 core profile 片元阶段的
最小可用形态，它同时演示了 \Tp{precise} 与顶点侧 \Tp{invariant} 的必要性。

\begin{lstlisting}[style=glslstyle,caption={顶点侧重声明 gl\_Position 为 invariant：两次不同 pass 的同一顶点必须落到同一像素网格},label={lst:ch11-invariant}]
#version 460 core
invariant gl_Position;            // 逐位一致，禁止驱动做等价重写

layout(location = 0) in vec3 aPos;
uniform mat4 uViewProj;

void main()
{
    gl_Position = uViewProj * vec4(aPos, 1.0);
}
\end{lstlisting}

\begin{lstlisting}[style=glslstyle,caption={core profile 片元着色器：precise 阻止 fma 折叠，保证 g-buffer 与光照 pass 的低位一致},label={lst:ch11-core-frag}]
#version 460 core
precision highp float;

layout(location = 0) in vec3 inNormal;
layout(location = 1) in vec3 inWorldPos;
layout(location = 2) in vec2 inUV;

layout(binding = 0) uniform sampler2D uAlbedo;
layout(binding = 1) uniform sampler2D uMetalRough;

layout(location = 0) out vec4 outColor;

// precise 的用途：阻止把 a * b + c 折成一条 fma。
// 折叠后低位不同，延迟管线里前一 pass 写出、后一 pass 读回的
// 打包标量会出现逐帧闪烁。
precise float RoughFromPack(vec2 mr)
{
    float r = mr.y * 0.5 + 0.001;
    return r;
}

void main()
{
    vec3 alb = texture(uAlbedo, inUV).rgb;
    vec2  mr = texture(uMetalRough, inUV).ba;
    vec3  n  = normalize(inNormal);

    // 线性空间出图：色调映射留给后处理，这里不编码 sRGB
    outColor = vec4(alb * n.y * RoughFromPack(mr), mr.x);
}
\end{lstlisting}

\section{什么时候继续用 OpenGL，什么时候必须走显式 API}

表~\ref{tab:ch11-gl-explicit} 把常见能力族按 ``GL 侧最低门槛'' 与 ``显式 API 对应物''
列出，用来估计迁移工作量：左列越接近底部，GL 侧的代码越像在给驱动做提示，而显式
API 侧越要自己做资源生命周期与同步。

\begin{longtable}{>{\raggedright\arraybackslash}p{0.17\textwidth}
                  >{\raggedright\arraybackslash}p{0.34\textwidth}
                  >{\raggedright\arraybackslash}p{0.28\textwidth}}
\caption{OpenGL 能力族、所需版本或扩展与显式 API 对应物}\label{tab:ch11-gl-explicit} \\
\toprule
\textbf{能力族} & \textbf{GL 侧门槛} & \textbf{显式 API 对应物} \\
\midrule
\endfirsthead
\toprule
\textbf{能力族} & \textbf{GL 侧门槛} & \textbf{显式 API 对应物} \\
\midrule
\endhead
\midrule
\endfoot
\bottomrule
\endlastfoot
可编程管线 & \GSL{} 1.20 起；\Lib{GL\_ARB\Bk\_vertex\_shader}、
\Lib{GL\_ARB\Bk\_fragment\_shader} & 图形管线对象一次性固化全部状态 \\
直接状态访问 & 4.5 核心；\Lib{GL\_ARB\Bk\_direct\Bk\_state\Bk\_access} & 无全局绑定态，句柄即对象 \\
离屏渲染 & 3.0；\Lib{GL\_ARB\Bk\_framebuffer\Bk\_object} & 帧缓冲、附件视图与交换链 \\
异步上传 & 2.1；\Lib{GL\_ARB\Bk\_pixel\Bk\_buffer\Bk\_object} & 暂存缓冲加设备侧拷贝命令 \\
持久映射 & 4.4；\Lib{GL\_ARB\Bk\_buffer\Bk\_storage} & 主机可见内存与显式刷新原语 \\
计算着色器 & 4.3；\Lib{GL\_ARB\Bk\_compute\Bk\_shader} & 计算管线与描述符集合 \\
设备与主机同步 & 3.2；\Lib{GL\_ARB\_sync} & fence 与信号量，队列显式化 \\
绑定无关纹理 & 4.5；\Lib{GL\_ARB\Bk\_bindless\Bk\_texture} & 描述符索引与设备地址 \\
\SPIV{} 作为输入 & 4.6；\Lib{GL\_ARB\_gl\_spirv} & 原生二进制着色器 \\
网格着色与间接 & \Lib{GL\_EXT\Bk\_mesh\_shader}、\Lib{GL\_NV\Bk\_mesh\_shader} & 放大与着色两级管线 \\
硬件光追 & \Lib{GL\_NV\Bk\_ray\_tracing}、\Lib{GL\_EXT\Bk\_ray\_tracing} & 加速结构与着色记录表 \\
免校验上下文 & \Lib{GL\_KHR\Bk\_no\_error} & 校验只在调试层，默认关闭 \\
\end{longtable}

一个粗略判断：需要多队列、显式内存管理、跨线程提交、光追或可预测的帧时间时，
显式 API 是唯一选项；需要快速起步、驱动代劳资源生命周期、代码量小、面向教学与
可视化时，core profile 的 OpenGL 仍然完全够用，而且 4.6 已经把不少 ``显式味''
的特性搬进了核心（见第~\ref{ch:ch09-vulkan} 章关于内存与队列的对比，以及
第~\ref{ch:ch10-d3d12-metal} 章关于命令列表的写法）。

\begin{bestpractice}[title={用 OpenGL 但不掉进遗产陷阱}]
\begin{itemize}
\item 只请求 core profile 且带 \Tp{GL\_CONTEXT\_FLAG\Bk\_FORWARD\_COMPATIBLE\_BIT}，
把兼容路径在启动时就
关掉，别让驱动同时维护两套前端。
\item 上传一律走 ``持久映射 PBO + 双半区轮转 + flush + \Fn{glMemoryBarrier}''
四件套；任何一处省略都会退化成同步拷贝，但只表现为帧时间抖动。
\item 能力检测以扩展名为准、以版本号做补充；每条分支都要有一条不依赖新扩展的
保底实现，并在 CI 里按最低配置跑一遍。
\item 需要跨线程或跨上下文共享资源时，把同步对象当成显式契约写进代码，
不要依赖 ``驱动通常会做对''。
\end{itemize}
\end{bestpractice}

% ============================================================
% ---- frag_ch12-window-text.tex ----
\chapter{窗口、事件与字形：渲染上下文的来处与去处}
\label{ch:ch12-window-text}

前面各章都在讲 GPU 拿到的数据从哪来、算完往哪去，但一条隐含的前提线一直没
摊开：GPU 算出的像素最终落进\textbf{谁}的缓冲区，输入事件以什么相位进入这一
帧，以及屏幕上的文字怎么从码点变成三角形。这三件事都发生在图形 API 之外——
窗口系统、事件循环与字体栈。本章把它们讲清楚：交换链参数与呈现调用的细节
属于第~\ref{ch:ch09-vulkan}、第~\ref{ch:ch10-d3d12-metal}、
第~\ref{ch:ch11-gl-legacy} 章，GPU 计时与抓帧工具属于第~\ref{ch:ch14-perf-tools}
章，这里只谈它们跨平台可见的行为与约束。

\section{窗口系统的几种抽象：谁持有像素}
\label{sec:ch12-window-abstraction}

桌面与移动平台的窗口系统表面上都是``建一个窗口、拿一块可绘制区域''，骨子里的
差别在三件事：\textbf{谁持有最终像素}、\textbf{合成器在不在呈现路径上}、以及
\textbf{surface 与 swapchain 的对应关系}。

Win32 的原语是 \Tp{HWND}：\Fn{CreateWindowExW} 注册窗口类并创建句柄，像素直接
进窗口自身。传统 GL 用 \Fn{wglCreateContext} 加 \Fn{wglMakeCurrent} 把上下文
绑到 \Tp{HDC}；Vulkan 用 \Fn{vk{\Bk}CreateWin32{\Bk}SurfaceKHR} 包同一个
\Tp{HWND}；Direct3D 跳过 surface 概念，\Fn{IDXGIFactory2::{\Bk}CreateSwap{\Bk}ChainForHwnd}
直接吃 \Tp{HWND}。DWM 合成器始终在路径上：应用写进交换链后台缓冲，呈现时
移交所有权，屏幕上的最终图像由 DWM 合成。

X11 的原语是 \Tp{xcb\_window\_t}（或遗留 \Tp{Window}）：X 服务器持有窗口树与
像素，客户端经 \Fn{xcb\_create\_window} 请求创建。GL 走 GLX，Vulkan 走
\Fn{vk{\Bk}CreateXlib{\Bk}SurfaceKHR} 或 \Fn{vk{\Bk}CreateXcb{\Bk}SurfaceKHR}。
合成式窗口管理器历史上是可选的：无合成时应用像素直达扫描输出，有合成时与
Wayland 类似多一次纹理拷贝。

Wayland 把``合成器必然在路径上''变成协议本身：客户端只有
\Fn{wl\_compositor\Bk\_create\Bk\_surface} 创建的 \Tp{wl\_surface}，窗口语义
由 \Tp{xdg\_surface}（\Lib{xdg\_shell} 扩展）承载；内容放进 \Tp{wl\_buffer}
（通常是 \Lib{linux-dmabuf} 导入的 DMA-BUF），\Fn{wl\_surface\_commit} 之后
客户端失去所有权，直到 frame callback 告知可复用。Vulkan 对应
\Fn{vk{\Bk}CreateWayland{\Bk}SurfaceKHR}。常见误解：Wayland 没有``交换链直写
窗口像素''，提交节奏由客户端自律，合成器只会对旧内容重复合成来回压。

macOS 上 \Tp{NSWindow} 的 contentView 挂 \Tp{CAMetalLayer}，drawable 就是
呈现单元；Metal 没有 surface 对象，若走 Vulkan 栈则
\Fn{vk{\Bk}CreateMetal{\Bk}SurfaceEXT} 包同一个 \Tp{NSView}。Android 的原语是
\Tp{ANativeWindow}：SurfaceFlinger 必然在路径上，应用经 BufferQueue 取缓冲，
NDK 另提供 \Tp{ASurface{\Bk}Transaction} 把缓冲、时间戳与变更打包成一次原子
提交；呈现统计用 \Fn{ASurfaceTransactionStats\Bk\_getLatchTime} 一类调用回读
排队、锁存与上屏时间。表~\ref{tab:ch12-surface-forms} 按这三问对齐各平台。

\begin{longtable}{@{}>{\raggedright\arraybackslash\footnotesize}p{0.17\textwidth} >{\raggedright\arraybackslash\footnotesize}p{0.21\textwidth} >{\raggedright\arraybackslash\footnotesize}p{0.23\textwidth} >{\raggedright\arraybackslash\footnotesize}p{0.19\textwidth}@{}}
\caption{跨平台 surface 形态：原生句柄、图形 API 入口与合成器位置}\label{tab:ch12-surface-forms} \\
\toprule
\textbf{平台} & \textbf{原生句柄} & \textbf{Surface/呈现入口} & \textbf{合成器在路径上?} \\
\midrule
\endfirsthead
\toprule
\textbf{平台} & \textbf{原生句柄} & \textbf{Surface/呈现入口} & \textbf{合成器在路径上?} \\
\midrule
\endhead
\bottomrule
\endfoot
\endlastfoot
Win32 & \Tp{HWND}、\Tp{HDC} & \Fn{vk{\Bk}CreateWin32{\Bk}SurfaceKHR}、\Fn{CreateSwap{\Bk}ChainForHwnd} & DWM 恒定合成 \\
X11/XCB & \Tp{xcb\_window\_t} & \Fn{vk{\Bk}CreateXcb{\Bk}SurfaceKHR}；GLX & 取决于合成式 WM \\
Wayland & \Tp{wl\_surface} 加 xdg\_shell & \Fn{vk{\Bk}CreateWayland{\Bk}SurfaceKHR} & 合成器是唯一通路 \\
Cocoa & \Tp{NSWindow}、\Tp{CAMetalLayer} & \Fn{vk{\Bk}CreateMetal{\Bk}SurfaceEXT}；drawable & WindowServer 恒定合成 \\
Android & \Tp{ANativeWindow} & BufferQueue、\Tp{ASurface{\Bk}Transaction} & SurfaceFlinger 恒定合成 \\
\end{longtable}

\begin{note}[title={surface 与 swapchain 不是同一层东西}]
Surface 是``向窗口系统声明我要往这块像素区画''的握手对象，除生命周期外没有
队列语义；swapchain 是图形 API 围绕这块区管理的图像集合。一个 surface 上可
以销毁重建多个 swapchain（尺寸或格式变更）。混淆二者会导致 resize 时错误地
重建 surface，把平台侧握手（Wayland 上尤其贵）做进每帧路径。交换链参数表见
第~\ref{ch:ch09-vulkan} 章。
\end{note}

\section{抽象层的取舍：GLFW、SDL、工具集与直接调平台 API}
\label{sec:ch12-abstraction-choice}

\Lib{GLFW} 与 \Lib{SDL} 是最常见的两头：前者刻意薄，后者厚到包含音频、手柄与
线程原语。GLFW 的核心暴露是窗口创建、回调注册、\Fn{glfw{\Bk}Make{\Bk}ContextCurrent}，
以及 \Fn{glfw{\Bk}Create{\Bk}Window{\Bk}Surface}：后者只把 \Tp{HWND}、\Tp{xcb\_{\Bk}window\_{\Bk}t}
或 \Tp{wl\_surface} 包装成 surface，实例扩展名单由
\Fn{glfwGet{\Bk}Required{\Bk}Instance{\Bk}Extensions} 查询，surface 到
swapchain 的那步仍归调用者。SDL 的对应物是 \Fn{SDL\_{\Bk}Vulkan\Bk\_{\Bk}Create{\Bk}Surface}
与 \Fn{SDL\_{\Bk}GL\_{\Bk}Create{\Bk}Context}，事件统一成 \Tp{SDL\_Event} 队列供
\Fn{SDL\_{\Bk}Poll{\Bk}Event} 与 \Fn{SDL\_{\Bk}Wait{\Bk}Event} 取用。两家的共同取舍：不暴露平台
原生句柄之外的深水区细节，需要 Wayland \Tp{wl\_display} 指针或
\Tp{ANative{\Bk}Window} 的回调时机时，要么用取值入口（GLFW 提供
\Fn{glfw{\Bk}Get{\Bk}Win32{\Bk}Window} 一类），要么绕过抽象层。

\Lib{Qt} 或 \Lib{GTK} 的内嵌路线（取 \Fn{QWindow::winId} 或
\Tp{gtk\_drawing\_area} 的原生句柄再建上下文）适合``渲染视图只是应用一小块''
的场景，代价是事件循环归工具集所有，帧循环必须寄生在它的定时器或信号里，
合成器的呈现时机你完全说不上话。Rust 生态的 \Lib{winit} 反证了 C 项目常直接
调平台 API 的原因：一旦需要输入透传子类化、Wayland 子表面、每显示器 DPI
变更通知这类能力，抽象层要么缺入口，要么只给你最小公分母语义。教学代码与
小工具也常手写这一层（自注册窗口类、自拆消息、自己绑上下文），几百行样板
换来的是对时序的完全掌控——这正是上面两类需求出现时的退路。

HiDPI 是抽象层最常糊弄的细节。\Fn{glfwGet{\Bk}WindowSize} 返回\textbf{逻辑}
尺寸（以``点''计，随内容缩放系数变化），\Fn{glfwGet{\Bk}Framebuffer{\Bk}Size}
返回\textbf{像素}尺寸；二者之比即缩放系数，Windows 上常见 1.0、1.25、1.5、
2.0 档，macOS 上常见 1.0 与 2.0 档。所有送进图形 API 的视口、渲染分辨率与
图集页大小必须用像素尺寸；字体的像素目标字号（见~\ref{sec:ch12-glyph-to-triangle}
节）也要按系数换算，否则高分屏上文字要么糊要么小。多显示器允许每屏不同系数，
窗口跨屏移动时 Windows 触发 \Fn{WM\_DPICHANGED} 换档：缩放比变化应当作一次
完整的 resize 走重建路径；GLFW 与 SDL 对该事件的封装程度不一，未封装时需
自行在窗口过程里接住。

\begin{warning}[title={最小化时 frame buffer 尺寸可能为零}]
最小化或被完全遮挡的窗口可能报告 $0\times0$ 的 frame buffer 尺寸。把它直接
送进视口计算与交换链重建会触发除零或非法 extent。所有``尺寸变更''回调入口
先判非零再重建；恢复显示时把``从 0 变非 0''当作一次完整重建事件处理，
包括重做与尺寸无关但可能已丢失的资源。
\end{warning}

\section{事件模型：回调、轮询队列与本帧输入对应哪一帧}
\label{sec:ch12-event-phase}

三种事件模型在``何时把输入交给应用''上分道：\textbf{回调式}（GLFW 的处理器在
平台消息上下文中被直接调用）、\textbf{轮询队列式}（队列保序，X11 应用要等
\Fn{XNextEvent} 或 \Fn{xcb\_poll\_event} 取出 \Tp{ConfigureNotify} 一类事件
才派发，SDL 把各家事件折叠进 \Fn{SDL\_PollEvent}），以及 Wayland 的
\textbf{流式派发}：\Fn{wl\_display\Bk\_dispatch} 从 socket 读一条事件就地调
处理器，而处理器内部可能再进派发（例如 roundtrip 等待回执）——这条路径是
\textbf{重入}的，处理器不能假设自己运行在唯一的顶层循环里。

相位决定输入延迟下限。惯例把 \Fn{glfwPollEvents} 放在帧首：本帧能看到的输入，
是到此刻为止到达的全部事件；放帧尾则本帧输入迟一帧才生效。更讲究的做法是给
输入打单调时钟时间戳，呈现后用平台回执（见~\ref{sec:ch12-present-vsync} 节）
反推``本帧呈现时用户已看见哪条输入''，射击类判定需要这一步。帧队列深度同样
吃延迟：present 之后驱动可以把帧扣在内部队列凑同步点，N 深队列意味着输入到
像素至少多走 N 帧；``低延迟模式''（调小 DXGI 帧延迟上限、Wayland 攒够新帧
尽快 commit）本质都在削这个队列。resize 与 swapchain 重建的竞态最常见：尺寸
回调线程与提交线程不同步，就会用旧 extent 画新窗口，或被呈现拒绝。

\begin{note}[title={焦点、最小化与呈现暂停}]
失去焦点或最小化时继续全速呈现是纯浪费：Wayland 合成器停止发 frame callback
后客户端拿不到驱动信号，Android 上 SurfaceFlinger 不再锁存缓冲，桌面系统
则可能被 DWM 或 WindowServer 降级。稳妥策略：未焦点时转 \Fn{glfwWaitEvents}
（线程睡到下一条事件），最小化时停产出帧，恢复时先重建可能已被丢弃的呈现
资源——各 API 的丢失语义见第~\ref{ch:ch09-vulkan}、第~\ref{ch:ch10-d3d12-metal} 章。
\end{note}

清单~\ref{lst:ch12-event-pump} 给出``帧首清事件、按 framebuffer 尺寸重建''
的骨架，输入时间戳与呈现回执的配对也放在同一处：

\begin{lstlisting}[style=cstyle,label={lst:ch12-event-pump},caption={帧首清事件、按 framebuffer 尺寸重建呈现缓冲的骨架}]
static int g_fb_w = 0, g_fb_h = 0;
static bool g_resized = false;
static uint64_t g_last_input_us, g_vsync_us;

static void on_window_size(GLFWwindow* w, int width, int height) {
    (void)w;
    if (width > 0 && height > 0) {  /* 忽略最小化的零尺寸 */
        g_fb_w = width; g_fb_h = height;
        g_resized = true;
    }
}
void frame_tick(void) {
    glfwPollEvents();               /* 帧首清队列：本帧可见最新输入 */
    g_last_input_us = latest_input_timestamp_us();
    if (g_resized) {
        recreate_present_buffers(g_fb_w, g_fb_h);
        g_resized = false;          /* 重建与提交同线程，避开竞态 */
    }
    if (glfwGetWindowAttrib(g_win, GLFW_ICONIFIED))
        return;                     /* 最小化：不产出帧，等事件唤醒 */
    draw_frame();
    present_and_estimate_vsync_us(&g_vsync_us);
    estimate_input_lag_us(g_last_input_us, g_vsync_us);
}
\end{lstlisting}

\section{呈现与垂直同步：只讲跨平台可见的行为}
\label{sec:ch12-present-vsync}

呈现模式在三大 API 里名字不同、语义相同（参数细节各回第~\ref{ch:ch09-vulkan}、
第~\ref{ch:ch10-d3d12-metal}、第~\ref{ch:ch11-gl-legacy} 章的交换链节，这里不重复
参数表）：\textbf{FIFO} 等下一个垂直消隐间隔交付，永不撕裂也永不掉帧，但
队列深度决定输入延迟；\textbf{mailbox} 多后台缓冲轮转，永不撕裂、不强制等待，
新帧在最近的 vsync 边界上屏；\textbf{立即}不等，画完就换，撕裂出现在换缓冲
时刻的扫描线位置——屏幕上同时有上一帧该线以上与当前帧该线以下的部分，
撕裂线随换缓冲时刻漂移。关掉 vsync 的可见后果就是这条会漂移的撕裂线，外加
帧率上限消失后功耗与帧队列抖动同时上升。被拒的呈现还会改相位：窗口被遮挡
或设备丢失时 present 返回失败码，恢复路径必须与 resize 同构处理。

限速与 pacing 是两件事：vsync 只保证``上屏对齐''，不保证``产出对齐''；产出
抖动（一帧 8 毫秒、下一帧 12 毫秒）叠在 60 赫兹屏上，表现为同一帧被扫描
两次、下一帧只扫一次的节奏跳（jitter），观感比稳定 30 帧更糟。通用削法是
把帧生成率压在刷新率的整数分频上（每 $N$ 个 vsync 只提交一帧）加 CPU 侧
帧预算等待；外部帧限速器与驱动的帧率上限是同一笔账，它们只是把产出抖动
推迟到队列边界，并不消除它。自适应刷新率把``产出对齐''的责任推给显示器，
但撕裂与亮度闪烁的取舍又回到应用手里，跨平台库对它的支持参差不齐。

回执在不同平台的深度差别很大：EGL 有 \Lib{EGL\_{\Bk}KHR\_{\Bk}present\Bk\_{\Bk}id} 的
\Fn{egl{\Bk}Get{\Bk}Present{\Bk}Id} 标识某帧被哪次呈现带走；Android BLAST 路径的 frame
timeline 给排队、锁存、上屏三个时刻，经 \Tp{ASurface{\Bk}TransactionStats}
读取；桌面厂商提供过更深的管线回执（\Lib{NV\_{\Bk}present\_{\Bk}video} 一族与各驱动
实现专有的呈现完成通知），跨平台项目通常只能退到``present 调用返回加最近
一次 vsync 估计''。回执的意义是把``本帧输入对应哪帧呈现''从猜测变成测量，
并形成下一帧预算的闭环——这正是第~\ref{ch:ch14-perf-tools} 章帧预算讨论的
输入端。

\begin{guideline}[title={呈现策略选择的最小决策表}]
\begin{itemize}
\item 交互无关的内容回放：FIFO，队列拉满，求吞吐不求延迟。
\item 一般实时应用：mailbox，或 vsync 开加浅队列；撕裂不可接受就别碰立即模式。
\item 竞技类低延迟：立即模式或厂商低延迟扩展，接受偶发撕裂换输入延迟。
\item 无论哪种：present 返回不等于已显示；不要把 present 调用时间当成
像素到达扫描器的时间。
\end{itemize}
\end{guideline}

\section{字形如何变成三角形：加载、hinting、轮廓化与 shaping}
\label{sec:ch12-glyph-to-triangle}

字体栈的职责链：从字体文件取轮廓、按字号与像素网格整形坐标、把坐标变成
多边形、交给 GPU。\Lib{FreeType} 覆盖前三级里除三角化外的全部。点阵或位图
字体直接给你 \Tp{FT\_Bitmap}（内嵌位图经默认加载即得，彩色字体另有
\Tp{FT\_LOAD\_COLOR}）；轮廓字体走 \Fn{FT\_Set\_Char\_Size}（点单位）或
\Fn{FT\_Set\_Pixel\Bk\_Sizes}（像素单位），\Fn{FT\_Load\_Char} 的
\Tp{FT\_LOAD\Bk\_TARGET\Bk\_LIGHT} 一族（还有 \Tp{NORMAL}、\Tp{HDRY}、
\Tp{VFREE} 三个目标）决定 hinting 把轮廓吸附到像素网格的程度——hinting 是
有损变换，同一字形在不同字号的结果不同。\Fn{FT\_Get\_Glyph} 取到的核心是
\Tp{FT\_Outline}：二次贝塞尔线段（on 点与 off 点控制点对）与直线段，闭合
轮廓带孔洞与方向约定。变三角形需要多边形化：曲线先细分成折线，再做带孔
简单多边形三角化（常用 ear-clipping 类算法），或者把覆盖判定搬进着色器
（见~\ref{sec:ch12-text-to-gpu} 节）。

\textbf{Shaping 必须在 CPU，且不是排版的同义词}。\Lib{HarfBuzz} 把码点序列
变成字形序列：\Fn{hb\_buffer\_add\_utf8} 灌文本，
\Fn{hb\_buffer\Bk\_guess\Bk\_segment\Bk\_properties} 定脚本与语言，
\Fn{hb\_shape} 产出 infos（字形编号加 cluster——每个输出字形映射回输入码点
位置，供光标命中与删改）与 positions（advance 与偏移）两组平行数组。字体
读侧经 \Fn{hb\_font\_set\_funcs} 把 FreeType 回调接进来，hb-ft 层的
\Fn{hb\_ft\_font\_create} 直接代劳，但要记得 \Fn{hb\_ft\_font\Bk\_set\Bk\_load\_flags}
指定 \Tp{FT\_LOAD\_NO\_SCALE}，让缩放归 HarfBuzz 管。为什么必须做：阿拉伯文
连字（同一字符在词首、词中、词尾字形不同）、印度文系音节重组（辅音连缀变
合字或下标形式）、kerning 字距对与 GPOS 表里的上标基线调整——跳过 shaping
逐字渲染，文本会错得不可读。双向文本要先按 Unicode 双向算法分段再分段整形；
哪里能折行、哪里要断词是另一件事，交给 \Lib{hb-linebreak} 或 ICU 的断行
迭代器，别自己按空格拆。

\begin{lstlisting}[style=cstyle,label={lst:ch12-shape-line},caption={HarfBuzz 整形一行文本并逐字形经 FreeType 取轮廓}]
void shape_line(FT_Face face, const char* utf8, uint32_t byte_len) {
    hb_font_t* hb_font = hb_ft_font_create(face, NULL);
    hb_ft_font_set_load_flags(hb_font,
                              FT_LOAD_DEFAULT | FT_LOAD_NO_SCALE);
    hb_buffer_t* buf = hb_buffer_create();
    hb_buffer_add_utf8(buf, utf8, (int)byte_len, 0, (int)byte_len);
    hb_buffer_guess_segment_properties(buf);
    hb_shape(hb_font, buf, NULL, 0);   /* 特性表传空：全用默认 */
    unsigned n = 0;
    const hb_glyph_info_t* infos = hb_buffer_get_glyph_infos(buf, &n);
    const hb_glyph_position_t* pos =
        hb_buffer_get_glyph_positions(buf, &n);
    for (unsigned i = 0; i < n; ++i) {
        /* NO_SCALE 下 codepoint 是字形索引，位移为字体单位 */
        FT_Load_Glyph(face, infos[i].codepoint,
                      FT_LOAD_DEFAULT | FT_LOAD_NO_HINTING);
        emit_outline(&face->glyph->outline,
                     pos[i].x_advance, pos[i].x_offset);
    }
    hb_buffer_destroy(buf);
    hb_font_destroy(hb_font);
}
\end{lstlisting}

清单~\ref{lst:ch12-shape-line} 里 \Tp{hb\_position\_t} 以 26.6 定点数计；
光标定位要用 cluster 回映射：同一 cluster 的多个字形共享一个输入位置，
光标必须落在整簇边界而不是簇中间。

\section{文本进 GPU 的四条路}
\label{sec:ch12-text-to-gpu}

第一条：\textbf{单通道灰度图集}。\Fn{FT\_Render\_Glyph} 以
\Tp{FT\_RENDER\Bk\_MODE\Bk\_NORMAL} 光栅化出 8 位覆盖率位图并打包进大纹理
（字块留 padding、禁 mip），每字形一个四边形，着色器取覆盖率乘色，简单可控。

第二条：\textbf{多通道图集}。LCD 次像素用 \Tp{FT\_RENDER\Bk\_MODE\_LCD}
得每子像素独立灰度的三通道位图，配合垂直方向超采样可解出次像素锐度，但
采样假设物理排列已知，旋转屏与合成器缩放会毁掉它（\Tp{FT\_RENDER\Bk\_MODE\_LCD\Bk\_V}
换轴同理）。位图字体的 \Tp{FT\_PIXEL\Bk\_MODE\Bk\_GRAY2} 等 1、2、4 位打包
模式要手工展开位序与行步长，和直接 memcpy 进纹理的直觉不同，是图集破图第一来源。

第三条：\textbf{（多通道）有向距离场}。轮廓算带符号距离场入图集，着色器以
阈值反解边缘：距离减 $0.5$，\Fn{fwidth} 把过渡带自适应到屏幕密度，经
\Fn{smoothstep} 得 $\alpha$（清单~\ref{lst:ch12-sdf-sample}）；放大不发虚，
描边与阴影只是改阈值与偏移。多通道版由 \Lib{msdf-atlas} 一族工具离线产出，
度量沿用 BMFont 的 \Tp{.fnt} 布局（每字形记 \Tp{x}、\Tp{y}、\Tp{width}、
\Tp{height}、\Tp{xoffset}、\Tp{yoffset}、\Tp{xadvance}）；字块间距不足会泄漏。

第四条：\textbf{轮廓直提交}。\Tp{FT\_Outline} 的点作顶点提交，细分或计算着色器
做带符号覆盖判定；任意缩放无损，代价是顶点内存（第~\ref{ch:ch03-vertex-input}
章）与曲线求交复杂度，计算侧思路见第~\ref{ch:ch08-compute-driven} 章，
对比如表~\ref{tab:ch12-text-gpu-routes}。

\begin{longtable}{@{}>{\raggedright\arraybackslash\footnotesize}p{0.16\textwidth} >{\raggedright\arraybackslash\footnotesize}p{0.22\textwidth} >{\raggedright\arraybackslash\footnotesize}p{0.23\textwidth} >{\raggedright\arraybackslash\footnotesize}p{0.19\textwidth}@{}}
\caption{文本进 GPU 的四条路：CPU 成本、GPU 画质与典型坑}\label{tab:ch12-text-gpu-routes} \\
\toprule
\textbf{路线} & \textbf{CPU 侧成本} & \textbf{GPU 侧画质特征} & \textbf{典型坑} \\
\midrule
\endfirsthead
\toprule
\textbf{路线} & \textbf{CPU 侧成本} & \textbf{GPU 侧画质特征} & \textbf{典型坑} \\
\midrule
\endhead
\bottomrule
\endfoot
\endlastfoot
灰度图集 & 光栅化加打包；缓存失效需重排 & 小字号锐利；放大出像素边 & mip 糊字；hinting 变更使图集错位 \\
多通道或 LCD 图集 & 同上，另加位序展开 & 次像素锐度依赖物理排列 & GRAY2 展开错；旋转缩放破坏次像素 \\
SDF 或 MSDF 图集 & 离线距离场生成与打包 & 放大不发虚；描边阴影近乎免费 & 间距不足致距离泄漏 \\
轮廓直提交 & 曲线细分或多边形化 & 真几何，任意缩放无损 & 顶点内存膨胀；孔洞方向处理错 \\
\end{longtable}

\begin{lstlisting}[style=glslstyle,label={lst:ch12-sdf-sample},caption={图集采样与 SDF 求色的片段着色器}]
#version 450
layout(binding = 0) uniform sampler2D u_atlas;  /* a=覆盖率或距离场 */
layout(location = 0) in vec2 v_uv;
layout(location = 1) in vec4 v_tint;
layout(location = 0) out vec4 o_color;
void main() {
    float d = texture(u_atlas, v_uv).a - 0.5;
    float w = fwidth(d);        /* 灰度图集则直接取 alpha 当覆盖率 */
    float a = smoothstep(-w, w, d);
    o_color = vec4(v_tint.rgb, v_tint.a * a);
}
\end{lstlisting}

\begin{compare}[title={什么时候走哪条路}]
小字号密排的 UI 正文选灰度或 LCD 图集；需旋转、缩放、发光描边的动态文本选
SDF 或 MSDF；字形集巨大且保真优先走轮廓直提交或多页图集换页。
\end{compare}

\begin{bestpractice}[title={窗口、事件与文字三条不可越的线}]
\begin{itemize}
\item 像素尺寸与逻辑尺寸分家：extent、视口与字号的像素目标一律取 framebuffer 尺寸；缩放系数变更当作 resize 走同一条重建路径。
\item 事件在帧首清、输入打时间戳；present 返回不等于已显示，``输入对应哪帧呈现''用平台回执测量，不靠猜。
\item 不要跳过 shaping 逐字取形；cluster 映射留好，断行交给 hb-linebreak 或 ICU；图集要么留足 padding 与采样钳制，要么走 SDF。
\end{itemize}
\end{bestpractice}

% ============================================================
% ---- frag_ch12.tex ----
\chapter{光照与阴影：PBR、IBL 与阴影管线}
\label{ch:ch12-lighting-pbr}

物理化着色（PBR）不是一种着色模型，而是一组约束：BRDF 必须与微表面几何自洽、
能量不能凭空产生、材质参数必须在设备之间可比。这组约束把 ``调参'' 变成了
``测量''，也把代价从美术挪到了运行时------一项采样、一张预过滤图、一级级联
都要按带宽和帧时间去衡量。

\section{微表面 BRDF 的三块：D、G、F}

微表面理论假设可见面由大量朝向不同的小镜面（microfacet）组成，半向量
\(H=(L+V)/\|L+V\|\) 决定哪些小镜面向观察者反射光源。Cook-Torrance 形式把这个
分布写成三项乘积：
\(f = D(H)\,G(L,V,H)\,F(V,H) / \bigl(4\,(N\!\cdot\!L)(N\!\cdot\!V)\bigr)\)。

\begin{itemize}
\item \(D\) 是法线分布项。工业界标准是 GGX（Trowbridge-Reitz）：
\(D = \alpha^{2} / \bigl(\pi ((N\!\cdot\!H)^{2}(\alpha^{2}-1)+1)^{2}\bigr)\)。
它的直觉是 ``长尾''：粗糙度升高时高光的尾部比 Beckmann 展得更开，
因而湿面、磨砂金属这类材质看起来更像实测数据。\(\alpha\) 是
\emph{表观}粗糙度的平方（\(\alpha = r^{2}\)），这一平方来自 Disney 的参数化经验：
线性粗糙度在视觉上不均匀。
\item \(G\) 是几何遮蔽与阴影项，描述小镜面互相挡住的比例。
Schlick-GGX 的单视角形式 \(G_{1}(x)=x/(x(1-k)+k)\) 只要换一个 \(k\) 就同时覆盖
直接光与 IBL：直接光取 \(k=\alpha/2\)，IBL 取同一式子并把 \(N\!\cdot\!V\) 与
视图无关项合并进预积分。Smith 形式是两次 \(G_{1}\) 相乘，代价是两次除法。
\item \(F\) 是菲涅尔项，描述掠射角反射率升高。Schlick 近似
\(F=F_{0}+(1-F_{0})(1-N\!\cdot\!V)^{5}\)；\(F_{0}\) 对介质取折射率决定的
菲涅尔常数 \(\bigl((n-1)/(n+1)\bigr)^{2}\)，默认 \(n\) 为 1.5 时得到 0.04，
对金属直接取 albedo 的线性空间值。工程上还会用到 \(F_{90}=50F_{0}\) 并夹到 1
的近似，用它把 \((1-N\!\cdot\!V)^{5}\) 换成一次乘加。
\end{itemize}

\subsection{metalness/roughness 参数化与能量守恒}

在 metalness/roughness 工作流里，金属度是镜面与漫反射之间的开关：
介质部分（metalness 为 0）有彩色漫反射与无色的 0.04 镜面反射；金属部分
（metalness 为 1）没有漫反射，镜面反射带上 albedo 的颜色。写出来就是
\(F_{0}\) 取 0.04 与 albedo 按金属度插值的结果，
\(k_{d} = (1-F)(1-m)\)，其中 \(m\) 是金属度。

单散射的微表面模型会丢能量：被 \(G\) 挡住的能量既没有反射也没有散射回来。
粗糙度越高、掠射角越大，丢得越多，最直观的症状是 ``金属在粗糙时发黑''。
常见的补法有两类：按能量偏差预计算一张二维补偿图（横轴粗糙度、纵轴
\(N\!\cdot\!V\)），把偏差加回漫反射或镜面项；或者用带归一化因子的
\(G\)（把 Smith 项除以单散射积分）。两者都是近似，但都比不做补偿更接近
实测的反射率曲线，而且几乎不要额外的每像素采样。

\section{直接光照与阴影采样}

一个可直接编译的最小核：GGX 的 D、G、F 与两级 CSM 采样。注意 bias 出现在两个
完全不同的地方------法线偏移在世界空间，深度偏移在光空间的归一化设备坐标里，
把它们混为一谈是阴影 bug 的头号来源。

\begin{lstlisting}[style=glslstyle,caption={阴影采样：法线偏移、深度偏移与 Poisson 盘 PCF},label={lst:ch12-shadow}]
#version 460 core
precision highp float;

uniform sampler2DShadow uShadow;      // 硬件比较 + 双线性滤波
uniform mat4  uLightVP[4];            // 各级联的光空间矩阵
uniform vec4  uCascadeFar;            // 各 split 的远端距离
uniform float uTexel;                 // 1 / 阴影贴图像素尺寸

const vec2 kPoisson[8] = vec2[8](
    vec2(-0.6288, -0.0624), vec2(0.2249,  0.5345),
    vec2(-0.4593,  0.1795), vec2(0.4593, -0.1795),
    vec2(0.2249, -0.5345),  vec2(-0.2249, 0.5345),
    vec2(0.6288,  0.0624),  vec2(-0.2249,-0.5345));

int CascadeOf(float dist)
{
    if (dist < uCascadeFar.x) return 0;
    if (dist < uCascadeFar.y) return 1;
    if (dist < uCascadeFar.z) return 2;
    return 3;
}

float Shadow(vec3 p, vec3 n, vec3 l, float dist)
{
    int c = CascadeOf(dist);
    // 法线偏移：把点推离表面，量随掠射角放大
    float nOff = uTexel / max(dot(n, l), 0.1);
    vec4 ls = uLightVP[c] * vec4(p + n * nOff, 1.0);
    vec3 sc = ls.xyz / ls.w * 0.5 + 0.5;
    if (sc.z >= 1.0) return 1.0;         // 超出远端，按无遮挡处理

    // 深度偏移同样随掠射角变化，随级联变大
    float bias = uTexel * 2.0 * (1.0 + float(c));
    float z = sc.z - bias;
    float sum = 0.0;
    for (int i = 0; i < 8; ++i) {
        vec2 o = kPoisson[i] * uTexel * 2.0;
        sum += texture(uShadow, vec3(sc.xy + o, z));
    }
    return sum * 0.125;
}
\end{lstlisting}

\subsection{Cook-Torrance 的着色核}

\begin{lstlisting}[style=glslstyle,caption={GGX 直接光与 IBL split-sum：D、G、F 的真实写法},label={lst:ch12-cook}]
#version 460 core
precision highp float;

layout(binding = 2) uniform samplerCube uIrradiance;   // 余弦卷积结果
layout(binding = 3) uniform samplerCube uPrefilter;    // GGX 预过滤 mip 链
layout(binding = 4) uniform sampler2D   uBrdfLut;      // 2D 预积分

uniform vec3 uCamPos;
uniform float uMaxLod;       // 预过滤图的最后一级
uniform vec4 uLightPos;      // w = 0 平行光，w = 1 点光
uniform vec3 uLightColor;

vec3 FresnelSchlick(float noV, vec3 f0)
{
    return f0 + (1.0 - f0) * pow(1.0 - noV, 5.0);
}

float D_GGX(float noH, float a)
{
    float a2 = a * a;
    float d = noH * noH * (a2 - 1.0) + 1.0;
    return a2 / (3.14159265 * d * d + 1e-7);
}

float G1(float noX, float k)                 // Schlick-GGX
{
    return noX / (noX * (1.0 - k) + k);
}

float G_SmithGGX(float noV, float noL, float a)
{
    float k = a * 0.5;                       // 直接光的 k
    return G1(noV, k) * G1(noL, k);
}

vec3 Direct(vec3 n, vec3 p, vec3 albedo, float rough, float metal, vec3 f0)
{
    vec3 v = normalize(uCamPos - p);
    vec3 l = normalize(uLightPos.w > 0.5
            ? uLightPos.xyz - p : uLightPos.xyz);
    vec3 h = normalize(v + l);
    float noV = max(dot(n, v), 1e-4), noL = max(dot(n, l), 0.0);
    float noH = max(dot(n, h), 0.0),  voH = max(dot(v, h), 0.0);

    float a = rough * rough;                 // 感知粗糙度到 alpha 的平方
    vec3  f = FresnelSchlick(voH, f0);
    float d = D_GGX(noH, a);
    float g = G_SmithGGX(noV, noL, a);
    vec3  spec = d * g * f / (4.0 * noV * noL + 1e-6);

    vec3 diff = albedo * (1.0 - metal) / 3.14159265;
    diff *= vec3(1.0) - f;                   // 能量守恒的漫反射
    return (diff + spec) * uLightColor * noL;
}

vec3 Indirect(vec3 n, vec3 v, float rough, vec3 f0, vec3 albedo, float metal)
{
    float noV = max(dot(n, v), 1e-4);
    // 环境镜面：split-sum = 预过滤辐射度 x (F0 x LUT.x + LUT.y)
    vec2  lut = texture(uBrdfLut, vec2(noV, rough)).xy;
    vec3  r   = reflect(-v, n);
    float lod = rough * uMaxLod;                    // 粗糙度到层级线性映射
    vec3  pre = textureLod(uPrefilter, r, lod).rgb;
    vec3  k   = FresnelSchlick(noV, f0);
    vec3  kd  = (vec3(1.0) - k) * (vec3(1.0) - metal);
    vec3  irr = texture(uIrradiance, n).rgb;
    return kd * albedo * irr + pre * (f0 * lut.x + lut.y);
}
\end{lstlisting}

\section{IBL 的离线烘焙}

运行时不能对半球积分，所以环境光拆成三张离线图。辐照度图是余弦加权的半球卷积
（等价于对入射辐射求漫反射响应），分辨率很低（每面 32 至 64 足够），也可以用
二阶（三个频带）球谐的 9 个系数替代；预过滤辐射度图按 GGX 的 \(D\) 项做重要性采样，把结果
写进立方体贴图的 mip 链，粗糙度到层级的映射取
\(l = r \times l_{\max}\)（\(r\) 为粗糙度，\(l_{\max}\) 为最后一级）；
BRDF LUT 则在每个
\((N\!\cdot\!V,\ r)\) 像素上把几何项与菲涅尔分解成
\(F_{0}\) 与 \(1-F_{0}\) 两个系数，运行时只做一次二维采样。

\begin{lstlisting}[style=pythonstyle,caption={离线烘焙要点：三张图的生成顺序与采样核（卷积本体在 compute 着色器里跑）},label={lst:ch12-ibl}]
import numpy as np

def ggx_d(noh, rough):
    """GGX 法线分布：alpha 取粗糙度的平方，故这里是 rough ** 4。"""
    a2 = rough ** 4
    deno = noh * noh * (a2 - 1.0) + 1.0
    return a2 / (np.pi * deno * deno + 1e-9)

def ggx_half_vector(xi1, xi2, rough):
    """GGX 重要性采样的半向量：与着色器里的式子逐项对应。"""
    phi = 2.0 * np.pi * xi1
    ct = np.sqrt((1.0 - xi2) / (1.0 + (rough * rough - 1.0) * xi2))
    st = np.sqrt(max(0.0, 1.0 - ct * ct))
    return np.array([st * np.cos(phi), st * np.sin(phi), ct])

def prefilter_face(texel_dir, rough, env_sample):
    """单个 texel、单个粗糙度的预过滤：返回 RGB。
    生产实现把这层循环搬到 compute 着色器，一个 texel 一个 workgroup；
    采样数随层级下降（层级越高波瓣越宽，少采样即可）。"""
    acc, wsum = np.zeros(3), 0.0
    for k in range(4096):
        xi1 = (np.float64(k) + 0.5) * 0.618034 % 1.0
        xi2 = (np.float64(k) + 0.5) * 0.754878 % 1.0
        h = ggx_half_vector(xi1, xi2, rough)
        nh = max(float(np.dot(h, texel_dir)), 0.0)
        if nh <= 0.0:
            continue
        v = 2.0 * nh * h - texel_dir          # 半向量镜像即采样方向
        vh = max(float(np.dot(v, h)), 0.0)
        if vh <= 0.0:
            continue
        w = ggx_d(nh, rough) * nh / (4.0 * vh)   # GGX 立体角权重
        acc += env_sample(v) * w
        wsum += w
    return acc / max(wsum, 1e-9)

# 顺序：全景图转立方面 -> 辐照度（余弦卷积）-> 预过滤（逐 mip）
#       -> BRDF LUT（与材质无关，全局共享一张）-> 打包成 KTX2 进资产库
\end{lstlisting}

\subsection{采样与精度上的两个坑}

预过滤要处理 mip 选择：低粗糙度时 GGX 的波瓣窄，直接采高层 mip 会糊，采
低层 mip 会漏能量，正确做法是按 GGX 波瓣立体角与 texel 立体角的比值估计
level of detail，再配合双三线性采样。辐照度卷积则相反------核很宽，
欠采样表现为色偏而不是噪点，所以用低分辨率加球谐是安全的。两张图都必须在
线性空间里烘焙，输出用浮点或至少 16 位每通道，否则暗部的量化台阶会被曝光放大。

\section{线性空间、色调映射与 bloom}

光照必须在线性辐亮度空间累加，随后才编码到显示传递函数。链路上的三个决定点：
输入纹理按 sRGB 解码（或用 \Lib{GL\_EXT\_texture\_srgb} 与
\Tp{GL\_SRGB8\_ALPHA8} 让硬件做）、累加用浮点 HDR 目标、出图前做一次色调映射。

色调映射族：Reinhard 的 \(c/(1+c)\) 保留高光但会把对比压平；带白点的扩展
形式 \(c(1+c/w^{2})/(1+c)\) 能把最亮处归一到 1；\Lib{ACES} 的工程近似
（Stephen Hill 给的 RRT 加 ODT 拟合多项式）是行业默认，它自带 Rec.709 的
色相偏移，比单纯压亮度更接近胶片；\Lib{AgX} 与 Hable/Uncharted2 属于同一族的
不同拟合。选曲线的实际差别在高亮饱和区的色相保持：物理正确的灯光强度下，
纯 RGB 曲线会把强光推成白色，而 ACES 与 AgX 保留色相。

bloom 的正确位置是 HDR 链里色调映射之前：先做带软膝盖的亮部提取，
再逐级降采样（大核、跨步滤波以避免阶梯），最后逐级升采样叠加；在 LDR 上做
bloom 只会得到一圈灰雾。

\begin{warning}[title={bias 与 peter-panning 是一对不可调和的误差}]
阴影贴图的深度比较有系统性偏差：不加 bias 时自遮挡产生 ``条纹''（acne），
加多了亮面整体脱离阴影（peter-panning，人物的脚与地面断开）。
两者不能靠一个全局常数同时消掉，因为误差随掠射角、距离与级联分辨率变化。
可用的手段按优先级排：法线偏移（沿 \(N\) 推离表面，掠射角自动放大）、
级联相关的深度偏移、斜率缩放偏移（\Fn{glPolygonOffset}，
Vulkan 侧对应 \Tp{depthBiasSlopeFactor} 与 \Tp{depthBiasConstantFactor}）、
以及把偏移放到光空间而不是
NDC 来减轻透视失真。真正稳定的方案是改用与 texel  footprint 相关的 bias，
即在光空间按像素覆盖面积推偏移量。
\end{warning}

\section{级联、PCF、VSM 与接触硬化}

透视失真来自透视投影：近处 texel 覆盖面积小、远处大，于是一张贴图无法同时给出
近处锐利与远处无噪点的阴影。级联阴影贴图（CSM）把视锥按距离切成 2 至 4 段，
每段独立正交包围盒。split 方案用 practical split scheme：纯均匀分法在远处浪费
精度，纯对数分法在近处浪费，两者的 \(\lambda\) 加权混合（\(\lambda\) 常取
0.75 至 0.95）是默认答案。稳定性比精度更难：级联边界每帧抖动会让阴影 ``爬行''，
解法是把包围盒半径与 texel 世界尺寸取整后对齐（texel snapping），并对旋转光源
固定焦距平面。

滤波族按代价排：单次采样（最锐但有锯齿）；盒式 PCF（\(n\times n\) 次采样，
台阶明显）；Poisson 盘 PCF（8 至 16 个各向同性点，需要权重夹以避免 acne，
即 Karis 的做法）；FCP（fractional coverage，4 次硬件比较采样近似一次 PCF，
性价比在主机平台上最好）；各向异性 PCF（沿光线方向拉长核，用 7 次采样换
更宽的软化半径）。接触硬化的物理来源是面光源：半影宽度约等于
\((d_{r}-d_{b})/d_{r}\) 乘光源尺寸（\(d_{r}\) 是到受照点的距离，
\(d_{b}\) 是到第一个遮挡物的距离），
因此需要一次 blocker 搜索（按深度找第一个遮挡者），核半径随深度差增长。
VSM/ESM 存的是深度的指数矩而非比较值，得到单次采样即可软化的阴影，
代价是光泄漏（漏光到不该亮的区域）与数值不稳，通常需要 MSAA 或多级指数来压。

\section{管线族与屏幕空间效果}

表~\ref{tab:ch12-pipelines} 是选型的主表。前向管线把光直接累加进片元，
光源数线性地乘在片元代价上，但天然支持透明与多重采样。 tiled 前向+ 把屏幕切成
16 至 32 像素的 tile，在 compute 阶段按深度球与视锥求交生成每 tile 光源列表，
片元阶段遍历列表；它保留了前向的透明与 MSAA，代价是列表遍历的分支与
每 tile 的共享内存。clustered 把第三维（视深）也切成 16 至 32 段，
于是每个体素的光源列表更短，是当前主机与桌面引擎的主流。

延迟管线先写 g-buffer 再逐光源算光照，把 ``像素 x 光源'' 变成
``有效像素 x 光源''，代价是三件事：g-buffer 带宽、多重采样不可用（需要
depth-only MSAA 或 post-pass 平滑）、以及透明必须回到前向。g-buffer 通道按
128 位每像素排布是通行做法：反照率加金属度打包、八面体编码的法线加粗糙度与
AO、再加一个自发光/高光强度通道；主机上也有压到 64 位每像素的紧凑布局，
靠的是把金属度塞进 alpha、把 AO 与粗糙度打包成 8 位。
光照预通道（light pre-pass）是延迟的变体：只存法线、深度与材质强度，
把可见性算成一张光照图，随后用前向材质重新着色，因而在需要透明与
不同材质管线的场景里更灵活。

\begin{longtable}{>{\raggedright\arraybackslash}p{0.13\textwidth}
                  >{\raggedright\arraybackslash}p{0.19\textwidth}
                  >{\raggedright\arraybackslash}p{0.18\textwidth}
                  >{\raggedright\arraybackslash}p{0.15\textwidth}
                  >{\raggedright\arraybackslash}p{0.15\textwidth}}
\caption{光照管线族在光源代价、带宽、透明与平台上的对比}\label{tab:ch12-pipelines} \\
\toprule
\textbf{管线族} & \textbf{多光源代价} & \textbf{带宽} & \textbf{透明支持} & \textbf{适用平台} \\
\midrule
\endfirsthead
\toprule
\textbf{管线族} & \textbf{多光源代价} & \textbf{带宽} & \textbf{透明支持} & \textbf{适用平台} \\
\midrule
\endhead
\midrule
\endfoot
\bottomrule
\endlastfoot
前向 & 片元乘光源数，几盏灯即饱和 & 最低，单次写 & 天然正确，可 MSAA & 移动与简单场景 \\
前向+（tiled） & 每 tile 一次列表，光数千级无压力 & 低，加列表读 & 正确，保留 MSAA & 移动端主流 \\
clustered & 体素级列表，遍历最短 & 低到中 & 正确 & 主机与桌面 \\
延迟 & 只算可见像素，与几何复杂度解耦 & 高，g-buffer 写读各一轮 & 需回前向或 pre-pass & 桌面主机，动态光多 \\
光照预通道 & 同延迟 & 中，g-buffer 更窄 & 灵活，材质二次着色 & 需多样材质的场景 \\
纯光追 & 与光源数近似无关 & 中，需加速结构驻留 & 天然处理透明与折射 & 见第~\ref{ch:ch13-raytracing} 章 \\
\end{longtable}

屏幕空间效果补的是 ``这一帧看不到但下一帧可能看得到'' 的那部分能量。SSAO 在
深度缓冲里做半球射线步进，估计环境遮挡，它应当只乘到环境项与间接光上，乘到
直接光会立刻显得脏；GTAO 用解析的圆积分替代射线，代价更低且边缘更干净。
SSR 在屏幕空间沿反射方向步进，出屏即失败，因此必须与环境探针或光追兜底混合；
深度重建可用透视正确的方式（按世界步长推进屏幕 y 再反投影）避免 ``无限远
拉伸''。时间维度的降噪与超分不在本章，见第~\ref{ch:ch14-perf-tools} 章的工具面
与第~\ref{ch:ch13-raytracing} 章的时间累积。

\begin{bestpractice}[title={把 PBR 链跑稳的顺序}]
\begin{itemize}
\item 先把线性空间坐实：输入解码、HDR 累加、单次色调映射、最后编码，
任何一处重复编码都会让粗糙度与金属度的观感整体漂移。
\item BRDF 只实现一套，放在共享着色器库里，前向与延迟共用同一函数与同一
能量补偿图，避免两种管线给出两种材质。
\item 阴影 bias 按级联与掠射角参数化，用固定场景（薄板、栅栏、细长物体）做
回归测试；每次调整 CSM split 都重跑，别只看一个角色。
\item IBL 三张图进版本库当资产管，与材质一起打 KTX2，烘焙参数记录在文件头；
换环境探针时先重烘焙再重调灯光强度，顺序反了会把问题归错给灯光。
\end{itemize}
\end{bestpractice}

% ============================================================
% ---- frag_ch13.tex ----
\chapter{光线追踪：加速结构、着色与降噪}
\label{ch:ch13-raytracing}

硬件光线追踪把 ``谁来遮挡'' 这个问题从各种近似里拿出来，交回给几何本身。
代价是三种新的负担：加速结构要占显存并随场景更新，着色程序被拆成一堆小记录
导致分支与局部性成为主瓶颈，而低采样密度带来的噪声必须有降噪器兜住。

\section{BVH 与两级加速结构}

求交本身不贵，贵的是 ``和谁求交''。暴力法每根光线要遍历所有三角形，
硬件实现因此都建 BVH：自顶向下把三角形分进包围盒树，节点划分用表面积启发
（SAH------划分代价与命中概率的乘积最小）。三角形数量翻倍只让遍历深度加一，
这就是 ``对数级'' 的全部含义。

两级结构是为场景的动静分离设计的。底层是几何级的 BLAS（bottom-level）：
一个网格的三角形或 AABB 集合，内部是 BVH；顶层是实例级的 TLAS：每个叶子里
存一个 3 乘 4 变换、24 位实例 ID、8 位可见性掩码与一个指向 BLAS 的地址，
内部再建一棵实例 BVH。分开之后：

\begin{itemize}
\item 同一 BLAS 可以被多个实例复用（植被、道具），显存按网格计而不是按实例计。
\item 实例动、几何不动时只重建 TLAS，代价与实例数相关，与三角形数无关。
\item 几何动（蒙皮、位移、morph 变形）时对 BLAS 做 refit：只更新包围盒与叶重排，
前提是该几何布局在构建时就带上 \emph{允许更新} 标志。
\item 允许更新的构建比不允许的慢、体积更大，且 refit 会逐步劣化树质量，
所以长动画网格通常按若干帧一次的节奏重建而非一直 refit。
\end{itemize}

\section{DXR、Vulkan RT 与 OptiX 的对象形态}

\subsection{DXR}

对象是 \Tp{D3D12\_{\Bk}STATE\_{\Bk}OBJECT}，类型分 collection 与 raytracing pipeline
两种，由一组 subobject 拼装：DXIL 库、着色器配置
（负载与属性大小）、全局与局部根签名、命中组、管线配置（含最大递归深度）。
入口在 \Tp{ID3D12{\Bk}Graphics{\Bk}Command{\Bk}List4} 上：
\Fn{Get{\Bk}Raytracing{\Bk}Acceleration\Bk Structure{\Bk}Prebuild{\Bk}Info} 先报尺寸，
\Fn{Build{\Bk}Raytracing{\Bk}Acceleration{\Bk}Structure} 建，\Fn{Dispatch{\Bk}Rays} 追。
命令列表接口的后续版本沿用同一组入口，只是把屏障与内联追踪补了进来
（Tier 1.1 加内联光线追踪，可在计算队列上走 \Fn{Dispatch{\Bk}Rays} 之外的
一条路径追短光线，并补了 micromap 与 cluster 层）。

\subsection{Vulkan RT}

三个扩展通常一起出现：\Lib{VK\_{\Bk}KHR{\Bk}\_{\Bk}acceleration{\Bk}\_{\Bk}structure}、
\Lib{VK\_{\Bk}KHR{\Bk}\_{\Bk}ray\_{\Bk}tracing\_{\Bk}pipeline} 与 \Lib{VK\_{\Bk}KHR{\Bk}\_{\Bk}ray\_{\Bk}query}，
外加 \Lib{VK\_{\Bk}KHR{\Bk}\_{\Bk}deferred{\Bk}\_{\Bk}host\_{\Bk}operations} 把构建挪到多主机线程。
对象类型是 \Tp{Vk{\Bk}Acceleration{\Bk}StructureKHR}，构建信息类型见下。
构建信息 \Tp{Vk{\Bk}Acceleration{\Bk}Structure{\Bk}Build{\Bk}Geometry{\Bk}InfoKHR} 里
区分类型（顶/底层级）与模式（重建或更新），几何数据以设备地址形式给出
（顶点、索引、可选的变换缓冲）。尺寸先问 \Fn{vk{\Bk}Get{\Bk}Acceleration{\Bk}Structure{\Bk}Build{\Bk}SizesKHR}，
它返回三个值：结果大小、scratch 大小与更新时的序列化大小。

构建标志决定性能形状，它们共享同一个长前缀，形如
\Tp{VK\_{\Bk}BUILD\_{\Bk}ACCELERATION\_{\Bk}STRUCTURE\_{\Bk}...\_{\Bk}BIT\_{\Bk}KHR}：
其中 \Tp{PREFER\_{\Bk}FAST\_{\Bk}TRACE} 与 \Tp{PREFER\_{\Bk}FAST\_{\Bk}BUILD} 是互斥取向，
\Tp{ALLOW\_{\Bk}UPDATE} 允许 refit，\Tp{ALLOW\_{\Bk}COMPACTION} 允许产出压缩版本。
压缩后要取实际大小，用 \Fn{vk{\Bk}Cmd{\Bk}Write{\Bk}Acceleration{\Bk}Structures{\Bk}PropertiesKHR}；
产出目标结构则用 \Fn{vk{\Bk}Cmd{\Bk}Copy{\Bk}Acceleration{\Bk}StructureKHR} 的 compaction 模式。

管线侧的组类型分 general 与 hit group 两类。\Fn{vk{\Bk}Cmd{\Bk}Trace{\Bk}RaysKHR} 收四个
\Tp{Vk{\Bk}Strided{\Bk}Device{\Bk}Address{\Bk}RegionKHR}，
分别对应 raygen、miss、hit、callable。记录表的对齐规则要看属性结构里的两个数，
属性结构叫 \Tp{Vk{\Bk}Physical{\Bk}Device{\Bk}Ray{\Bk}Tracing{\Bk}Pipeline{\Bk}PropertiesKHR}，
两个数是 \Tp{shader{\Bk}Group{\Bk}Handle{\Bk}Size} 与 \Tp{shader{\Bk}Group{\Bk}Base{\Bk}Alignment}：
步长必须是前者的整数倍，起始地址必须按后者对齐------把这两个数写死是最常见的
跨厂商崩溃源。另有一个管线标志值得记住，它把 ``命中组里一定有对应着色器''
变成契约，换来更快的遍历：
\Tp{VK\_{\Bk}PIPELINE\_{\Bk}CREATE{\Bk}RAY\_{\Bk}TRACING{\Bk}NO\_{\Bk}NULL{\Bk}CLOSEST\_{\Bk}HIT{\Bk}SHADERS\_{\Bk}BIT\_{\Bk}KHR}。

\subsection{OptiX}

OptiX 把重点放在程序组织而非命令流：上下文（\Tp{OptixDeviceContext}）下
建模块（从 PTX 或 CUDA 源编译）、程序组（\Tp{OptixProgramGroup}，类型分
raygen、miss、hitgroup、callable，以入口函数名绑定）与管线
（\Tp{OptixPipeline}，含负载大小与递归深度上限）。加速结构叫 GAS（几何层）
与 IAS（实例层），\Fn{optixAccelComputeMemoryUsage} 报输出与临时缓冲尺寸，
\Fn{optixAccelBuild} 用 \Tp{OptixBuildInput}（三角形、自定义 AABB、实例）
构建并返回 \Tp{OptixTraversableHandle}；更新走同一条 API，只把操作改成 update。
它与 DXR、Vulkan 的差别是 ``没有绑定表''：程序组由 handle 索引，
递归与模块链接由管线编译期解析。

\section{从预构建到 DispatchRays}

\begin{lstlisting}[style=cppstyle,caption={DXR 的两级加速结构构建：先问尺寸，再建 BLAS 与 TLAS},label={lst:ch13-blas}]
// ---- 几何描述：一个网格 = 一个 BLAS ----
D3D12_RAYTRACING_GEOMETRY_DESC geo = {};
geo.Type = D3D12_RAYTRACING_GEOMETRY_TYPE_TRIANGLES;
geo.Flags = D3D12_RAYTRACING_GEOMETRY_FLAG_OPAQUE; // 免掉 any-hit 调度
geo.Triangles.IndexFormat = DXGI_FORMAT_R32_UINT;
geo.Triangles.IndexCount = mesh.indexCount;
geo.Triangles.IndexBuffer = ibAddress;
geo.Triangles.VertexFormat = DXGI_FORMAT_R32G32B32_FLOAT;
geo.Triangles.VertexCount = mesh.vertCount;
geo.Triangles.VertexBuffer.StartAddress = vbAddress;
geo.Triangles.VertexBuffer.StrideInBytes = sizeof(MeshVertex);

// ---- 预构建信息：尺寸与 scratch 全由驱动报 ----
D3D12_BUILD_RAYTRACING_ACCELERATION_STRUCTURE_DESC blas = {};
blas.Inputs.Type = D3D12_RAYTRACING_ACCELERATION_STRUCTURE_TYPE_BOTTOM_LEVEL;
blas.Inputs.Flags =
    D3D12_RAYTRACING_ACCELERATION_STRUCTURE_BUILD_FLAG_PREFER_FAST_TRACE |
    D3D12_RAYTRACING_ACCELERATION_STRUCTURE_BUILD_FLAG_ALLOW_UPDATE;
blas.Inputs.NumDescs = 1;
blas.Inputs.DescsLayout = D3D12_ELEMENTS_LAYOUT_ARRAY;
blas.Inputs.pGeometryDescs = &geo;

D3D12_RAYTRACING_ACCELERATION_STRUCTURE_PREBUILD_INFO pre = {};
cmdList4->GetRaytracingAccelerationStructurePrebuildInfo(&blas.Inputs, &pre);
// pre.ResultDataMaxSizeInBytes 是上界；带 ALLOW_UPDATE 时体积明显更大

// ---- 实例描述：只存变换、掩码与 BLAS 地址 ----
D3D12_RAYTRACING_INSTANCE_DESC inst = {};
inst.InstanceID = mesh.slotId;        // 24 bit，命中时进 InstanceIndex()
inst.InstanceMask = 0xFF;             // 8 bit 可见性掩码
memcpy(inst.Transform, worldMatrix, sizeof(inst.Transform));
inst.BLAS = blasCommittedAddress;

blas.DestAccelerationStructureData = blasAddr;
blas.ScratchAccelerationStructureData = scratchAddr;
cmdList4->BuildRaytracingAccelerationStructure(&blas, 0, nullptr);
// 目标资源需转成 RT-ACCELERATION_STRUCTURE 状态后才能被 TLAS 引用

tlas.Inputs.Type = D3D12_RAYTRACING_ACCELERATION_STRUCTURE_TYPE_TOP_LEVEL;
tlas.Inputs.DescsLayout = D3D12_ELEMENTS_LAYOUT_ARRAY;
tlas.Inputs.NumDescs = instanceCount;
tlas.Inputs.InstanceDescs = instanceBufferAddr;
tlas.DestAccelerationStructureData = tlasAddr;
tlas.ScratchAccelerationStructureData = scratchAddr;
cmdList4->BuildRaytracingAccelerationStructure(&tlas, 0, nullptr);
\end{lstlisting}

\begin{lstlisting}[style=cppstyle,caption={记录表布局与批量发射：步长按 16 字节对齐，命中区按材质展开},label={lst:ch13-sbt}]
struct SbtLayout {
    UINT64 raygenAddr, missAddr, hitAddr;
    UINT32 stride, missCount, hitCount;
};
// 上传缓冲布局：[raygen 一条][miss 若干][hit group 每材质一条]
// 步长取 16 的倍数；着色器组句柄从状态对象取出后按步长铺进空位。
UINT32 Align16(UINT32 v) { return (v + 15u) & ~15u; }

void TraceReflections(UINT w, UINT h)
{
    D3D12_DISPATCH_RAYS_DESC d = {};
    d.RaygenRegion.StartAddress   = sbt.raygenAddr;
    d.RaygenRegion.SizeInBytes    = sbt.stride;
    d.RaygenRegion.StrideInBytes  = sbt.stride;
    d.MissRegion.StartAddress     = sbt.missAddr;
    d.MissRegion.SizeInBytes      = sbt.stride * sbt.missCount;
    d.MissRegion.StrideInBytes    = sbt.stride;
    d.HitGroupRegion.StartAddress = sbt.hitAddr;
    d.HitGroupRegion.SizeInBytes  = sbt.stride * sbt.hitCount;
    d.HitGroupRegion.StrideInBytes = sbt.stride;
    d.Width = w; d.Height = h; d.Depth = 1;

    cmdList4->SetPipelineState1(pipelineState.get());
    // 根参数沿用 compute 风格的 root 参数与描述符表设置方式
    cmdList4->DispatchRays(&d);
}
\end{lstlisting}

\subsection{追踪域的组织与预算}

一条命中路径由 ``实例 \(\to\) 几何 \(\to\) 材质'' 三级索引决定用哪条记录，
所以记录表的行数等于命中组数，而不是网格数------把材质而不是网格作为记录单位，
表会小一个数量级。四种程序各司其职：raygen 是唯一能发射线与写结果的入口；
miss 处理落空（通常写环境探针颜色）；closest-hit 是确定命中后的着色，
只跑一次；any-hit 用于遮挡判定与 alpha 测试，但它会在每个候选命中上调用，
是把遍历变成串行的大户。能用射线标志解决的就别留 any-hit 路径：
\Tp{RAY\_FLAG\_FORCE\_OPAQUE} 让实例掩码之外的几何一律视为不透明，
\Tp{RAY\_FLAG\_SKIP\_ANY\_HIT} 直接跳过候选命中的着色器调用。

负载（payload）在追踪时按值复制进每一条射线，大小由着色器配置限定；
递归深度则受管线配置与栈约束双重限制。工程上的做法是把 ``剩余弹跳数''
与 ``方差预算'' 放进负载，在 raygen 里按像素分配采样数：靠近相机、
高光锐利的像素多给一次弹跳，其余直接落到探针。把预算摊平比平均分配更省时间，
因为最坏像素决定整批的耗时。

\subsection{raygen 与 closest-hit 的一对}

\begin{lstlisting}[style=hlslstyle,caption={DXR 的 raygen 与 closest-hit：反射一次弹跳，材质表按实例索引取},label={lst:ch13-hlsl}]
#define PAYLOAD_BYTES 28        // RayPayload 的字节数，须为 4 的倍数

struct RayPayload { float3 color; float travel; uint bounceLeft; };
struct Material   { float4 albedo; float2 params; uint textureSetId; };
struct CBVertex   { float3 pos; float3 nrm; };

RaytracingAccelerationStructure gTLAS;
StructuredBuffer<CBVertex>      gVerts;      // 每三角形 3 个顶点
StructuredBuffer<Material>      gMaterials;  // 按 InstanceIndex 索引
Texture2D<float4>               gGBuffer;    // 光栅阶段写的法线与视向
RWTexture2D<float4>             gOut;
[shader("raygeneration")]
void RayGen()
{
    uint2 px = DispatchRaysIndex().xy;
    float3 n = normalize(gGBuffer.Load(uint3(px, 0)).xyz);
    float3 v = gGBuffer.Load(uint3(px, 1)).xyz;

    RayDesc ray;
    ray.Origin    = ReconstructWorldPos(px) + n * 0.01;  // 抬起原点防自交
    ray.Direction = reflect(v, n);
    ray.TMin = 0.0;
    ray.TMax = 200.0;

    RayPayload p;
    p.color = 0.0; p.travel = 0.0; p.bounceLeft = 1;
    TraceRay(gTLAS, RAY_FLAG_CULL_BACK_FACING_TRIANGLES, 0xFF,
             px.x & 7u, 0.0, PAYLOAD_BYTES, 0, p, ray);

    gOut[px] = float4(p.color, 1.0);
}

[shader("closesthit")]
void ClosestHit(inout RayPayload p,
                BuiltInTriangleIntersectionAttributes attr)
{
    Material m = gMaterials[InstanceIndex()];
    uint   i0 = PrimitiveIndex() * 3u;
    float2 b  = attr.barycentrics;
    float3 n = normalize(
        gVerts[i0].nrm * (1.0 - b.x - b.y) +
        gVerts[i0 + 1u].nrm * b.x + gVerts[i0 + 2u].nrm * b.y);

    // 只算到这里；第二次弹跳交给下一层 raygen pass，避免深递归
    p.color  = ShadeHit(n, -WorldRayDirection(), m, p.bounceLeft);
    p.travel = RayTCurrent();
    p.bounceLeft -= 1u;
}
\end{lstlisting}

\section{分支与局部性为什么成了瓶颈}

光栅化的瓶颈在像素数量与带宽；追踪的瓶颈在 ``一批光线是否走同一条路''。
硬件按 warp 或 wave 批量推进遍历：同一批里如果一半光线进了子树 A、另一半进了
子树 B，两条路径只能串行，实际吞吐按分支数成倍下降。三个后果：

\begin{itemize}
\item 树的遍历是随机访存。节点在显存里按构建顺序而非屏幕顺序排布，
所以光线越发散，缓存命中率越低；向天空的像素与向地面的像素混在同一批里，
是最典型的浪费。
\item 命中组是独立着色程序。一个 wave 覆盖到多种材质时，每种材质的记录
都要跑一遍（其他 lane 空转），因此记录表里材质的排列顺序直接影响耗时。
\item 属性求值与顶点取回是一次额外的随机读。蒙皮、位移与逐顶点法线如果
在 closest-hit 里现算，代价往往高于求交本身。
\end{itemize}

对应的优化都是 ``让同批光线更相似''：按光线类型分 pass（阴影、反射、
半透明各一次 \Fn{DispatchRays}，各自宽度可调）；按屏幕 tile 分组发射；
把 TLAS 实例按材质或屏幕占比排序，使相邻像素倾向同一记录；远处网格退到
更小的 BLAS 或探针；能用 closest-hit 就绝不引入 any-hit；把属性烘焙成
三角形索引可直取的形式；用 micromap 或 cluster 层压几何体积。
判断手段只有实测：拆成按类型的 pass 后分别看耗时与每像素采样数
（第~\ref{ch:ch14-perf-tools} 章给了工具面）。

\begin{note}[title={加速结构是显存里的第二份几何}]
BLAS 与 TLAS 常驻显存，且构建时需要 scratch 缓冲；同时需要保留原始顶点与
索引数据供 refit 使用。典型场景里这三份之和可能与几何本体同量级，
\Tp{PREFER\_FAST\_TRACE} 与 \Tp{ALLOW\_UPDATE} 都会让它更大。
上线前的显存预算要单列这一项，并按实例最多的那一幕计算，而不是默认场景。
\end{note}

\section{降噪：稀疏采样加时空重建}

反射与间接光通常只有每像素一次到四次采样，噪声是设计前提而非事故。
降噪因此是渲染链的一部分，而不是可选后处理。三类做法按信息量递增：

\begin{itemize}
\item 单帧空间滤波：\Lib{Open Image Denoise}（OIDN）以 CPU 为主，输入
颜色加上反照率、法线、深度这类辅助通道，适合离线与对延迟不敏感的交互；
它不需要运动矢量，但也不借用时间信息。
\item 时域累积：把同一像素的历史加权合并，最简形式是
\(c = (c_{p} \times n + s) / (n + 1)\) 的滑动平均（\(c_{p}\) 是历史值，
\(n\) 是已累积帧数），配合最大累积次数上限
来防鬼影。必须同时维护方差或置信度，并用深度与运动矢量做历史拒绝；
邻域裁剪（在亮度色度空间里用前帧的包围盒限制候选值）是性价比最高的一步，
但它会把能量压暗，需要用方差引导的补偿。低粗糙度镜面与细小高光最容易
留下 ``萤火虫''，通常靠钳制单次采样能量与核半径处理。
\item 重采样与引导采样：\Lib{RTXDI} 提供带学习的直接光照采样
（用辐射分布估计代替逐光源遍历，把上千盏灯的问题降成几十个候选），
配合 \Lib{ReSTIR} 这类时空重要性重采样（在像素上存一个带权样本的
``蓄水池''，跨帧与邻域合并），能在极低采样数下给出可用结果。
它们改变的是 ``采样放在哪''，降噪器只是收尾。
\end{itemize}

OptiX 侧的降噪器以 AOV（渲染目标视图）为输入，接受历史帧、光流与
辅助通道，属于第二类；把它与自定义的时域累积叠用时要确认谁负责
反锯齿，否则会出现能量被两次平滑的 ``奶油感''。

\section{混合形态、超分与选型}

实时全路径追踪在绝大多数产品里不是终点形态。常见的混合是：光栅出
G-buffer（深度、法线、材质、运动矢量），追踪只做 ``近似做不了的那几件''------
反射、阴影、环境遮蔽与一次间接漫反射。阴影尤其划算：一根到光源的
射线只需一次 any-hit 判定并提前终止；Vulkan 的 \Lib{VK\_KHR\_ray\_query}
与 OpenGL 的 \Lib{GL\_EXT\_ray\_query} 允许在计算着色器里
直接查询加速结构而不用建完整记录表，实现成本远低于镜面反射链。
这套形态的好处是保留了既有的前向或延迟管线、透明与 MSAA 语义，
只把带宽与显存代价加在真正需要的位置。

超分与降噪解决的是不同问题，但共用同一份历史。时间重建族
（\Lib{DLSS}\Bk/{\Bk}\Bk\Lib{FSR}\Bk/{\Bk}\Bk\Lib{XeSS}）以低分辨率渲染的颜色、
深度、运动矢量为输入重建高分辨率帧，通常要求在线性 HDR、色调映射之前介入，
并且对 alpha 测试类效果（植被、粒子）要求排在它们之后。它们与光追降噪
叠加时的顺序是：追踪采样 \(\to\) 降噪 \(\to\) 时间重建 \(\to\) 色调映射。
顺序颠倒会让重建器把噪声当细节放大。

表~\ref{tab:ch13-choose} 是给排期用的判断表。

\begin{longtable}{>{\raggedright\arraybackslash}p{0.13\textwidth}
                  >{\raggedright\arraybackslash}p{0.17\textwidth}
                  >{\raggedright\arraybackslash}p{0.16\textwidth}
                  >{\raggedright\arraybackslash}p{0.17\textwidth}
                  >{\raggedright\arraybackslash}p{0.16\textwidth}}
\caption{全光追、光栅加光追与纯光栅在质量、显存、帧预算与硬件面\Bk{}上的取舍}\label{tab:ch13-choose} \\
\toprule
\textbf{形态} & \textbf{质量上限} & \textbf{显存} & \textbf{帧预算} & \textbf{硬件面} \\
\midrule
\endfirsthead
\toprule
\textbf{形态} & \textbf{质量上限} & \textbf{显存} & \textbf{帧预算} & \textbf{硬件面} \\
\midrule
\endhead
\midrule
\endfoot
\bottomrule
\endlastfoot
全实时路径追踪 & 最高：多次弹跳、透明与焦散一致 & 加速结构常驻，几何三份数据 & 最紧，反射与 GI 各需降噪 & 需专用 BVH 单元与高显存 \\
光栅加光追反射 & 反射与半透明正确，其余靠探针 & 中，只常驻需要的层级 & 可控，按像素分配采样 & 主流桌面独显 \\
光栅加光追阴影 & 消除 CSM 的漏光与失真 & 低到中 & 每像素一根射线，很省 & 桌面与部分新移动平台 \\
纯光栅加探针 & 静态环境可信，动态物体露馅 & 最低 & 最宽裕 & 全平台，含无光追硬件 \\
离线全追踪 & 与实时形态无关，用于烘焙资产 & 不计 & 不计 & 主机或云端 \\
\end{longtable}

三条不上光追的硬理由：显存已经吃紧（加速结构那份是净增），目标平台里有相当
份额没有可用的硬件 BVH（软件兜底差一到两个数量级，只适合资源极少的小场景），
以及帧预算留给降噪后已不足 2 毫秒。反过来，如果画面的关键缺陷就是
``反射与遮挡不对''（玻璃、水面、金属走廊、动态光源的软阴影），
探针与屏幕空间方案的修补成本会比一次混合形态的改造更高。

\begin{bestpractice}[title={接入光追的顺序}]
\begin{itemize}
\item 先只做光追阴影，用 \Tp{RayQuery} 走计算着色器，不动既有管线；
这一步能验证加速结构的内存、更新节奏与画质收益。
\item 再加反射，并把它当成两次弹跳预算以内的独立 pass；记录表按材质组织，
降噪与超分的接入点一次定死，别留到后期改顺序。
\item 每一层加完都重跑三个固定回归：显存峰值、实例最密集关卡的构建耗时、
以及静止相机下的噪声方差曲线（时域累积是否收敛的直接指标）。
\item 给纯光栅路径留一条能出的分支，并在最低配置上跑通；这不是可选项，
它是发行清单的一部分。
\end{itemize}
\end{bestpractice}

% ============================================================
% ---- frag_ch14.tex ----
\chapter{性能、调试与选型}
\label{ch:ch14-perf-tools}

前三章把光照、材质与光线追踪一路铺到了能出画面。这一章不再加新特效，
而是回到最朴素的问题：帧时间为什么慢、慢在哪一段、以及遇到花屏黑屏时该
拿哪把工具去量。渲染优化没有玄学，只有度量；把 ``测得准'' 做成习惯，选型
判断才有依据。本章依次讲测量路径、瓶颈判据、带宽与显存预算、跨平台的调试
工具族，最后收束到 ``自研还是用库'' 的取舍。

\section{怎么测：CPU 帧时间与 GPU 时间戳}

先看两条最容易读错的时间线。\DomCPU 侧的帧时间好办：在应用主循环里对
\Fn{now} 打点，减去上一帧的打点即可，但它测的是 ``从提交到下一次提交'' 的
墙钟，包含了 VSync 等待与提交气泡，不能等同于 GPU 真正干活的时间。\DomGPU
侧要用硬件时间戳：在命令串起止各写一次 \Tp{vkCmdWriteTimestamp}，读回两个
tick 相减，再乘以设备报告的粒度 \Tp{timestampPeriod}（单位 ns/tick）换算成
毫秒。关键点在于时间戳的绝对值只在同一次 fence 完成的提交内比较才有意义，
跨帧直接相减是错的。

\subsection{CPU 帧时间的采样与分位数}

\DomCPU 帧时间要按 ``一帧一帧'' 采，不要取平均。平均值会把偶发的尖峰抹平，
而卡顿恰恰来自尖峰。做法是把每帧的墙钟差写进环形缓冲，导出成 \Tp{.csv}，
再统计中位数与 p95、p99：中位数反映常态，p95/p99 反映掉帧的概率。测量时
还要区分 ``开启 VSync'' 与 ``关闭 VSync'' 两种模式：前者帧时间被钳在刷新
周期，慢的表现是丢帧计数上升而不是单帧变长；后者才能看到真实的每帧耗时。

\subsection{GPU 时间戳的读回路径}

\DomGPU 时间戳的读回有两条路径：同步路径用
\Tp{VK\_QUERY\_RESULT\_WAIT\_BIT}，驱动自旋到查询就绪再返回，简单但会制造
气泡；异步路径用 \Tp{VK\_QUERY\_RESULT\_WITH\_INDEX\_BIT} 加轮询
\Fn{vkGetQueryPoolResults}，返回 \Tp{VK\_NOT\_READY} 就下一帧再读，把等待
摊到流水里。下面的例子走同步路径，因为它最适合建立 ``可信的第一版基线''；
上线前的连续采样应换成异步路径。

\begin{lstlisting}[style=cppstyle,caption={Vulkan GPU 时间戳的写入与读回（真实调用序列）},label={lst:ch14-timestamp}]
// 1) 设备就绪后拿到 timestamp 粒度：ns 每 tick
float periodNs = 0.f;
{
    VkPhysicalDeviceProperties props{};
    vkGetPhysicalDeviceProperties(physDev, &props);
    periodNs = props.limits.timestampPeriod;
}

// 2) 建一个 timestamp 查询池，两个槽位对应帧的起止
VkQueryPoolCreateInfo qci{};
qci.sType       = VK_STRUCTURE_TYPE_QUERY_POOL_CREATE_INFO;
qci.queryType   = VK_QUERY_TYPE_TIMESTAMP;
qci.queryCount  = 2;
VkQueryPool tsPool = VK_NULL_HANDLE;
vkCreateQueryPool(device, &qci, nullptr, &tsPool);

// 3) 录制：命令串两端各写一个时间戳，中间是被测的绘制
vkCmdResetQueryPool(cmd, tsPool, 0, 2);
vkCmdWriteTimestamp(cmd, VK_PIPELINE_STAGE_TOP_OF_PIPE_BIT,      tsPool, 0);
DrawScene(cmd);                       // 被测的整个帧
vkCmdWriteTimestamp(cmd, VK_PIPELINE_STAGE_BOTTOM_OF_PIPE_BIT,   tsPool, 1);
vkEndCommandBuffer(cmd);

// 4) 提交并等本帧 fence，确保结果可读
vkQueueSubmit(queue, 1, &submit, frameFence);
vkWaitForFences(device, 1, &frameFence, VK_TRUE, UINT64_MAX);

// 5) 读回：WAIT_BIT 让驱动自旋到就绪；未就绪不要当成 0
uint64_t ticks[2] = {0, 0};
VkResult r = vkGetQueryPoolResults(device, tsPool, 0, 2,
    sizeof(ticks), ticks, sizeof(uint64_t),
    VK_QUERY_RESULT_64_BIT | VK_QUERY_RESULT_WAIT_BIT);
if (r == VK_SUCCESS) {
    double gpuMs = (double)(ticks[1] - ticks[0]) * periodNs * 1e-6;
    printf("GPU frame: %.3f ms\n", gpuMs);   // 仅在同一次提交内相减
}
\end{lstlisting}

\subsection{管线统计与逐 pass 计时}

除了起止时间戳，还有 ``管线内部分布'' 的粒度：\DomGPU 会把若干时间戳查询
点放在固定阶段（如片元着色前后），Vulkan 的 \Tp{VK\_QUERY\_TYPE\Bk\_PIPELINE\Bk\_*}
系列与 OpenGL 的 \Tp{EXT\_disjoint\_timer\_query} 提供分段计时；OpenGL 侧要注意
disjoint 标志，一旦驱动做了频率缩放，这段计时不可信，要整段丢弃。

\begin{warning}[title={occlusion query 的副作用}]
遮挡查询（\Tp{VK\_QUERY\_TYPE\Bk\_OCCLUSION\_BIT}、OpenGL 的
\Tp{glGetQueryObjectuiv}）能让驱动把看不见的事件 ``丢进垃圾桶''，省掉后续
片元处理。但读回是异步的：本帧提交、往往下一帧才能拿到结果。若为 ``拿结果''
而强制同步，就把并行度抹平了。成熟做法是把查询结果推迟一帧使用，并接受
``一帧误差''，切勿在热路径里每帧都 \Fn{Wait}。
\end{warning}

\section{瓶颈判据：ALU、带宽、提交与同步气泡}

拿到帧时间后，要回答的是 ``慢在哪一类资源''。四类瓶颈各有可操作的判据，
也有各自的首选修法，汇总如下。

\begin{longtable}{>{\raggedright\arraybackslash}p{0.14\textwidth}
                  >{\raggedright\arraybackslash}p{0.23\textwidth}
                  >{\raggedright\arraybackslash}p{0.23\textwidth}
                  >{\raggedright\arraybackslash}p{0.23\textwidth}}
\caption{瓶颈类型 $\times$ 现象 $\times$ 度量方法 $\times$ 首选修法}\label{tab:ch14-bottleneck} \\
\toprule
瓶颈类型 & 典型现象 & 度量方法 & 首选修法 \\
\midrule
\endfirsthead
\caption[]{瓶颈类型 $\times$ 现象 $\times$ 度量方法 $\times$ 首选修法（续）}\\
\toprule
瓶颈类型 & 典型现象 & 度量方法 & 首选修法 \\
\midrule
\endhead
\bottomrule
\endfoot
ALU 受限 & 帧率随着色复杂度线性下降，改分辨率几乎无感 & GPU 时间戳随像素与顶点数增长，ALU 占用高 & 降精度到 mediump，合并指令，砍循环与分支，用快速近似函数 \\
带宽受限 & 分辨率越高越慢，降采样立刻变快 & 采样率与带宽类计数上升，纹理 cache 命中低 & 压缩纹理（BC、ASTC、ETC2），正确选 mip，减少 overdraw，格式退到 RG8 或 R8 \\
提交受限 & \DomCPU 帧时间明显高于 \DomGPU & 每帧 DrawCall 与 \Fn{vkQueueSubmit} 计数偏高 & 合并命令缓冲，减少状态切换，用实例化与批处理摊薄提交 \\
同步气泡 & \DomGPU 出现空闲窗口，整体卡顿 & 时间戳区间与 \DomCPU 对不齐，fence 等待时间长 & 提高帧队列深度，双缓冲或三缓冲，compute 与图形重叠，延迟读回 \\
overdraw & 半透明叠加与大三角形拖慢光栅 & 片元处理次数远超屏幕像素数 & 先画不透明，加深度预通道，对贴花与粒子做裁剪 \\
\end{longtable}

判据的核心是两条对照：\DomCPU 时间对 \DomGPU 时间（决定是提交受限还是资源受限），
以及 ``改分辨率'' 与 ``改着色复杂度'' 两种扰动各自的敏感度（决定是带宽受限还是
ALU 受限）。扰动法比任何单一指标都可靠。

\subsection{扰动法：三步把瓶颈钉死}

单看一张时间轴容易误判，用三次扰动交叉验证：
\begin{enumerate}
  \item 把渲染分辨率减半（着色复杂度不变）：帧时间同比例下降，指向带宽或
        overdraw；几乎不动，指向 ALU 或提交。
  \item 把片元着色器换成输出常色（分辨率不变）：帧时间大跌，指向 ALU；
        不动，说明瓶颈在着色之外。
  \item 把 DrawCall 数减半而像素覆盖不变（用实例化批处理）：帧时间下降
        明显，指向提交受限；不动，说明提交不是瓶颈。
\end{enumerate}
三次扰动各自单独跑，取交集就是可动手的位置，避免 ``看着像带宽其实是要修
同步气泡'' 的错判。

\begin{compare}[title={ALU 与带宽，一句话区分}]
ALU 受限怕 ``算得多''，加像素或加指令才慢；带宽受限怕 ``搬得多''，加分辨率
或加采样才慢。把这两句话当成快速分诊的入口。
\end{compare}

\section{带宽的三个大头：overdraw、纹理与写回}

带宽瓶颈很少是单点，而是三处的叠加。第一是 overdraw：同一片屏幕像素被多次
写入，半透明、粒子、无深度预通道的全屏 pass 都是元凶；先不透明后透明的排序、
early-Z、贴花的 alpha test 能直接砍掉冗余写入。第二是纹理：未压缩的 RGBA8
一张 4K 就是六十几 MB，取样的字节量决定带宽；用 BC、ASTC、ETC2 把每像素
压到不到一字节，配合 mip 让远处只读低分辨率层，是收益最稳的一招。第三是
写回与 resolve：MSAA 的 resolve、HDR 目标、后处理的多次全屏读写，都在往
显存带宽上叠加；能降到 RG16F 或 R11G11B10F 的中间缓冲就别留 RGBA16F，
能合并的 pass 就别拆成两次全屏读写。

\subsection{纹理压缩格式与各自的适用面}

压缩既省带宽也省显存，但格式选错会掉质量或掉兼容。按 ``块大小 / 通道 /
适用场景'' 记：
\begin{itemize}
  \item BC 系列（桌面与主机）：固定 $4\times4$ 块。BC1 是无独立 alpha 的
        RGB，压得最狠；BC3 给 RGBA；BC4 单通道、BC5 双通道，适合法线贴图
        这类只有两三个通道的数据；BC6H 面向 HDR 浮点 RGB；BC7 是高质量
        RGBA，通道模式多，代价是编码更慢。
  \item ETC2（移动基线）：OpenGL ES 3.1 起为必修格式，适合不带或带少量
        alpha 的场景纹理；它的单通道变体适合遮罩类贴图。
  \item ASTC（移动端与主机）：块尺寸从 $4\times4$ 到 $12\times12$ 可选，
        能在质量与体积间连续权衡，同时覆盖 LDR 与 HDR，跨平台一致性最好。
  \item PVRTC（已淘汰）：早期移动 GPU 的格式，质量与通道适配都落后，
        新项目不应再选它，仅做旧资产兼容时才碰到。
\end{itemize}
选型口诀：桌面/主机优先 BC7 与 BC6H，移动优先 ASTC，需要极小体积再退到
ETC2，绝不新增 PVRTC。

\section{显存预算、纹理流送与 residency}

显存不是越大越好的黑盒。先做预算表：把常驻资源按用途分成渲染目标、纹理
图集、顶点与索引、加速结构几档，各自设上限；再留一档 headroom 给驱动内部
与 swapchain。超预算时不要靠祈祷，而要显式 ``驻留''（residency）管理：
\Tp{VK\_{\Bk}MEMORY\_{\Bk}PROPERTY\_{\Bk}FLAG\_{\Bk}DEVICE\Bk\_{\Bk}LOCAL\_{\Bk}BIT} 放常驻热数据，
host-visible 放上传中的冷数据，配合 \DomGPU 侧的 eviction 与纹理流送，让
远处与低优先级的纹理停留在低 mip 或压缩态，被视到再升级。流送的判据是
``这一帧实际采样到的 mip 级别''，据此决定哪些 mip 可以换出。

\subsection{residency 的三步落地}

\begin{enumerate}
  \item 分类：把资源标成 ``常驻热''（当前帧一定用到）、``可换出''（远处或
        低优先级）与 ``在途''（正在上传或下载）三类，分别放不同内存属性。
  \item 度量：跟踪每帧真实的显存峰值与 swapchain 占用，留出驱动内部与
        fragmentation 的 headroom，别把预算顶满。
  \item 换出：按可见性与相机距离决定升/降级哪个 mip，纹理流送只搬用到的
        mip 层，避免整张 mipmap 链一起上传打爆带宽。
\end{enumerate}
流送与 residency 的失败模式通常是 ``抖动''：一帧升一帧降，反复搬同一张图。
解决办法是给升降级加迟滞（hysteresis），宁可晚一点升，也不要来回搬。

\begin{note}[title={测带宽要先测采样率而非频率}]
多数 ``感觉是带宽'' 的问题，根因是采样次数或写次数超标，而不是显存频率不够。
先数纹理取样次数与片元写入次数，再谈硬件上限。
\end{note}

\section{调试工具族：按平台看能拿到什么}

工具按平台分工，但思路一致：抓帧看静态状态、跑 profile 看动态时间。
\Lib{RenderDoc} 跨 Vulkan、D3D11/12 与 OpenGL，抓一帧成
\Tp{.rdc} 后能逐 DrawCall 回放、看 pipeline state、绑定资源、shader 反汇编
与单个像素的调试——它给的是 ``这一帧的真相''，看不到跨帧的时间分布与真实
频率缩放。时间线要看各平台自带或厂商的 profiler：Windows 的
GPUView、AMD 的 Radeon GPU Profiler、Intel 的 GPA，macOS/iOS 用
Xcode 的 GPU frame capture；D3D12 还应在开发期开 Agility SDK 的
验证层、配合 timestamp query heap 铺时间轴。能看到的：命令序列、绑定、
每个 pass 的耗时；看不到的：驱动内部实现细节与编译器之后的真实调度。

\subsection{抓一帧之后先看哪三样}

\begin{enumerate}
  \item pipeline state：混合方程、深度测试/写入、cull 模式、拓扑——花屏
        多半在这里（例如 blend 关掉导致覆盖异常、depth write 缺失导致排序错）。
  \item 绑定资源：这一 draw 实际采样/写入的是哪张纹理、哪个 buffer 区间，
        黑屏常是绑到了空的或 layout 不对的资源。
  \item 单个像素/顶点的 shader 回放：给定屏幕坐标，看片元执行到的分支与
        中间量，把 ``算错'' 与 ``没算'' 区分开。
\end{enumerate}
时间分布不在抓帧工具里，要用平台 profiler 的时间轴与 GPU 时间戳补上；两类
工具配合，静态状态和动态耗时才能对上。

\begin{warning}[title={验证层消息不等于性能问题}]
Vulkan validation、D3D12 Agility SDK 的 debug layer 报的是 ``用法违例''
（资源状态、同步、描述符范围），它会显著拖慢帧时间。别把开验证层后的慢当成
性能瓶颈——那多半是同步检查把异步摊平了的副作用。基线测量要在关掉验证层的
构建上做。
\end{warning}

\begin{lstlisting}[style=shellstyle,caption={抓帧、压缩与帧时间统计的命令组合},label={lst:ch14-tools}]
# 1) 抓一帧为 .rdc，再用 RenderDoc GUI 打开看 pipeline state
renderdoccmd capture ./bin/demo
# 2) 采集能力：确认支持的压缩格式与 timestamp 粒度
vulkaninfo --summary
# 3) 离线压纹理并生成 mip：KTX2/UASTC（GPU 侧也可 BC、ASTC、ETC2）
toktx --t2 --2d --genmipmap --encode UASTC hero.ktx2 hero.png
etcpak -3 hero.png hero_etc2.pkm
# 4) 帧时间分位数：从采样日志算中位与 p95，避开单帧尖峰的误导
awk -F, 'NR>1{print $3}' frames.csv | sort -n |
  awk '{a[NR]=$1} END{printf "med=%.2fms p95=%.2fms\n", a[int(NR/2)], a[int(NR*0.95)]}'
\end{lstlisting}

\section{选型判断：自研、抽象层还是引擎内改}

把上面的度量流程跑通后，回到路线选择。若瓶颈是 ``提交气泡与显式控制''，
且需要跨帧的精细资源编排，直接自研显式 \Tp{Vulkan}/{\Bk}\Tp{D3D12}/{\Bk}\Tp{Metal}
才划算；抽象层库（如 \Lib{bgfx}、\Lib{wgpu}、\Lib{Filament}）适合要覆盖
多后端、又不想自己维护同步与内存的中小团队；若目标本就在引擎里，优先在
引擎的渲染抽象接口内改，收益来自减少它的默认开销，而不是绕开它。三条路的
共同前提都是本章的度量能力——先能量出瓶颈，再谈选型。

\subsection{决策路径速记}

\begin{itemize}
  \item 需要跨帧资源编排、显式多队列、极致提交控制 $\rightarrow$ 自研显式
        API（\Tp{Vulkan}/{\Bk}\Tp{D3D12}/{\Bk}\Tp{Metal}），代价是自己维护同步与内存。
  \item 要覆盖多后端但团队不大 $\rightarrow$ 抽象层库（\Lib{bgfx}、\Lib{wgpu}、
        \Lib{Filament}），把同步与移植交给库，牺牲一部分控制粒度。
  \item 目标本就在引擎内 $\rightarrow$ 在引擎的渲染抽象接口内改，优先削
        默认开销（多余 pass、未压缩纹理、全屏读写），而不是绕开引擎。
\end{itemize}
具体库与许可见附录 \ref{ap:apC-library-choice}，语言与 API 对照见附录
\ref{ap:apB-shading-langs}。

\begin{guideline}[title={测量先于优化，扰动验证结论}]
优化前先用时间戳与抓帧建立 ``CPU/GPU 两条时间线 + 一张瓶颈类型表''；任何
结论都用扰动法复验（改分辨率、改着色复杂度、改提交粒度各测一次）。看不懂的
抖动不要猜，抓一帧下来按 pass 拆开看，比翻代码更快定位。
\end{guideline}

% ============================================================
% ---- frag_apA.tex ----
\appendix
\chapter{术语与缩写}
\label{ap:apA-glossary}

本册的术语按\emph{读者会在哪一处撞上它}来分组，而不是按字母顺序：管线与阶段回答
``这段代码在哪一级执行''，着色语言与编译链回答``这个名字属于哪一侧''，图形 API 与对象
回答``句柄是谁创建的''，性能与调试回答``指标在量什么''。
每条定义都写成能独立读懂的一句话，必要时给出对应的 API 名字；括号里的编号是相关章节。
术语本身不重复正文的推导，遇到分歧写法以本册用词约定（本章末尾的要点框）为准。

\section{管线与阶段}

``阶段''指可编程或固定功能的执行环节，一个绘制调用会依次穿过它们。
下表的定义强调\emph{输入与输出是什么}，因为绝大多数关于管线的争论其实来自把某一级的
输入当成另一级的输入。
另外要注意``阶段存在''与``阶段可用''是两件事：几何着色器与曲面细分在不少移动驱动上
要么不提供、要么贵到不该用，本册提到它们时都当作可选能力，具体取舍见 \ref{ch:ch01-boundary}。

\begin{longtable}{>{\raggedright\arraybackslash}p{0.20\textwidth}
                  >{\raggedright\arraybackslash\small}p{0.11\textwidth}
                  >{\raggedright\arraybackslash}p{0.40\textwidth}
                  >{\raggedright\arraybackslash}p{0.09\textwidth}}
\caption{管线与阶段术语}\label{tab:apA-pipeline} \\
\toprule
术语 & 缩写 & 一句定义 & 相关章 \\
\midrule
\endfirsthead
\toprule
术语 & 缩写 & 一句定义 & 相关章 \\
\midrule
\endhead
\bottomrule
\endfoot
输入装配 & IA & 读取顶点缓冲与索引，把它们拼成线段或三角形这类图元的阶段。 & \ref{ch:ch02-raster-pipeline} \\
顶点着色器 & VS & 每个顶点执行一次的可编程阶段，唯一职责是写出裁剪空间坐标。 & \ref{ch:ch04-transform-space} \\
属性取指 & Vertex Fetch & 由输入装配把顶点属性从缓冲里读出来并按声明类型转换的硬件动作。 & \ref{ch:ch03-vertex-input} \\
变换后缓存 & PTC & 缓存已通过顶点着色器的图元顶点，供共享同一顶点的相邻三角形复用。 & \ref{ch:ch02-raster-pipeline} \\
曲面细分 & Tess & 由外壳与内部着色器在运行时把图元分裂成更多图元的两级阶段。 & \ref{ch:ch02-raster-pipeline} \\
几何着色器 & GS & 以整个图元为单位、可以增删图元并按层选择输出目标的阶段。 & \ref{ch:ch02-raster-pipeline} \\
网格着色器 & MS & 用网格与任务两级阶段取代输入装配、顶点着色器与几何着色器的新模型。 & \ref{ch:ch08-compute-driven} \\
流输出 & TF & 把顶点阶段的输出写回缓冲而不是继续走光栅化的反馈路径。 & \ref{ch:ch08-compute-driven} \\
图元 & Prim & 点、线、三角形这三类拓扑的统称，是光栅化的输入单位。 & \ref{ch:ch02-raster-pipeline} \\
绕序 & Winding & 用顶点在屏幕上的旋转方向定义正反面、供剔除与混合决策的规则。 & \ref{ch:ch04-transform-space} \\
片元丢弃 & Discard & 在片元阶段提前结束执行、既不写颜色也不写深度的动作。 & \ref{ch:ch05-glsl} \\
光栅化 & Rast & 确定哪些像素被图元覆盖，并为每个覆盖位置生成片元的固定功能阶段。 & \ref{ch:ch02-raster-pipeline} \\
片元着色器 & FS & 每个候选像素位置执行一次、产出颜色与深度并可丢弃该片元的阶段。 & \ref{ch:ch05-glsl} \\
像素着色器 & PS & 片元着色器在 \HSL\ 生态里的同义叫法，强调它输出到渲染目标。 & \ref{ch:ch06-hlsl} \\
早深度测试 & Early-Z & 在片元着色器之前用插值深度做测试，用来跳过必然不可见的片元。 & \ref{ch:ch02-raster-pipeline} \\
隐藏面消除 & HSR & 靠深度缓冲或顺序无关的办法去掉被遮挡像素的统称。 & \ref{ch:ch14-perf-tools} \\
输出合并 & OM & 把片元结果与已有内容按混合、深度、模板规则合成一处的末级阶段。 & \ref{ch:ch02-raster-pipeline} \\
混合 & Blend & 按源与目的的因子组合片元颜色与帧缓冲原值的逐像素运算。 & \ref{ch:ch02-raster-pipeline} \\
多重采样 & MSAA & 一个像素内放多个覆盖样本以降低边缘锯齿，颜色仍只算一次的方案。 & \ref{ch:ch02-raster-pipeline} \\
解析 & Resolve & 把多重采样附件的多个样本平均成一张单采样纹理的动作。 & \ref{ch:ch09-vulkan} \\
计算着色器 & CS & 不产生图元、按工作组在任意缓冲上并行读写的可编程阶段。 & \ref{ch:ch08-compute-driven} \\
工作组 & WG & 计算调度里共享局部存储并可能共享同步点的一批调用。 & \ref{ch:ch08-compute-driven} \\
绘制调用 & DC & 从应用提交一次绘制所引发的全部驱动与硬件开销的统称。 & \ref{ch:ch14-perf-tools} \\
实例化 & Inst & 用同一段几何配不同数据重复绘制多份的阶段能力。 & \ref{ch:ch03-vertex-input} \\
剔除 & Cull & 在光栅化前按背面朝向或视锥关系丢掉图元的动作总称。 & \ref{ch:ch04-transform-space} \\
平铺延迟渲染 & TBDR & 把屏幕切成小块、片上完成着色与合并的架构家族，移动端最典型。 & \ref{ch:ch01-boundary} \\
子通道 & Subpass & 同一个渲染通道内共享平铺存储、按依赖顺序衔接的阶段划分。 & \ref{ch:ch09-vulkan} \\
加载与存储动作 & L/S & 声明每个通道开始时附件如何初始化、结束时是否写回内存的两组开关。 & \ref{ch:ch09-vulkan} \\
\end{longtable}

\section{着色语言与编译链}

这一节的术语横跨两侧：语言名字属于作者，中间表示与阶段名属于工具链。
下表刻意不比较优劣，只把``哪个词属于哪一层''说清楚，避免把编译期概念当成运行期能力。

\begin{note}[title={前端、中端、后端的分界}]
  把源码编译成中间表示的那一步叫前端，改写中间表示的那一步叫中端，
  由驱动把中间表示编成硬件指令的那一步发生在运行期、叫后端。
  本册说``编译''时默认指前两步，因为它们是应用能自己控制的；
  涉及驱动那一步时一律写成``运行期编译''或``驱动优化''，避免和离线编译混为一谈。
\end{note}

\begin{longtable}{>{\raggedright\arraybackslash}p{0.20\textwidth}
                  >{\raggedright\arraybackslash\small}p{0.11\textwidth}
                  >{\raggedright\arraybackslash}p{0.42\textwidth}
                  >{\raggedright\arraybackslash}p{0.09\textwidth}}
\caption{着色语言与编译链术语}\label{tab:apA-lang} \\
\toprule
术语 & 缩写 & 一句定义 & 相关章 \\
\midrule
\endfirsthead
\toprule
术语 & 缩写 & 一句定义 & 相关章 \\
\midrule
\endhead
\bottomrule
\endfoot
着色语言（OpenGL 系） & \GSL & 面向 OpenGL 与 Vulkan 的着色语言，以 \Tp{layout} 限定符描述接口。 & \ref{ch:ch05-glsl} \\
着色语言（Direct3D 系） & \HSL & 面向 Direct3D 的着色语言，以语义（semantic）绑定输入输出。 & \ref{ch:ch06-hlsl} \\
着色语言（Metal 系） & \MSL & Metal 的着色语言，语法接近 \CPP\ 且带地址空间限定符。 & \ref{ch:ch10-d3d12-metal} \\
着色中间表示 & \SPIV & 二进制、显式控制流的着色中间表示，Vulkan 的官方输入格式。 & \ref{ch:ch07-spirv-toolchain} \\
入口 & Entry & 一个编译单元里被管线某一级调用的那个函数及其阶段标签。 & \ref{ch:ch07-spirv-toolchain} \\
语义 & Sem & \HSL\ 里附加在变量名后面、用来标识接口位置的标注。 & \ref{ch:ch06-hlsl} \\
位置 & Loc & \GSL\ 里标识输入输出槽位编号的限定符取值。 & \ref{ch:ch05-glsl} \\
内建变量 & Builtin & 语言预先声明、携带固定含义（如顶点下标、片元坐标）的变量。 & \ref{ch:ch05-glsl} \\
精度限定 & Precision & 声明浮点运算位宽的限定符，在移动设备上直接影响可用性与功耗。 & \ref{ch:ch05-glsl} \\
着色器模型 & SM & Direct3D 版本与可用特性组合的编号体系。 & \ref{ch:ch10-d3d12-metal} \\
编译数据库 & CDB & 记录每个翻译单元的真实编译参数、供工具复现编译现场的清单。 & \ref{ch:ch07-spirv-toolchain} \\
预处理器 & PP & 在真正编译前展开宏与条件包含、决定变体分支的那一步。 & \ref{ch:ch07-spirv-toolchain} \\
反射 & Reflect & 从编译产物里读回接口（名称、类型、占用）以便布局绑定描述符的过程。 & \ref{ch:ch07-spirv-toolchain} \\
变体 & Variant & 同一份源码用不同宏组合编译出的多个产物之一。 & \ref{ch:ch07-spirv-toolchain} \\
重定位 & Remap & 批量改写产物里的标识符以合并重复资源的处理步骤。 & \ref{ch:ch07-spirv-toolchain} \\
规范化 & Norm & 把前端产物的方言差异抹平、使其符合目标运行环境约束的改写。 & \ref{ch:ch07-spirv-toolchain} \\
反编译 & Decomp & 从中间表示还原成某种可读源码语言的动作，常用于跨 API 移植。 & \ref{ch:ch07-spirv-toolchain} \\
能力 & Cap & 中间表示里声明这段代码需要哪类硬件特性的条目。 & \ref{ch:ch07-spirv-toolchain} \\
执行模型 & Exec & 标识这段代码属于哪一级阶段、如何取输入的那份声明。 & \ref{ch:ch07-spirv-toolchain} \\
内存模型 & Memory Model & 规定并行执行者之间可见性与顺序约束的那套规则。 & \ref{ch:ch08-compute-driven} \\
子组 & Subgroup & 比工作组更细一层、能做廉价内部通信的执行粒度，跨子组不保证任何顺序。 & \ref{ch:ch08-compute-driven} \\
通道遮蔽 & Swizzle & 用 \Tp{.xyzw}\ 的重排写法取用向量分量，本身不产生额外指令。 & \ref{ch:ch05-glsl} \\
组内共享存储 & Shared & 同一工作组内互相可见的低延迟存储，容量与生命周期都受硬件限制。 & \ref{ch:ch08-compute-driven} \\
接口块 & Block & 把一组变量打包声明、整体绑定到缓冲的语言结构。 & \ref{ch:ch03-vertex-input} \\
行优先与列优先 & RM 与 CM & 矩阵在内存里按行或按列连续的排布方向，编译选项与语言默认不同。 & \ref{ch:ch04-transform-space} \\
原子操作 & Atomic & 对缓冲元素做不可分割读改写的指令，并行写同一位置的前提。 & \ref{ch:ch08-compute-driven} \\
汇编与反汇编 & ASM & 把二进制的 \SPIV\ 模块与可读文本互相转换的一对动作。 & \ref{ch:ch07-spirv-toolchain} \\
\end{longtable}

\begin{lstlisting}[style=shellstyle,caption={编译链上最常出现的六个命令名：读它们的输出时请对齐上面的术语},label={lst:apA-tools}]
# 前端：源码 -> SPIR-V
glslangValidator -V -S vert vertex.vert -o vert.spv
glslc --target-env=vulkan1.2 --target-spv=spv1.5 shader.vert -o shader.spv
dxc -T vs_6_0 -E main -spirv fx.hlsl -Fo fx.spv

# 检查与改写
spirv-val vert.spv
spirv-dis vert.spv | head
spirv-opt -O -o vert.opt.spv vert.spv

# 反向：SPIR-V -> 可读源码
spirv-cross --msl --entry main fx.spv
\end{lstlisting}

\section{图形 API 与对象}

句柄的名字在各 API 里互相错开，但职责几乎一一对应；下表按职责而不是按厂商排列，
缩写列给出的是最常见的两家写法。
定义里出现的具体枚举名一律按原样书写，下划线不属于排版装饰。
读表时请记住一条共同的坑：同一个对象被多个线程同时使用时，多数规范都要求应用自己做
外部同步；能否跨帧复用属于对象语义的一部分，而不是实现细节，见 \ref{ch:ch09-vulkan}。

\begin{longtable}{>{\raggedright\arraybackslash}p{0.20\textwidth}
                  >{\raggedright\arraybackslash\small}p{0.11\textwidth}
                  >{\raggedright\arraybackslash\small}p{0.42\textwidth}
                  >{\raggedright\arraybackslash}p{0.09\textwidth}}
\caption{图形 API 与对象术语}\label{tab:apA-api} \\
\toprule
术语 & 缩写 & 一句定义 & 相关章 \\
\midrule
\endfirsthead
\toprule
术语 & 缩写 & 一句定义 & 相关章 \\
\midrule
\endhead
\bottomrule
\endfoot
物理设备 & Adapter & 系统里一块真实图形硬件的枚举表示，用于查询能力与极限。 & \ref{ch:ch09-vulkan} \\
逻辑设备 & Device & 应用在硬件上的会话句柄，几乎所有对象都从它创建。 & \ref{ch:ch09-vulkan} \\
队列 & Queue & 提交给硬件执行的命令流通道，可分图形、传输等用途。 & \ref{ch:ch09-vulkan} \\
命令缓冲 & CmdBuf & 录制好的一串绘制与状态切换，可重复提交给队列。 & \ref{ch:ch09-vulkan} \\
命令编码器 & Encoder & Metal 里边录制边维护状态的命令记录对象。 & \ref{ch:ch10-d3d12-metal} \\
渲染通道 & RP & 声明附件、加载存储动作与子通道依赖的帧结构描述。 & \ref{ch:ch09-vulkan} \\
渲染目标 & RT & 可被输出合并阶段写入的纹理级别，颜色或深度皆可。 & \ref{ch:ch10-d3d12-metal} \\
帧缓冲 & FB & 一组附件按层级绑定成的可绘制整体。 & \ref{ch:ch09-vulkan} \\
管线对象 & PSO & 把阶段、布局与固定功能状态合并后的一次编译结果。 & \ref{ch:ch09-vulkan} \\
管线布局 & PL & 声明描述符集与推送常量分层的静态布局。 & \ref{ch:ch09-vulkan} \\
根签名 & RS & Direct3D 12 里等价的参数布局描述，绑定表在此定义。 & \ref{ch:ch10-d3d12-metal} \\
描述符 & Descr & 指向缓冲或纹理实际位置的小型不透明记录，按集合批量更新。 & \ref{ch:ch09-vulkan} \\
参数缓冲 & AB & Metal 里把一组资源与常量打包提交的等价机制。 & \ref{ch:ch10-d3d12-metal} \\
绑定无关 & Bindless & 在着色器里用下标直接寻址资源、绕开逐绘制绑定的做法。 & \ref{ch:ch09-vulkan} \\
交换链 & SC & 与窗口关联、负责轮转与送显的图像集合。 & \ref{ch:ch09-vulkan} \\
信号量 & Sem & 队列之间或队列与主机之间的顺序同步原语。 & \ref{ch:ch09-vulkan} \\
围栏 & Fence & 主机可以等待、用来确认一批提交已完成的同步原语。 & \ref{ch:ch09-vulkan} \\
事件 & Event & 可以在录制中被置位与等待的细粒度同步对象。 & \ref{ch:ch10-d3d12-metal} \\
图像布局 & Layout & 一张图在某一时刻的物理排布与用途声明，转换需屏障。 & \ref{ch:ch09-vulkan} \\
缓冲视图 & BV & 把一段缓冲按指定格式解释成纹理或采样数据的窗口。 & \ref{ch:ch09-vulkan} \\
采样器状态 & SS & 描述滤波、寻址与各向异性的一组不可变参数。 & \ref{ch:ch02-raster-pipeline} \\
查询池 & QP & 存放时间戳、遮挡结果与管线条目计数的对象。 & \ref{ch:ch14-perf-tools} \\
内存属性位 & Flags & 声明这段内存该被主机还是设备访问、是否缓存的一组位，例如
  \Tp{VK\_MEMORY\_PROPERTY\_\Bk HOST\_VISIBLE\_BIT}。 & \ref{ch:ch09-vulkan} \\
内存屏障 & Barrier & 声明两侧访问在顺序与可见性上的约束、把乱序执行钉住的原语。 & \ref{ch:ch09-vulkan} \\
子分配 & Suballoc & 在一大块设备内存里按偏移与长度切出多个对象的管理手法。 & \ref{ch:ch09-vulkan} \\
对象池 & Pool & 预先批量分配、再按请求切分描述符或命令对象的容器。 & \ref{ch:ch10-d3d12-metal} \\
扩展 & Ext & 版本之外由标准组织或厂商提供的可选能力，启用前必须查支持位。 & \ref{ch:ch09-vulkan} \\
校验层 & Layer & 在提交时检查参数与状态机合法性、可选装进实例的调试组件。 & \ref{ch:ch14-perf-tools} \\
\end{longtable}

\section{性能与调试}

性能术语的危险在于同一个数字在不同架构里含义不同：``占用率高''在桌面意味着并行度充足，
在平铺架构里可能只是片上存储被占满。下表尽量把定义绑到\emph{可直接观测}的对象上。
另一个通用的读法纪律是：先确定量的是 \DomCPU\ 还是 \DomGPU\ 的时间，再谈大小；
两侧的数字放进同一张柱状图对比是最常见的误判来源，做法见 \ref{ch:ch14-perf-tools}。

\begin{longtable}{>{\raggedright\arraybackslash}p{0.20\textwidth}
                  >{\raggedright\arraybackslash\small}p{0.11\textwidth}
                  >{\raggedright\arraybackslash}p{0.42\textwidth}
                  >{\raggedright\arraybackslash}p{0.09\textwidth}}
\caption{性能与调试术语}\label{tab:apA-perf} \\
\toprule
术语 & 缩写 & 一句定义 & 相关章 \\
\midrule
\endfirsthead
\toprule
术语 & 缩写 & 一句定义 & 相关章 \\
\midrule
\endhead
\bottomrule
\endfoot
帧时长 & FT & 连续两帧之间的墙上时间，比每秒帧数更能暴露卡顿。 & \ref{ch:ch14-perf-tools} \\
帧预算 & Budget & 单帧允许消耗的时间上限，由刷新率倒推得出。 & \ref{ch:ch14-perf-tools} \\
瓶颈侧 & Bound & 判断是 \DomCPU\ 提交跟不上还是 \DomGPU\ 执行跟不上。 & \ref{ch:ch14-perf-tools} \\
气泡 & Bubble & 提交流水线里 \DomGPU\ 空转、等 \DomCPU\ 产活的那段空隙。 & \ref{ch:ch09-vulkan} \\
停顿 & Stall & 因为等待前序结果而无法继续执行的阻塞。 & \ref{ch:ch09-vulkan} \\
过度绘制 & Overdraw & 同一像素被多个片元重复着色的浪费量。 & \ref{ch:ch02-raster-pipeline} \\
带宽 & BW & 单位时间在显存与芯片之间搬运的字节数，常是移动端的硬上限。 & \ref{ch:ch03-vertex-input} \\
占用率 & Occ & 一个计算单元上实际驻留的执行组数与理论上限之比。 & \ref{ch:ch08-compute-driven} \\
分支发散 & Div & 同一执行组内的调用走了不同分支而被迫串行的现象。 & \ref{ch:ch08-compute-driven} \\
特殊函数单元 & SFU & 负责三角函数、倒数平方根等超越运算的专用硬件。 & \ref{ch:ch12-lighting-pbr} \\
时间戳查询 & TS & 在命令流里插入标记以测量各段 \DomGPU\ 耗时的机制。 & \ref{ch:ch14-perf-tools} \\
性能计数器 & PC & 硬件统计单元，给出图元数、片元数、缓存命中这类原始量。 & \ref{ch:ch14-perf-tools} \\
帧捕获 & Capture & 抓取一帧的全部调用与资源以便回放与逐绘制检查的调试手段。 & \ref{ch:ch14-perf-tools} \\
调试标签 & Label & 应用主动打上的名字，用来在工具里辨认通道与资源。 & \ref{ch:ch14-perf-tools} \\
垂直同步 & VSync & 把送显对齐到刷新周期以消除撕裂、代价是延迟上升的开关。 & \ref{ch:ch09-vulkan} \\
呈现模式 & Present & 在立即、信箱、先进先出这三档里权衡延迟、吞吐与功耗的选择。 & \ref{ch:ch09-vulkan} \\
多重缓冲 & NBUF & 用两到三份帧资源让绘制与显示重叠进行的常用做法。 & \ref{ch:ch09-vulkan} \\
细节层次 & LOD & 按屏幕占比或距离切换网格与贴图精度以省算力的策略。 & \ref{ch:ch14-perf-tools} \\
分配粒度 & Gran & 驱动要求资源在内存里对齐的最小单位，跨它访问需要额外屏障。 & \ref{ch:ch09-vulkan} \\
帧图 & FG & 把一帧内所有通道与它们的依赖显式建成一张图再统一调度的框架。 & \ref{ch:ch09-vulkan} \\
微基准 & Micro & 只重复测量单一操作的极小循环，用来看单价而不能当结论。 & \ref{ch:ch14-perf-tools} \\
提交批次 & Batch & 把多次绘制合并进一次编码与一次提交，用来摊薄驱动开销。 & \ref{ch:ch10-d3d12-metal} \\
未定义行为 & UB & 规范不做任何保证、不同驱动结果可以不同的那一类错误。 & \ref{ch:ch01-boundary} \\
设备丢失 & Lost & 驱动因执行超时或故障重置硬件、应用需要重建全部对象的事件。 & \ref{ch:ch10-d3d12-metal} \\
回归基线 & Baseline & 用来判断一次改动是否让帧时长变差的历史数据点。 & \ref{ch:ch14-perf-tools} \\
\end{longtable}

\begin{warning}[title={三个最容易用错的词}]
  一、``绑定无关''（bindless）描述的是一种\emph{寻址方式}，不是某个 API 的对象，
  不要用它指代描述符集本身。
  二、``子组''（subgroup）不等于工作组：前者的通信代价低但可见范围小，
  跨子组不带来任何顺序与可见性保证，同步必须显式做。
  三、``占用率''来自计算与硬件调度的语境，用在光栅化阶段没有可观测的定义，
  想表达 ``\DomGPU\ 够忙'' 请改用帧时长与时间戳查询的结果。
\end{warning}

\begin{guideline}[title={本册用词约定}]
  四种语言与中间表示一律写成 \GSL、\HSL、\MSL、\SPIV，不写全拼、不写大小写混排的变形，
  也不在复数或所有格上做任何加字；正文、表头与图注同一写法。
  ``CPU 侧''与``GPU 侧''一律用 \DomCPU\ 与 \DomGPU\ 两个标签，避免与
  \Lang{Host}、\Lang{Device} 之类同义词混用；描述统一内存时改用``主机可见''这一具体条件名。
  缩写只在第一次出现时给全称，之后直接用缩写；表格的缩写列固定使用本节的形式。
  \GSL\ 与 \HSL\ 的差异一律用``两侧''表述，不写成``新旧''；\SPIV\ 是中间表示而不是语言，
  任何段落都不把它与 \GSL、\HSL\ 并列成第三种源码。
  凡引用规范里的枚举与结构体名，一律原样保留大小写与下划线，长名允许在行末断字。
\end{guideline}

% ============================================================
% ---- frag_apB.tex ----
\chapter{着色语言与 API 速查}
\label{ap:apB-shading-langs}

正文把 \GSL、\HSL、\MSL 三种着色语言与 \Tp{Vulkan}/{\Bk}\Tp{D3D12}/{\Bk}\Tp{Metal}
三套 API 分开讲。本附录把它们并排到一张桌上，方便移植与读别家代码时快速对
表。三张主表分别覆盖：语言语法与语义对照、API 对象与概念对照、编译与校验
命令速查。所有条目以头文件与官方规范为准，版本差异请核对对应版本官方文档。

\section{着色语言：GLSL、HLSL、MSL 对照}

下表按 ``维度'' 逐行对齐三语言。单元格里的标识符只是常见写法示例，并非穷举。

\begin{longtable}{>{\raggedright\arraybackslash}p{0.15\textwidth}
                  >{\raggedright\arraybackslash}p{0.22\textwidth}
                  >{\raggedright\arraybackslash}p{0.22\textwidth}
                  >{\raggedright\arraybackslash}p{0.22\textwidth}}
\caption{\GSL、\HSL、\MSL 语法与语义对照}\label{tab:apB-lang} \\
\toprule
维度 & \GSL & \HSL & \MSL \\
\midrule
\endfirsthead
\caption[]{\GSL、\HSL、\MSL 语法与语义对照（续）}\\
\toprule
维度 & \GSL & \HSL & \MSL \\
\midrule
\endhead
\bottomrule
\endfoot
入口 & 固定 \Tp{main} & 任意命名函数，用 \Tp{-E} 指定 & 任意命名函数，按 stage 传入 \\
输出（片元） & \Tp{out vec4}，旧式为 \Tp{gl\_FragColor} & 返回带语义的结构，如 \Tp{SV\_Target} & \Tp{[[color(0)]]} 的 \Tp{float4} \\
顶点位置 & 写 \Tp{gl\_Position} & 语义 \Tp{SV\_Position} & 属性 \Tp{[[position]]} \\
绑定方式 & \Tp{layout} 内 \Tp{set} 与 \Tp{binding}（\Tp{Vulkan}），\Tp{binding} 单层（OpenGL） & \Tp{register} 分 \Tp{b}/{\Bk}\Tp{t}/{\Bk}\Tp{s}/{\Bk}\Tp{u} 四类 & \Tp{[[buffer]]}、\Tp{[[texture]]}、\Tp{[[sampler]]} \\
精度限定 & \Tp{highp}、\Tp{mediump}、\Tp{lowp}，ES 片元必填 & 无等价三档；\Tp{half}、\Tp{min16float}、\Tp{precise} & 无 \Tp{lowp} 档；靠 \Tp{half} 与 \Tp{float} 类型 \\
矩阵存储 & 默认列主序，可 \Tp{row\_major} 覆盖 & cbuffer 默认行主序，\Tp{dxc} 可切列主序 & \Tp{float4x4} 默认列主序，可 \Tp{row\_major} \\
向量与通道 & \Tp{vec2} 至 \Tp{vec4}，\Tp{xyzw} 通道 & \Tp{float2} 至 \Tp{float4}，通道同左 & \Tp{float2} 至 \Tp{float4}，通道同左 \\
裁剪与深度 & OpenGL 下 $z$ 为 \Tp{[-1,1]}；\Tp{Vulkan} 下 $y$ 翻转、$z$ 为 \Tp{[0,1]} & 裁剪 $x,y$ 为 \Tp{[-1,1]}，$z$ 为 \Tp{[0,1]}，$y$ 向上 & 裁剪 $x,y$ 为 \Tp{[-1,1]}，$z$ 为 \Tp{[0,1]}，帧缓冲 $y$ 向下 \\
可用阶段 & 顶点、片元、几何、细分控制与求值、计算 & 顶点、像素、几何、域与壳、计算、网格与放大、光追各阶段 & 顶点、片元、计算；无几何着色器，支持细分与光追 \\
采样与纹理 & \Tp{sampler2D} 可 combined，也可 \Tp{texture} 与 \Tp{sampler} 分离 & \Tp{Texture2D} 与 \Tp{SamplerState} 天生分离 & \Tp{texture2d} 与 \Tp{sampler} 分离 \\
内建函数 & \Tp{texture}、\Tp{normalize}、\Tp{mix} & \Tp{Sample}、\Tp{normalize}、\Tp{lerp} & \Tp{sample}、\Tp{normalize}、\Tp{mix} \\
扩展机制 & \Tp{\#extension} 加 \Tp{GL\_} 前缀扩展 & 靠 shader model 与 \Tp{dxc} 能力档 & 靠 \MSL 版本与 \Tp{[[...]]} 属性 \\
\end{longtable}

读这张表时最容易踩的是 ``裁剪与深度'' 一行：\GSL 的默认约定取决于目标
API，同一份 \Tp{gl\_Position} 在 OpenGL 与 \Tp{Vulkan} 下的 $y$ 方向与 $z$
区间不同，投影矩阵要按目标 API 生成，不能一份通吃。下面把这行的两个易错点
展开。

\subsection{入口、输出与系统值的写法}

三语言对 ``顶点输出到裁剪空间'' 这件事的写法不同：\GSL 直接写内建变量
\Tp{gl\_Position}；\HSL 靠语义名（system-value semantic）标注，顶点位置是
\Tp{SV\_Position}、片元颜色是 \Tp{SV\_Target}、深度是 \Tp{SV\_Depth}；\MSL
则用属性标注，位置写 \Tp{[[position]]}、颜色写 \Tp{[[color(0)]]}。片元输出
的 ``多重渲染目标'' 也各自一套：\GSL 用多个 \Tp{out} 变量或
\Tp{layout(location=n)}，\HSL 用 \Tp{SV\_Target0}/{\Bk}\Tp{SV\_Target1} 这样的
编号语义，\MSL 用 \Tp{[[color(n)]]}。移植时最容易漏的是 ``语义名的编号'' 与
``location 的编号'' 必须和 pipeline 的输出布局一一对齐，写错就静默丢通道。

\subsection{精度限定与可移植性}

\GSL 的 \Tp{highp}/{\Bk}\Tp{mediump}/{\Bk}\Tp{lowp} 是显式精度档，在 ES 片元着色器里
必须给默认精度；\HSL 没有同名三档，靠 \Tp{half} 类型与 \Tp{min16float} 之类
的最小精度注解，行为依 GPU 而定；\MSL 只有 \Tp{half} 与 \Tp{float}，靠类型
而非限定符。可移植的写法是：位置与法线用 \Tp{highp}/{\Bk}\Tp{float}，颜色与方向
向量可降到 \Tp{mediump}/{\Bk}\Tp{half}，但要在低端 GPU 上实测半精度是否够，别
假设三语言会给出同样的舍入。

\subsection{向量、矩阵与通道的类型写法}

三语言的向量与矩阵类型名基本一一对应：\GSL 的 \Tp{vec3} 对应 \HSL 的
\Tp{float3} 与 \MSL 的 \Tp{float3}，通道写法都用 \Tp{.xyzw}（也可 \Tp{.rgba}
与 \Tp{.stpq} 复用）。差别在矩阵：\GSL 用 \Tp{mat3}、\Tp{mat4x3} 这样的定名，
\HSL 用 \Tp{float3x3}、\Tp{float4x3}，\MSL 用 \Tp{float3x3} 加行列主序模板
参数。矩阵乘向量时，行列主序与 ``行/列向量约定'' 必须成对匹配，否则转置错误
会让光照整体镜像，而这类 bug 在一张图上很难看出来。移植时优先固定一种约定
（列向量左乘最常用），并在 CPU 侧统一构造矩阵。

\section{图形 API：对象与概念对照}

同一件事在三个 API 里名字不同，下表把它们对齐。左列是概念，右三列是各 API
里的对应对象或机制。

\begin{longtable}{>{\raggedright\arraybackslash}p{0.15\textwidth}
                  >{\raggedright\arraybackslash}p{0.22\textwidth}
                  >{\raggedright\arraybackslash}p{0.22\textwidth}
                  >{\raggedright\arraybackslash}p{0.21\textwidth}}
\caption{\Tp{Vulkan}、\Tp{D3D12}、\Tp{Metal} 对象与概念对照}\label{tab:apB-api} \\
\toprule
概念 & \Tp{Vulkan} & \Tp{D3D12} & \Tp{Metal} \\
\midrule
\endfirsthead
\caption[]{\Tp{Vulkan}、\Tp{D3D12}、\Tp{Metal} 对象与概念对照（续）}\\
\toprule
概念 & \Tp{Vulkan} & \Tp{D3D12} & \Tp{Metal} \\
\midrule
\endhead
\bottomrule
\endfoot
设备与实例 & \Tp{VkInstance} 到 \Tp{VkDevice} & \Tp{ID3D12Device} & \Tp{MTLDevice} \\
队列与命令 & \Tp{VkQueue} 与 \Tp{Vk}\allowbreak\Tp{CommandBuffer} & CommandQueue 与 CommandList & \Tp{MTL}\allowbreak\Tp{CommandQueue} 与 \Tp{MTL}\allowbreak\Tp{CommandBuffer} \\
同步 & \Tp{VkFence}（CPU 侧）与 \Tp{VkSemaphore}（GPU 侧） & \Tp{ID3D12Fence} 加资源状态迁移 & 自然编码顺序加 \Tp{MTL}\allowbreak\Tp{SharedEvent} \\
内存与资源 & \Tp{Vk}\allowbreak\Tp{DeviceMemory} 绑定 \Tp{VkBuffer} 与 \Tp{VkImage} & \Tp{ID3D12Heap} 与 resource & \Tp{MTLHeap}，或 buffer 与 texture 的 storage mode \\
资源绑定 & descriptor set 加 pipeline layout & root signature 加 descriptor heap & argument buffer 或编码器直接设态 \\
管线状态 & \Tp{VkPipeline}（图形与计算两类） & PSO（\Tp{ID3D12}\allowbreak\Tp{PipelineState}） & Render 与 Compute PipelineState \\
命令缓冲池 & \Tp{Vk}\allowbreak\Tp{CommandPool} 分配与回收 command buffer & command allocator 复用 command list & 队列直接建 command buffer \\
内存属性 & \Tp{Vk}\allowbreak\Tp{Memory\allowbreak PropertyFlags}（如 device-local） & heap properties 与 heap flags & storage mode（private 或 shared） \\
\end{longtable}

概念上，\Tp{Vulkan} 把显存、同步、状态都摊成显式对象，交给调用者；\Tp{D3D12}
与它最接近，但用 ``资源状态迁移'' 替代 \Tp{Vulkan} 的 layout；\Tp{Metal} 默认
做更多隐式工作（命令缓冲按编码顺序自然同步），写起来最省样板，但可编排的
精细度低于前两者。

\subsection{同步模型的关键差别}

三者要保证 ``生产与消费不抢跑''，用的原语不同。\Tp{Vulkan} 分两个方向：跨
\DomCPU 与 \DomGPU 用 \Tp{VkFence}，\DomGPU 内部跨队列用 \Tp{VkSemaphore}，
资源可用性靠 pipeline barrier 表达。\Tp{D3D12} 用 \Tp{ID3D12Fence} 做
\DomCPU 侧等待，资源可用性靠显式的 resource barrier 加状态迁移（例如从
render-target 态搬到 shader-read 态）。\Tp{Metal} 大部分依赖命令编码器在同一
command buffer 内的自然顺序，跨 command buffer 用 \Tp{MTLSharedEvent}。共同
点是：忘了插 barrier 或漏了状态迁移，就是竞态与花屏的头号来源。

\subsection{资源绑定：三种心智模型}

\Tp{Vulkan} 把绑定描述成 descriptor set（一组带 layout 的槽位），换 set 换
绑定；\Tp{D3D12} 用 root signature 声明根参数、descriptor heap 承载具体视图；
\Tp{Metal} 既可在编码器上直接 \Tp{setFragmentTexture} 之类逐个绑，也可打包成
argument buffer 一次提交。移植一份渲染代码时，这一层的差异最能决定工作量：
set/heap/argument buffer 的粒度与更新代价各不相同，批处理与逐帧更新要分开设计。

\subsection{管线状态：一次定型还是逐步拼装}

\Tp{Vulkan} 与 \Tp{D3D12} 都要求把渲染状态（混合、深度、拓扑、着色器组合）
提前烘进一个不可变的管线对象（\Tp{VkPipeline} 或 PSO），运行时只切换对象；
变体多时对象数会爆炸，需要变体缓存与预热。换 set/heap/argument buffer 的粒度与更新代价各不相同。
\Tp{Metal} 也提供 \Tp{RenderPipelineState}，但许多状态（视口、裁剪矩形、部分
绑定）可在编码器上逐帧改，不用重建管线，迭代更快、批量切换更慢。移植时的经验：
把 ``每帧都变的小状态'' 留在编码器/root signature 层，把 ``着色器与固定功能
的组合'' 烘进管线，既避免对象爆炸，又保住提交效率。

\subsection{对象寿命与所有权}

显式 API 把 ``谁分配、谁复用、何时能安全销毁'' 变成调用者的责任。
\Tp{Vulkan} 里 command buffer 从 \Tp{VkCommandPool} 分配与回收，绑定到
\Tp{VkDeviceMemory} 的资源要等所有引用它的提交完成才能释放；\Tp{D3D12}
用 command allocator 复用 command list，且同一 allocator 在其记录的命令
未执行完前不能重置；\Tp{Metal} 多用引用计数（ARC）自动管寿命，手动负担最轻，
但也最容易被 ``自然顺序'' 掩盖真实的同步问题。跨端移植时，对象寿命是隐性 bug
的温床：把 \Tp{Metal} 的 ``自动就行'' 直接搬到 \Tp{Vulkan}，往往表现为偶发
崩溃或验证层报 ``资源仍在使用''。

\section{编译与校验命令速查}

工具链按 ``源语言 $\rightarrow$ 中间或目标 $\rightarrow$ 校验与优化'' 排布。
下表给出常用命令与产物；参数只是常见写法，完整选项以工具自带的帮助为准。

\begin{longtable}{>{\raggedright\arraybackslash}p{0.20\textwidth}
                  >{\raggedright\arraybackslash}p{0.16\textwidth}
                  >{\raggedright\arraybackslash}p{0.28\textwidth}
                  >{\raggedright\arraybackslash}p{0.16\textwidth}}
\caption{着色器编译、转换与校验命令速查}\label{tab:apB-tools} \\
\toprule
工具 & 用途 & 常见参数（示例） & 产物 \\
\midrule
\endfirsthead
\caption[]{着色器编译、转换与校验命令速查（续）}\\
\toprule
工具 & 用途 & 常见参数（示例） & 产物 \\
\midrule
\endhead
\bottomrule
\endfoot
\Tp{glslang}\allowbreak\Tp{Validator} & \GSL 到 \SPIV，兼做校验 & \Tp{-V} \Tp{-e} \Tp{main} \Tp{-o} \Tp{out.spv}，可 \Opt{target-env}\allowbreak\Tp{=vulkan1.2} & \SPIV 模块 \\
\Tp{glslc} & \GSL 到 \SPIV（Shaderc） & \Opt{target-env}\allowbreak\Tp{=vulkan1.2} \Tp{-g} \Tp{-o} \Tp{out.spv} & \SPIV 模块 \\
\Tp{dxc} & \HSL 到 DXIL 或 \SPIV & \Tp{-T}\allowbreak\Tp{ps\_6\_0} \Tp{-E} \Tp{main} \Tp{-Fo} \Tp{out.dxil}，或 \Tp{-spirv} & DXIL 容器或 \SPIV \\
\Tp{slangc} & Slang 或 \HSL 到 \SPIV 或 DXIL & \Tp{-target}\allowbreak\Tp{spirv} \Tp{-stage}\allowbreak\Tp{vertex} \Tp{-entry}\allowbreak\Tp{main} \Tp{-o} \Tp{out.spv} & \SPIV 或 DXIL \\
\Tp{spirv-dis} & \SPIV 反读为文本 & \Tp{-o} \Tp{out.txt} 或打到标准输出 & 可读汇编 \\
\Tp{spirv-cross} & \SPIV 再源出 \GSL、\HSL、\MSL & \Tp{--glsl}、\Tp{--hlsl}、\Tp{--msl}，配 \Opt{shader-model} 与 \Opt{output} & 目标语言源码 \\
\Tp{spirv-opt} & \SPIV 优化与规整 & \Tp{-O} \Tp{-o} \Tp{out.spv}，或 \Opt{legalize-types} 加 \Opt{strip-reflect} & 优化后 \SPIV \\
\Tp{spirv-val} & \SPIV 合法性校验 & \Opt{target-env}\allowbreak\Tp{=vulkan1.2} \Tp{in.spv} & 通过或以退出码报错 \\
\Tp{fxc}（旧） & \HSL 到 DXBC（\Tp{D3D11} 时代） & \Tp{/T}\allowbreak\Tp{vs\_5\_0} \Tp{/E} \Tp{main} \Tp{/Fo} \Tp{out.cso} & DXBC（\Tp{.cso}） \\
\end{longtable}

\subsection{一条典型流水线}

把三种工具串起来：\GSL 源码先编成 \SPIV，再校验、优化，最后按目标后端反源。
下面这条命令链就是 ``一份 \GSL、多处落地'' 的最小骨架。

\begin{lstlisting}[style=shellstyle,caption={从 GLSL 到目标语言源码的一条命令链},label={lst:apB-chain}]
# 1) GLSL 源 -> SPIR-V，并固定目标环境
glslangValidator -V --target-env vulkan1.2 -e main -o scene.vert.spv scene.vert
# 2) 校验产物是否合法（不合法即以非零退出码失败）
spirv-val --target-env vulkan1.2 scene.vert.spv
# 3) 需要时做无损优化并去掉反射信息，压缩体积
spirv-opt -O --strip-reflect -o scene.vert.opt.spv scene.vert.spv
# 4) 反源到其它后端：给 Metal 与 HLSL 各出一份
spirv-cross --msl --output scene.vert.metal scene.vert.opt.spv
spirv-cross --hlsl --shader-model 60 --output scene.vert.hlsl scene.vert.opt.spv
\end{lstlisting}

\subsection{产物格式与目标后缀}

不同工具的产物要按后端选对格式与扩展名，否则加载期才暴露问题：
\begin{itemize}
  \item \SPIV 二进制：\Tp{Vulkan} 与跨平台路线的通用产物，常见后缀
        \Tp{.spv}，可直接塞进 \Tp{VkShaderModule}。
  \item DXBC 与 DXIL：\Tp{D3D11} 用 DXBC，\Tp{D3D12} 的着色模型 6.0 起用
        DXIL（\Tp{dxc} 产 \Tp{.dxil} 或封进容器）。
  \item 目标语言源码：\Tp{spirv-cross} 出的 \Tp{.metal}、\Tp{.hlsl}、
        \Tp{.glsl} 需在各自平台用原生编译器再编一次。
\end{itemize}
选路线时的实质区别：走 \SPIV 意味着一次编译、多端加载；走反源意味着为每个
后端各留一条原生编译链。二者都真实可用，取舍在于你更想要 ``编译一次'' 还是
``每端最优''。

\subsection{校验工具与常见报错的对应}

编译期与运行期的报错来源不同，先分清再查：
\begin{itemize}
  \item \Tp{glslangValidator} 报错：多半是 \GSL 语法或 \Tp{layout} 组合非法，
        例如缺默认精度、\Tp{set}/{\Bk}\Tp{binding} 冲突；它也能做 \SPIV 侧的静态检查。
  \item \Tp{spirv-val} 报错：产物不符合目标 \SPIV 规范或 \Tp{target-env}，
        常见于用了目标环境不允许的能力或存储类。
  \item \Tp{spirv-cross} 报错：反源时遇到后端不支持的构造（例如某些
        \SPIV 能力在低版本 \MSL/{\Bk}\HSL 无对应），需要调整目标版本或改着色器。
  \item 运行期验证层（\Tp{Vulkan} validation、\Tp{D3D12} debug layer）报错：
        这是 API 用法违例（资源状态、同步、描述符范围），不是着色器编译问题。
\end{itemize}
把 ``编译能过但验证层报错'' 当成常态来防：编译器管语法与规范，验证层管运行时
约束，两者覆盖不同维度。

\subsection{跨版本核对清单}

\begin{enumerate}
  \item \SPIV 版本与目标环境：确认 \Tp{target-env} 与设备支持的 \SPIV 版本一致。
  \item \HSL 的 shader model：\Tp{ps\_6\_0} 一类目标要匹配 \Tp{dxc} 与运行库版本。
  \item \MSL 的 language version：新语法（如光追、arg tier）依赖较新的 \MSL。
  \item 深度与裁剪约定：换 API 后端时先核对 $y$ 方向与 $z$ 区间，再改投影矩阵。
  \item 绑定编号：\Tp{location}、语义名编号、\Tp{set}/{\Bk}\Tp{binding} 三者要和
        pipeline 的输入输出布局逐一对齐。
\end{enumerate}

\begin{guideline}[title={速查以头文件与官方规范为准}]
本附录只是并排对照的``地图''：语言的关键字集合、语义名、绑定方式与深度约定，
API 的对象名与同步模型，工具的选项与产物，都会随版本变化。凡与本表冲突处，
一律以头文件与官方规范为准，版本差异以对应版本官方文档为准；跨版本移植前
先用 \Tp{spirv-val} 与目标平台的验证层各跑一遍。
\end{guideline}

% ============================================================
% ---- frag_apC.tex ----
\chapter{库、引擎与选型矩阵}
\label{ap:apC-library-choice}

正文第 \ref{ch:ch14-perf-tools} 章给了 ``先测量再选型'' 的方法；本附录把可选
的库、引擎与工具摊开成矩阵，供立项时对照。范围只到 ``接口与许可'' 层面：
抽象层库给的是它暴露的语义，引擎只标到渲染抽象接口，不深入其内部实现。许可
一律以随源码分发的 LICENSE 与官方声明为准。读法很简单：先按 ``平台与许可''
把不合格的行划掉，再在剩下的里面按 ``控制粒度与迁移成本'' 挑，最后用第
\ref{ch:ch14-perf-tools} 章的测量方法验证它真的能跑够帧率。

\section{抽象层与渲染库对比}

下表按 ``控制粒度 / 平台覆盖 / 许可 / 着色器语言路径'' 四个维度对齐常见路线。
控制粒度越高越贴近硬件、越费工程；越低越省心、越受库的天花板约束。

\begin{longtable}{>{\raggedright\arraybackslash}p{0.12\textwidth}
                  >{\raggedright\arraybackslash}p{0.14\textwidth}
                  >{\raggedright\arraybackslash}p{0.16\textwidth}
                  >{\raggedright\arraybackslash}p{0.14\textwidth}
                  >{\raggedright\arraybackslash}p{0.21\textwidth}}
\caption{抽象层与渲染库对比}\label{tab:apC-lib} \\
\toprule
项目 & 控制粒度 & 平台覆盖 & 许可 & 着色器语言路径 \\
\midrule
\endfirsthead
\caption[]{抽象层与渲染库对比（续）}\\
\toprule
项目 & 控制粒度 & 平台覆盖 & 许可 & 着色器语言路径 \\
\midrule
\endhead
\bottomrule
\endfoot
\Lib{bgfx} & 中，库接管提交与内存 & 跨桌面与移动，后端含 Vulkan、D3D、Metal、GL & MIT & 自带 shadinglang，内置编译到各后端 \\
\Lib{wgpu} & 中偏高，走 WebGPU 语义 & 桌面、移动与 Web（wasm），含 \Lib{Dawn} & Apache-2.0 与 MIT 双许可 & WGSL，经 naga 到 \SPIV、MSL、DXIL \\
\Lib{Filament} & 中高，实体渲染器但自动管资源 & 桌面与移动，后端含 Vulkan、Metal、GL & BSD-3-Clause & 自带 \Tp{filam}，编译到 \SPIV 与各后端 \\
原生显式 API & 完全，自管内存、同步、管线 & Vulkan 跨平台、D3D12 仅 Windows 与 Xbox、Metal 仅 Apple & 随 API 与 SDK，无第三方运行时 & GLSL 到 \SPIV、\HSL 到 DXIL、\MSL \\
遗留 \Lib{OpenGL} & 低，状态机与隐式同步 & 覆盖广但 Apple 与移动端逐步弃用 & 随驱动 & \GSL，通常不产出 \SPIV \\
\end{longtable}

\subsection{怎么选控制粒度}

粒度是一条连续谱：左端是原生显式 API（一切自己来），右端是遗留 \Tp{OpenGL}
（一切交给状态机）。中间按 ``是否想自管内存与同步'' 划分——想省心就用
\Tp{wgpu}/{\Bk}\Tp{bgfx} 这类抽象层，想在移动端拿到高质量 PBR 又少造轮子就用
\Tp{Filament}，要极致帧编排与多队列重叠就回到原生 API。抽象层省下的移植
成本，通常以 ``拿不到最新特性'' 和 ``多一层间接开销'' 为代价。

\subsection{许可与分发的注意点}

按许可风险给这几条分档：\Lib{bgfx} 的 MIT、\Lib{Filament} 的 BSD-3-Clause、
\Lib{wgpu} 的 Apache-2.0 与 MIT 双许可都属于宽松许可，商用与闭源分发都友好，
只需保留版权声明；Apache-2.0 额外给了专利授权，对担心专利的团队更稳。真正的
红线是 \Lib{OptiX} 这类专有库：它绑定特定硬件、许可限制再分发，属于 ``厂商
锁定 + 非开源'' 双重约束，用之前必须确认目标端硬件与分发条款都允许。随 API
分发的能力（DXR、Vulkan 光追）不引入第三方运行时许可，但受平台与驱动版本
约束，这也是一种 ``软锁定''。

\subsection{引擎侧只到渲染抽象接口}

若项目已在成熟引擎里，选型往往是 ``用引擎的渲染抽象'' 而非 ``引入新库''。
本附录不展开引擎内部，只标注接口层的共性：现代引擎一般把 ``材质、管线状态、
资源与提交'' 抽象成一套自有的 RHI/渲染接口，把底层后端（Vulkan、D3D12、
Metal）藏在接口之后。在引擎内做优化的正确姿势是沿这套接口改默认开销——
减少多余 pass、别留未压缩纹理、合并全屏读写——而不是绕过接口另起一条后端。
接口越成熟，绕开它的收益通常越小、维护成本越大。

\section{光线追踪与降噪相关库的真实许可面}

光追与降噪的许可差异很大，别把 ``随 API'' 当成 ``开源''。下表区分专有、随 API
与开源三类。

\begin{longtable}{>{\raggedright\arraybackslash}p{0.16\textwidth}
                  >{\raggedright\arraybackslash}p{0.20\textwidth}
                  >{\raggedright\arraybackslash}p{0.20\textwidth}
                  >{\raggedright\arraybackslash}p{0.22\textwidth}}
\caption{光线追踪与降噪相关库的许可与依赖}\label{tab:apC-rt} \\
\toprule
项目 & 许可与开放度 & 依赖与硬件 & 说明 \\
\midrule
\endfirsthead
\caption[]{光线追踪与降噪相关库的许可与依赖（续）}\\
\toprule
项目 & 许可与开放度 & 依赖与硬件 & 说明 \\
\midrule
\endhead
\bottomrule
\endfoot
\Lib{OptiX} & NVIDIA 专有许可，非开源 & 绑定 NVIDIA RTX GPU 与驱动 & 硬件加速光追的高层接口，自带降噪 \\
\Lib{DXR} & 随 D3D12 API 分发 & Windows、Xbox 与支持 DX12 的 GPU & 状态对象加 \HSL 光追阶段，是 API 的一部分 \\
\Lib{Vulkan} 光追 & 随 \Tp{Vulkan} 规范，Khronos 免版税 & 支持 ray-tracing pipeline 的 GPU 与驱动 & 以扩展或核心特性提供 BLAS、TLAS 与光追阶段 \\
\Lib{OIDN}（Open Image Denoise） & Apache-2.0 开源 & 跨平台，CPU 与 GPU 后端 & 实时与离线降噪，常配路径追踪结果 \\
\Lib{RTXDI} & MIT 开源 & 较新 GPU 上的重光照采样 & 用重要性采样降低光线与光源求和成本 \\
\end{longtable}

要点：\Lib{OptiX} 是非开源的 NVIDIA 专有库，用了它就等于绑定单一厂商；而
\Lib{DXR} 与 \Tp{Vulkan} 光追是 API 自带的能力，跟随各自平台的驱动，不额外
引入许可义务；降噪侧 \Lib{OIDN} 与 \Lib{RTXDI} 是真正可自由取用的开源件。

\subsection{把光追与降噪接进现有管线}

工程上光追很少是 ``全场景光追''，而是与光栅混合：主体走光栅管线，反射、阴影
或全局光的一路交给光追，再叠一层降噪。据此选件：跨厂商的实时混合优先 DXR 或
Vulkan 光追（随 API，无额外许可）；离线或允许绑定 NVIDIA 时用 \Lib{OptiX}
拿更省事的高层接口。降噪是独立一环，\Lib{OIDN} 偏通用、\Lib{RTXDI} 偏
``重采样以降光线预算''，两者都可与上面任一条组合。别把降噪当成掩盖噪声的
默认手段——它补的是采样不足的账，根子上还是光线预算与带宽的权衡（见第
\ref{ch:ch14-perf-tools} 章）。

\section{资产与纹理工具}

纹理与压缩是带宽与显存预算（见第 \ref{ch:ch14-perf-tools} 章）的上游，工具链
选对能省掉大量运行期开销。

\begin{longtable}{>{\raggedright\arraybackslash}p{0.28\textwidth}
                  >{\raggedright\arraybackslash}p{0.34\textwidth}
                  >{\raggedright\arraybackslash}p{0.14\textwidth}}
\caption{资产与纹理工具速查}\label{tab:apC-asset} \\
\toprule
工具 & 用途 & 许可 \\
\midrule
\endfirsthead
\caption[]{资产与纹理工具速查（续）}\\
\toprule
工具 & 用途 & 许可 \\
\midrule
\endhead
\bottomrule
\endfoot
\Tp{toktx}、libktx（KTX-Software） & 生成 KTX2 容器，压缩并带 mip & Apache-2.0 \\
Basis Universal（basisu） & 超压缩到 ETC1S、UASTC，跨端转码 & Apache-2.0 \\
\Tp{astcenc} & ASTC 块压缩编码，质量与体积可权衡 & Apache-2.0 \\
DirectXTex & BC 系列编解码、mip、格式转换 & MIT \\
NVIDIA Texture Tools（nvtt） & BC 与 GPU 压缩编码 & MIT \\
\end{longtable}

选压缩工具前先定目标格式：桌面与主机走 BC（用 DirectXTex 或 nvtt），移动走
ASTC（用 \Tp{astcenc}）或超压缩走 Basis 的 ETC1S/UASTC 装进 KTX2 再由运行期
转码。KTX2 的价值在于把 ``压缩格式 + mip + 元数据'' 打包成一种 GPU 可直读
的容器，避免加载期再压一次。

\subsection{离线压缩与运行期转码的取舍}

一条链是 ``离线直接压成目标端原生格式''：桌面压 BC、移动压 ASTC，运行期 GPU
直读，零转码开销，但每种平台要各出一份资产。另一条是 ``压成 Basis（ETC1S 或
UASTC）打进包，运行期再转码到当前 GPU 的原生格式''：一份资产通吃多端、包体
小，代价是转码的一次性 CPU 时间与转码器的集成。资产量大、平台杂时，Basis 加
KTX2 往往更划算；平台单一、追求零运行期开销时，离线原生压缩更稳。

\section{选型矩阵：场景与推荐栈}

把前面几节收敛成一张可执行的对照表。理由列说明为什么这么选，风险列提醒要
提前防的坑。

\begin{longtable}{>{\raggedright\arraybackslash}p{0.20\textwidth}
                  >{\raggedright\arraybackslash}p{0.22\textwidth}
                  >{\raggedright\arraybackslash}p{0.24\textwidth}
                  >{\raggedright\arraybackslash}p{0.17\textwidth}}
\caption{场景 $\times$ 推荐栈 $\times$ 理由 $\times$ 风险}\label{tab:apC-choice} \\
\toprule
场景 & 推荐栈 & 理由 & 风险 \\
\midrule
\endfirsthead
\caption[]{场景 $\times$ 推荐栈 $\times$ 理由 $\times$ 风险（续）}\\
\toprule
场景 & 推荐栈 & 理由 & 风险 \\
\midrule
\endhead
\bottomrule
\endfoot
主机或桌面自研引擎 & 原生 Vulkan 或 D3D12 & 完全掌控内存、同步与提交，便于精细帧编排 & 开发成本高，跨平台要各写一套后端 \\
多平台一次覆盖、团队不大 & \Lib{wgpu} 或 \Lib{bgfx} & 抽象层把移植与同步交给库，接口稳定 & 抽象层有天花板，新特性滞后于原生 \\
移动端要高质量 PBR 又少造轮子 & \Lib{Filament} & 自带材质系统与光照，接口收敛 & 深度定制要改源码，包体偏大 \\
实时或离线光追 & \Lib{OptiX} 或 DXR、\Tp{Vulkan} 光追 & 硬件加速光追接口成熟，降噪有配套 & \Lib{OptiX} 绑定 NVIDIA 且非开源，存在厂商锁定 \\
教学或快速原型 & 遗留 \Lib{OpenGL} & API 简单、资料多、上手快 & 现代平台逐步弃用，控制粒度低 \\
Web 或浏览器 3D & \Lib{wgpu}（含 wasm 与 Dawn） & 与 WebGPU 路线一致，一份代码多处跑 & 浏览器与移动端成熟度参差 \\
\end{longtable}

\subsection{迁移成本与锁定}

选型矩阵里的 ``推荐栈'' 不是终身制。往原生迁，成本在同步、内存与多后端适配；
往抽象层迁，成本在把已经手写的显式资源管理改回库的语义；绑 \Lib{OptiX}，
成本是未来要换厂商时重写光追上层。降低锁定风险的可操作做法：把 ``后端相关
的部分'' 收在一层薄接口后面，让上层只依赖接口；许可敏感的件（专有库）用
可选特性开关隔离，保证拿掉它时其余功能仍可编译运行。

\subsection{什么时候反过来该自研}

三种信号出现时，抽象层的省心会被它的天花板反超：其一，帧预算已经紧到必须
用多队列、异步 compute 与显式资源编排才能塞下；其二，需要的最新 API 特性
迟迟不在库的后端里落地；其三，产品差异化就在渲染层本身，值得为它养一支
图形团队。反之，只要这些信号都不明显，用成熟库把移植、同步与内存交给别人
维护，几乎总是更划算——本附录所有 ``推荐栈'' 都建立在 ``先量出瓶颈'' 的前提上。

\section{排障决策树}

遇到问题先分症状，再按 ``先看什么 $\rightarrow$ 用什么确认'' 走：
\begin{itemize}
  \item 花屏：多半是状态或绑定错。先看抓帧里的 pipeline state（混合、深度
        写入、cull），再用第 \ref{ch:ch14-perf-tools} 章的抓帧流确认实际绑定
        的资源是不是预期的那张。
  \item 黑屏：先看输出通道与矩阵。确认片元确实写了颜色、投影矩阵的 $y$ 方向
        与 $z$ 区间匹配目标 API（见附录 \ref{ap:apB-shading-langs} 的深度约定），
        再查资源 layout 或状态是否可采样。
  \item 性能抖动：先分 \DomCPU 与 \DomGPU 两条时间线。若 GPU 稳而 CPU 抖，
        查提交与同步气泡；若 GPU 抖，查带宽与频率缩放，用分位数别用平均。
  \item 验证层报错：这是用法违例，不是性能问题。逐条读 validation 消息，
        先修资源状态、同步与描述符范围，再谈优化。
\end{itemize}

\begin{lstlisting}[style=shellstyle,caption={排障常用命令：验证层、抓帧与校验},label={lst:apC-debug}]
# 1) 打开 Vulkan 验证层跑一遍，报错先归类到 API 用法
VK_INSTANCE_LAYERS=VK_LAYER_KHRONOS_validation ./bin/demo
vulkaninfo --summary
# 2) 抓一帧，用 RenderDoc GUI 看 pipeline state 与绑定资源
renderdoccmd capture ./bin/demo
# 3) 黑屏或坐标错：校验产物合法并反源核对
spirv-val --target-env vulkan1.2 scene.vert.spv
spirv-cross --msl --output check.metal scene.vert.spv
# 4) 带宽抖动：确认没把未压缩纹理带进包里
toktx --t2 --2d --genmipmap --encode UASTC hero.ktx2 hero.png
\end{lstlisting}

\subsection{抓帧与 profile 的分工}

抓帧工具（如 \Lib{RenderDoc}）回答的是 ``这一帧对不对''：绑定、状态、单个
像素怎么算出来的。平台 profiler 与 GPU 时间戳回答的是 ``这段时间花在哪''：
哪个 pass 长、\DomCPU 与 \DomGPU 谁在等。正确顺序是先抓帧确认没有用法错误，
再用 profile 找时间大头；反过来做，会在一个本来就画错的帧上白调性能。

\subsection{什么时候把问题交给度量}

排障走到 ``状态对、绑定对、验证层干净'' 却仍慢，就已经出了 ``对不对'' 的范畴，
进入第 \ref{ch:ch14-perf-tools} 章的瓶颈判据：按 ALU、带宽、提交、同步四类
用扰动法定位，再回到本附录的选型矩阵看是不是当前栈的天花板。

决策树的使用纪律：一次只动一处、动完立刻复测。同时改绑定又改矩阵，就无法
把 ``修好了'' 归因给哪一步。先让验证层干净，再让抓帧干净，最后才是让时间轴
干净。

\subsection{四种症状的最小确认动作}

\begin{itemize}
  \item 花屏的最小确认：抓帧后关掉混合与深度写入各跑一次，看异常是否消失，
        以定位是状态问题还是资源内容问题。
  \item 黑屏的最小确认：把片元输出改成常量色再跑，画面出现说明是着色计算或
        绑定问题，仍全黑说明是矩阵、裁剪或 layout 让几何根本没进屏幕。
  \item 抖动的最小确认：只测 GPU 时间戳的 p95，不看平均值；把分辨率减半对比，
        先分清是带宽还是 ALU。
  \item 验证层报错的最小确认：读第一条消息的对象与阶段，改一处再重跑，别一次
        修多条导致无法归因。
\end{itemize}

\begin{bestpractice}[title={选型先看许可与平台，排障先让验证层归零}]
选库先确认两件事：许可是否匹配分发策略（尤其 \Lib{OptiX} 这类非开源件），
平台覆盖是否含全部目标端；两者都过关，再谈控制粒度与性能。排障的顺序固定：
验证层 $\rightarrow$ 抓帧状态与绑定 $\rightarrow$ 时间轴与带宽。跳步会让人
把 ``用法违例'' 当成 ``性能瓶颈'' 去优化，方向一错就全盘皆输。
\end{bestpractice}

\end{document}
